% Generated mechanically from frozen input canon/ko-kr/integration-20260908-r10-p03-register/inputs/p03/ko/smoothing.tex.
% Do not edit this generated file; edit the bound locale source instead.

% OK, start here.
%

\setcounter{chapter}{15}
\chapter{환 준동형의 평활화}



\phantomsection
\label{smoothing-section-phantom}


\section{서론}
\label{smoothing-section-introduction}

\noindent
이 장의 주된 결과는 다음과 같다.
$$
\fbox{뇌터 환 사이의 정칙 준동형은 매끄러운 준동형들의 여과 여극한이다.}
$$
이 정리는 Popescu의 결과이다. \cite{popescu-letter}를 보라.
Richard Swan은 Popescu의 증명을 읽기 쉽게 설명하였다.
\cite{swan}을 보라. Swan은 Andr\'e의 노트와 Ogoma의 논문
\cite{Ogoma}를 이용하였다.

\medskip\noindent
여기서는 Swan의 설명을 따르되, 조금 더 단순한 상황에서 논의할 수 있게 해 주는
중간 결과를 먼저 증명한다. 전체 흐름은 다음과 같다.
먼저 다음 “올림 문제”를 해결한다. 매끄러운 대수들의 여과 여극한의 평탄한 무한소
변형은 매끄러운 대수들의 여과 여극한이다. 이 결과는 본질적으로 주정리를
축약된 뇌터 환 사이의 준동형에 대해서만 증명하면 충분하다는 뜻이다.
다음으로 “올림 보조정리”와 “특이점 해소 보조정리”라고 부르는
두 가지 매우 기교적인 보조정리를 증명한다. 이들을 결합하면 주정리는
체 $k$ 위의 기하적으로 정칙인 뇌터 대수 $\Lambda$가 매끄러운 $k$-대수들의
여과 여극한임을 증명하는 문제로 귀착됨을 보인다.
그다음 Popescu-Ogoma-Swan-Andr\'e 증명에 필요한 국소적 기법들을 설명한다.
마지막 세 절에서는 증명을 마무리한다.

\medskip\noindent
서론을 마치기 전에 관련 문헌을 안내한다.
$A$를 헨젤 뇌터 국소환이라 하자.
임의의 $f_1, \ldots, f_m \in A[x_1, \ldots, x_n]$에 대하여,
연립방정식 $f_1 = 0, \ldots, f_m = 0$이 $A$의 완비화에서 해를 가질 필요충분조건이
$A$에서 해를 갖는 것일 때, $A$가 {\it 근사 성질}을 갖는다고 한다.
이 정의는 Artin이 도입하였다.
Artin은 해석적 연립방정식에 대해 근사 성질을 처음 증명하였다.
\cite{Artin-Analytic-Approximation}을 보라.
\cite{Artin-Algebraic-Approximation}에서 Artin은 우수 이산 값매김환 위에서
본질적으로 유한형인 국소환에 대해 근사 성질을 증명하였다.
Artin은 모든 우수 헨젤 국소환이 근사 성질을 가질 것이라고 추측하였다
(\cite{Artin-Algebraic-Approximation}의 26쪽).

\medskip\noindent
어느 시점부터 뇌터 환 사이의 모든 정칙 준동형이 매끄러운 대수들의
여과 여극한이라는 추측이 제기되었다. 예를 들어
\cite{Raynaud-Rennes}, \cite{popescu-global}, \cite{Artin-power-series},
\cite{Artin-Denef}를 보라.
이 추측\footnote{\cite{Artin-power-series}, \cite{Artin-Denef},
\cite{popescu-global}에서 정식화된 질문 또는 추측은 더 강한 명제이며,
\cite{Cipu}에서 원래 형태와 동치임이 증명되었다.}을 누가 처음 제기했는지는
확실하지 않다. 근사 성질과의 관계는 다음과 같다.
$A$가 위와 같을 때 $A \to A^\wedge$가 매끄러운 대수들의 여극한이면
근사 성질이 성립한다(추후 참조를 여기에 삽입).
또한 $A$가 우수 국소환이면 주정리를 사상 $A \to A^\wedge$에 적용할 수 있다.
우수 국소환의 조건 중 하나가 형식적 섬유의 기하적 정칙성이기 때문이다.
우수 국소환은 Grothendieck이 정의하였으며, 그 정의는 1965년에 출판되었음을
덧붙인다.

\medskip\noindent
\cite{Artin-power-series}에서는 체 위에서 본질적으로 유한형인 임의의 국소환
$R$에 대해 $R \to R^\wedge$가 매끄러운 대수들의 여과 여극한임을 증명하였다.
\cite{Rotthaus-Artin}에서는 우수 이산 값매김환 위에서 본질적으로 유한형인
임의의 국소환 $R$에 대해 $R \to R^\wedge$가 매끄러운 대수들의 여과 여극한임을
증명하였다. 끝으로, 앞서 설명한 것처럼
\cite{popescu-GND}, \cite{popescu-GNDA}, \cite{popescu-letter},
\cite{Ogoma}, \cite{swan}에서 주정리가 증명되었다.

\medskip\noindent
반대로, \cite{Rotthaus-excellent}에서는 위 결과 중 일부를 이용하여,
근사 성질을 갖는 모든 뇌터 국소환이 우수환임을 증명하였다.

\medskip\noindent
논문 \cite{Spivakovsky}는 주정리에 대한 다른 접근법을 제시하지만, 읽기가
어려운 것으로 보인다. 예를 들어 \cite[Lemma 5.2]{Spivakovsky}는
\cite[Lemma 3]{Elkik}를 잘못 다시 정식화한 것으로 보인다.
이 내용에 대한 Bourbaki 강연도 있다. \cite{Teissier}를 보라.












\section{특이 아이디얼}
\label{smoothing-section-singular-ideal}

\noindent
$R \to A$를 환 준동형이라 하자. $R$ 위의 $A$의 특이 아이디얼은
사상 $\Spec(A) \to \Spec(R)$의 특이점 집합을 정의하는 $A$의 근기 아이디얼이다.
정확한 정의는 다음과 같다.

\begin{definition}
\label{smoothing-definition-singular-ideal}
$R \to A$를 환 준동형이라 하자.
{\it $R$ 위의 $A$의 특이 아이디얼}은 $H_{A/R}$로 표기하며,
다음을 만족하는 유일한 근기 아이디얼 $H_{A/R} \subset A$이다.
$$
V(H_{A/R}) = \{\mathfrak q \in \Spec(A) \mid R \to A
\text{가 }\mathfrak q\text{에서 매끄럽지 않다}\}
$$
\end{definition}

\noindent
$R \to A$가 매끄러운 소 아이디얼들의 집합은 열려 있으므로 이 정의는 의미가 있다.
Algebra, Definition \ref{algebra-definition-smooth-at-prime}을 보라.
특이 아이디얼의 생성원들을 명시적으로 구하기 위해 먼저 다음 보조정리를 증명한다.

\begin{lemma}
\label{smoothing-lemma-find-strictly-standard}
$R$을 환이라 하고, $A = R[x_1, \ldots, x_n]/(f_1, \ldots, f_m)$이라 하자.
$\mathfrak q \subset A$를 소 아이디얼이라 하자.
$R \to A$가 $\mathfrak q$에서 매끄럽다고 가정하자.
그러면 $a \in A$, $a \not \in \mathfrak q$인 원소와
$0 \leq c \leq \min(n, m)$인 정수 $c$, 그리고 크기가 $c$인 부분집합
$U \subset \{1, \ldots, n\}$, $V \subset \{1, \ldots, m\}$이 존재하여,
어떤 $a' \in A$에 대해
$$
a = a' \det(\partial f_j/\partial x_i)_{j \in V, i \in U}
$$
이고, 모든 $\ell \in \{1, \ldots, m\}$에 대해
$$
a f_\ell \in (f_j, j \in V) + (f_1, \ldots, f_m)^2
$$
이다.
\end{lemma}

\begin{proof}
$I = (f_1, \ldots, f_m)$으로 놓으면, $R$ 위의 $A$의 소박한 여접복합체는
$I/I^2 \to \bigoplus A\text{d}x_i$와 호모토피 동치이다.
Algebra, Lemma \ref{algebra-lemma-NL-homotopy}를 보라.
소박한 여접복합체의 구성이 국소화와 가환한다는 사실을 사용한다.
Algebra, Section \ref{algebra-section-netherlander},
특히 Algebra, Lemma \ref{algebra-lemma-localize-NL}을 보라.
Algebra, Definitions \ref{algebra-definition-smooth}와
\ref{algebra-definition-smooth-at-prime}에 의해, 어떤 $a \in A$,
$a \not \in \mathfrak q$에 대해
$(I/I^2)_a \to \bigoplus A_a\text{d}x_i$는 분할 단사이다.
$x_1, \ldots, x_n$과 $f_1, \ldots, f_m$의 번호를 다시 매기면,
$f_1, \ldots, f_c$가 벡터공간 $I/I^2 \otimes_A \kappa(\mathfrak q)$의 기저를
이루고, $\text{d}x_{c + 1}, \ldots, \text{d}x_n$의 상이
$\Omega_{A/R} \otimes_A \kappa(\mathfrak q)$의 기저를 이룬다고 가정할 수 있다.
따라서 어떤 $a' \in A$, $a' \not \in \mathfrak q$에 대해 $a$를 $aa'$로
바꾸면, $f_1, \ldots, f_c$가 $(I/I^2)_a$의 기저를 이루고
$\text{d}x_{c + 1}, \ldots, \text{d}x_n$의 상이 $(\Omega_{A/R})_a$의
기저를 이룬다고 가정할 수 있다. 이 상황에서는 충분히 큰 정수 $N$에 대한
$a^N$이 보조정리의 조건을 만족한다($U = V = \{1, \ldots, c\}$로 둔다).
\end{proof}

\noindent
$R$ 위의 $A$에서 {\it 엄밀 표준} 원소라는 개념을 사용한다.
여기서의 개념은 Swan의 논문 \cite{swan}에 나오는 개념보다 조금 약하다.
또한 위 보조정리에서 구한 유형의 원소를 {\it 초등 표준} 원소로 정의한다.
이 서로 다른 유형의 원소들을 보조정리 \ref{smoothing-lemma-compare-standard}에서 비교한다.

\begin{definition}
\label{smoothing-definition-strictly-standard}
$R \to A$를 유한 표시 환 준동형이라 하자.
원소 $a \in A$가 {\it $R$ 위의 $A$에서 초등 표준}이라는 것은,
표시 $A = R[x_1, \ldots, x_n]/(f_1, \ldots, f_m)$과
$0 \leq c \leq \min(n, m)$이 존재하여, 어떤 $a' \in A$에 대해
\begin{equation}
\label{smoothing-equation-elementary-standard-one}
a = a' \det(\partial f_j/\partial x_i)_{i, j = 1, \ldots, c}
\end{equation}
이고, $j = 1, \ldots, m - c$에 대해
\begin{equation}
\label{smoothing-equation-elementary-standard-two}
a f_{c + j} \in (f_1, \ldots, f_c) + (f_1, \ldots, f_m)^2
\end{equation}
인 것을 말한다. 원소 $a \in A$가 {\it $R$ 위의 $A$에서 엄밀 표준}이라는 것은,
표시 $A = R[x_1, \ldots, x_n]/(f_1, \ldots, f_m)$과
$0 \leq c \leq \min(n, m)$이 존재하여, 어떤 $a_I \in A$에 대해
\begin{equation}
\label{smoothing-equation-strictly-standard-one}
a = \sum\nolimits_{I \subset \{1, \ldots, n\},\ |I| = c}
a_I \det(\partial f_j/\partial x_i)_{j = 1, \ldots, c,\ i \in I}
\end{equation}
이고, $j = 1, \ldots, m - c$에 대해
\begin{equation}
\label{smoothing-equation-strictly-standard-two}
a f_{c + j} \in (f_1, \ldots, f_c) + (f_1, \ldots, f_m)^2
\end{equation}
인 것을 말한다.
\end{definition}

\noindent
다음 보조정리는 (\ref{smoothing-equation-strictly-standard-one})의 귀결을 알아내는 데 유용하다.

\begin{lemma}
\label{smoothing-lemma-parse-equation-strictly-standard-one}
$R$을 환이라 하고, $A = R[x_1, \ldots, x_n]/(f_1, \ldots, f_m)$이라 하자.
$I = (f_1, \ldots, f_m)$으로 쓰고, $a \in A$라 하자.
그러면 (\ref{smoothing-equation-strictly-standard-one})로부터 다음을 만족하는
$A$-선형사상
$\psi : \bigoplus\nolimits_{i = 1, \ldots, n} A \text{d}x_i \to A^{\oplus c}$가
존재한다. 합성
$$
A^{\oplus c} \xrightarrow{(f_1, \ldots, f_c)}
I/I^2 \xrightarrow{f \mapsto \text{d}f}
\bigoplus\nolimits_{i = 1, \ldots, n} A \text{d}x_i
\xrightarrow{\psi}
A^{\oplus c}
$$
은 $a$를 곱하는 사상이다. 반대로, 이러한 $\psi$가 존재하면
$a^c$는 (\ref{smoothing-equation-strictly-standard-one})을 만족한다.
\end{lemma}

\begin{proof}
이는 Algebra, Lemma \ref{algebra-lemma-matrix-left-inverse}의 특수한 경우이다.
\end{proof}

\begin{lemma}[Elkik]
\label{smoothing-lemma-elkik}
$R \to A$를 유한 표시 환 준동형이라 하자.
특이 아이디얼 $H_{A/R}$는 $R$ 위의 $A$에서 엄밀 표준인 원소들로 생성되는
아이디얼의 근기이며, 또한 $R$ 위의 $A$에서 초등 표준인 원소들로 생성되는
아이디얼의 근기이다.
\end{lemma}

\begin{proof}
$a$가 $R$ 위의 $A$에서 엄밀 표준이라고 가정하자.
$A_a$가 $R$ 위에서 매끄럽다고 주장한다. 이를 보이면 $a \in H_{A/R}$가 따른다.
구체적으로, $A = R[x_1, \ldots, x_n]/(f_1, \ldots, f_m)$, $c$, $a' \in A$를
정의 \ref{smoothing-definition-strictly-standard}에서와 같이 잡는다.
$I = (f_1, \ldots, f_m)$으로 놓으면 $R$ 위의 $A$의 소박한 여접복합체는
$I/I^2 \to \bigoplus A\text{d}x_i$로 주어진다.
가정 (\ref{smoothing-equation-strictly-standard-two})로부터
$(I/I^2)_a$는 $f_1, \ldots, f_c$의 잉여류들로 생성됨을 알 수 있다.
가정 (\ref{smoothing-equation-strictly-standard-one})로부터 미분
$(I/I^2)_a \to \bigoplus A_a\text{d}x_i$가 왼쪽 역을 갖는다는 것이 따른다.
보조정리 \ref{smoothing-lemma-parse-equation-strictly-standard-one}을 보라.
따라서 정의와 Algebra, Lemma \ref{algebra-lemma-localize-NL}에 의해
$R \to A_a$는 매끄럽다.

\medskip\noindent
$H_e, H_s \subset A$를 각각 $R$ 위의 $A$에서 초등 표준인 원소들,
엄밀 표준인 원소들로 생성되는 아이디얼의 근기라 하자.
정의와 방금 증명한 사실에 의해 $H_e \subset H_s \subset H_{A/R}$이다.
포함관계 $H_{A/R} \subset H_e$는 보조정리
\ref{smoothing-lemma-find-strictly-standard}에서 따른다.
\end{proof}

\begin{example}
\label{smoothing-example-not-quasi-compact}
유한 표시 환 준동형이 매끄러운 점들의 집합은 준콤팩트 열린집합일 필요가 없다.
예를 들어 $R = k[x, y_1, y_2, y_3, \ldots]/(xy_i)$,
$A = R/(x)$라 하자. 그러면 $R \to A$의 매끄러운 점들의 집합은
$\bigcup D(y_i)$이며, 이는 준콤팩트가 아니다.
\end{example}

\begin{lemma}
\label{smoothing-lemma-strictly-standard-base-change}
$R \to A$를 유한 표시 환 준동형이라 하고, $R \to R'$를 환 준동형이라 하자.
$a \in A$가 $R$ 위의 $A$에서 초등 표준(각각 엄밀 표준)이면,
$a \otimes 1$은 $R'$ 위의 $A \otimes_R R'$에서 초등 표준(각각 엄밀 표준)이다.
\end{lemma}

\begin{proof}
$A = R[x_1, \ldots, x_n]/(f_1, \ldots, f_m)$이 $R$ 위의 $A$의 표시이면,
$A \otimes_R R' = R'[x_1, \ldots, x_n]/(f'_1, \ldots, f'_m)$은
$R'$ 위의 $A \otimes_R R'$의 표시이다.
여기서 $f'_j$는 $R'[x_1, \ldots, x_n]$에서의 $f_j$의 상이다.
따라서 결론은 정의에서 따른다.
\end{proof}

\begin{lemma}
\label{smoothing-lemma-final-solve}
$R \to A \to \Lambda$를 환 준동형들이라 하고,
$A$는 $R$ 위에서 유한 표시라고 하자.
$H_{A/R} \Lambda = \Lambda$라고 가정하자.
그러면 $R$ 위에서 매끄러운 $B$를 통한 분해 $A \to B \to \Lambda$가 존재한다.
\end{lemma}

\begin{proof}
$\Lambda$에서 $\sum f_i\lambda_i = 1$이 되도록
$f_1, \ldots, f_r \in H_{A/R}$와
$\lambda_1, \ldots, \lambda_r \in \Lambda$를 고른다.
$B = A[x_1, \ldots, x_r]/(f_1x_1 + \ldots + f_rx_r - 1)$로 놓고,
$x_i$를 $\lambda_i$로 보내어 $B \to \Lambda$를 정의한다.
$B$가 $R$ 위에서 매끄러움을 확인하려면, $H_{A/R}$의 정의에 의해
$A_{f_i}$가 $R$ 위에서 매끄럽고 $B_{f_i}$가 $A_{f_i}$ 위에서 매끄럽다는 것을
사용하면 된다. 자세한 내용은 생략한다.
\end{proof}





\section{대수의 표시}
\label{smoothing-section-presentations}

\noindent
이 절의 결과 중 일부는 Elkik의 결과이다. 다음 보조정리에 나오는 대수
$C$는 $A$ 위의 대칭대수임에 유의하자. 또한 $R$가 뇌터 환이면
$C$는 $R$ 위에서 유한 표시를 갖는다.

\begin{lemma}
\label{smoothing-lemma-improve-presentation}
$R$를 환, $A$를 유한 표시 $R$-대수라 하자.
다음 두 성질을 가지며 수축 사상을 갖는 유한형 $R$-대수 사상
$A \to C$가 존재한다.
\begin{enumerate}
\item $R \to A_a$가 국소 완전교차인
(More on Algebra, Definition
\ref{more-algebra-definition-local-complete-intersection})
각 $a \in A$에 대하여, 환 $C_a$는 $A_a$ 위에서 매끄럽고
$J/J^2$가 $C_a$ 위의 자유 가군이 되는 표시
$C_a = R[y_1, \ldots, y_m]/J$를 갖는다.
\item $A_a$가 $R$ 위에서 매끄러운 각 $a \in A$에 대하여,
가군 $\Omega_{C_a/R}$는 $C_a$ 위의 자유 가군이다.
\end{enumerate}
\end{lemma}

\begin{proof}
표시 $A = R[x_1, \ldots, x_n]/I$를 택하고
$I = (f_1, \ldots, f_m)$으로 쓰자. 짧은 완전열
$$
0 \to K \to A^{\oplus m} \to I/I^2 \to 0
$$
로 $A$-가군 $K$를 정의한다. 여기서 가운데 항의 $j$번째 기저 벡터
$e_j$는 오른쪽의 $f_j$의 동치류로 보내진다.
$$
C = \text{Sym}^*_A(I/I^2).
$$
로 놓자. 수축 사상은 단순히 $C$의 차수 $0$ 부분으로의 사영이다.
$y_j$를 $I/I^2$에서의 $f_j$의 동치류로 보내는 전사 사상
$R[x_1, \ldots, x_n, y_1, \ldots, y_m] \to C$가 있다.
이 사상의 핵 $J$는 원소 $f_1, \ldots, f_m$과,
$\sum h_j e_j$가 $K$의 원소를 정의하는
$h_j \in R[x_1, \ldots, x_n]$에 대한 원소 $\sum h_j y_j$로 생성된다.
$R \to A \to C$와 위의 표시들에 적용한
Algebra, Lemma \ref{algebra-lemma-exact-sequence-NL} 및
More on Algebra, Lemma
\ref{more-algebra-lemma-cotangent-complex-symmetric-algebra}에 의해,
$C$-가군의 짧은 완전열
\begin{equation}
\label{smoothing-equation-sequence}
I/I^2 \otimes_A C \to J/J^2 \to K \otimes_A C \to 0
\end{equation}
이 존재한다. $h \in R[x_1, \ldots, x_n]$를 상이 $a \in A$인
원소라 하자. 국소화한 환들의 표시로
$$
A_a = R[x_0, x_1, \ldots, x_n]/I'
\quad\text{및}\quad
C_a = R[x_0, x_1, \ldots, x_n, y_1, \ldots, y_m]/J'
$$
를 사용한다. 여기서 $I' = (hx_0 - 1, I)$이고 $J' = (hx_0 - 1, J)$이다.
따라서 $A_a$-가군으로서 $I'/(I')^2 = A_a \oplus (I/I^2)_a$이고,
$C_a$-가군으로서 $J'/(J')^2 = C_a \oplus (J/J^2)_a$이다.
그러므로 $R \to A_a \to C_a$와 위의 표시들에 대응하는
Algebra, Lemma \ref{algebra-lemma-exact-sequence-NL}의 완전열로
\begin{equation}
\label{smoothing-equation-sequence-localized}
C_a \oplus I/I^2 \otimes_A C_a \to
C_a \oplus (J/J^2)_a \to
K \otimes_A C_a \to 0
\end{equation}
을 얻는다.

\medskip\noindent
이제 $a \in A$가 $A_a$를 $R$ 위의 국소 완전교차로 만드는 원소라고
가정하자. 그러면 $(I/I^2)_a$는 $A_a$ 위의 유한 사영 가군이다.
More on Algebra, Lemma
\ref{more-algebra-lemma-quasi-regular-ideal-finite-projective}를 보라.
따라서 $K_a \oplus (I/I^2)_a \cong A_a^{\oplus m}$은 자유 가군이다.
특히 $K_a$도 유한 사영 가군이다.
More on Algebra, Lemma \ref{more-algebra-lemma-transitive-lci-at-end}에 의해
완전열 (\ref{smoothing-equation-sequence-localized})은 왼쪽에서도 완전하다.
따라서
$$
J'/(J')^2 \cong
C_a \oplus I/I^2 \otimes_A C_a \oplus K \otimes_A C_a \cong
C_a^{\oplus m + 1}
$$
이다. 이로써 (1)이 증명된다. 끝으로 $A_a$가 추가로 $R$ 위에서
매끄럽다고 가정하자. 그러면 같은 표시에서 $\Omega_{C_a/R}$가 사상
$$
J'/(J')^2 \longrightarrow
\bigoplus\nolimits_i C_a\text{d}x_i \oplus \bigoplus\nolimits_j C_a\text{d}y_j
$$
의 여핵임을 알 수 있다. 위의 분해에서 $J'/(J')^2$의 직합인자 $C_a$는
$hx_0 - 1$에 대응하므로 직합인자 $C_a\text{d}x_0$로 동형으로 보내진다.
$J'/(J')^2$의 직합인자 $I/I^2 \otimes_A C_a$는
$\bigoplus_{i = 1, \ldots, n} C_a\text{d}x_i$로 단사적으로 보내지며,
그 몫은 $\Omega_{A_a/R} \otimes_{A_a} C_a$이다.
직합인자 $K \otimes_A C_a$는
$\bigoplus_{j \geq 1} C_a\text{d}y_j$로 단사적으로 보내지며,
그 몫은 $I/I^2 \otimes_A C_a$와 동형이다.
따라서 마지막에 표시한 사상의 여핵은 가군
$I/I^2 \otimes_A C_a \oplus \Omega_{A_a/R} \otimes_{A_a} C_a$이다.
$(I/I^2)_a \oplus \Omega_{A_a/R}$가 자유 가군이므로
(매끄러운 환 준동형의 정의에 의해), (2)가 성립한다.
\end{proof}

\noindent
다음 명제는 헨젤 쌍 위의 매끄러운 환 준동형에 대해 Elkik이
\cite{Elkik}에서 증명하였다. 매끄러운 환 준동형에 대한 결과는
\cite{Arabia}에서도 찾을 수 있으며, 그곳에서는 매끄러운 대수 사이의
환 준동형도 올릴 수 있음을 증명한다.

\begin{proposition}
\label{smoothing-proposition-lift-smooth}
\begin{slogan}
매끄러운 대수와 신토믹 대수는 전사 사상을 따라 올릴 수 있다
\end{slogan}
$R \to R_0$를 핵이 $I$인 전사 환 준동형이라 하자.
\begin{enumerate}
\item $R_0 \to A_0$가 신토믹(syntomic) 환 준동형이면,
$A/IA \cong A_0$인 신토믹 환 준동형 $R \to A$가 존재한다.
\item $R_0 \to A_0$가 매끄러운 환 준동형이면,
$A/IA \cong A_0$인 매끄러운 환 준동형 $R \to A$가 존재한다.
\end{enumerate}
\end{proposition}

\begin{proof}
$R_0 \to A_0$가 신토믹이라고 가정하자. 특히 이는 국소 완전교차이다
(More on Algebra, Lemma \ref{more-algebra-lemma-syntomic-lci}).
표시 $A_0 = R_0[x_1, \ldots, x_n]/J_0$를 택하고
$C_0 = \text{Sym}^*_{A_0}(J_0/J_0^2)$으로 놓자.
$J_0/J_0^2$는 유한 사영 $A_0$-가군임에 유의하자
(Algebra, Lemma
\ref{algebra-lemma-syntomic-presentation-ideal-mod-squares}).
보조정리 \ref{smoothing-lemma-improve-presentation}에 의해 환 준동형
$A_0 \to C_0$는 매끄러우며, $K_0/K_0^2$가 $C_0$ 위의 자유 가군인
표시 $C_0 = R_0[y_1, \ldots, y_m]/K_0$를 찾을 수 있다.
Algebra, Lemma \ref{algebra-lemma-huber}에 의해,
$\overline{f}_1, \ldots, \overline{f}_c$가
$C_0$ 위의 $K_0/K_0^2$의 기저로 보내지는
$C_0 = R_0[y_1, \ldots, y_m]/(\overline{f}_1, \ldots, \overline{f}_c)$의
형태라고 가정할 수 있다.
$\overline{f}_1, \ldots, \overline{f}_c$를 올리는
$f_1, \ldots, f_c \in R[y_1, \ldots, y_c]$를 택하고
$$
C = R[y_1, \ldots, y_m]/(f_1, \ldots, f_c)
$$
로 놓자. 구성에 의해 $C_0 = C/IC$이다.
Algebra, Lemma
\ref{algebra-lemma-localize-relative-complete-intersection}에 의해,
$C$를 $C_g$로 바꾼 뒤에는 $C$가 $R$ 위의 상대적 전역 완전교차라고
가정할 수 있다.
결국 어떤 신토믹 $R$-대수 $C$에 대해
$C_0 = \text{Sym}^*_{A_0}(P_0)$가 $C/IC$와 동형이 되는
유한 사영 $A_0$-가군 $P_0$가 존재한다.

\medskip\noindent
정수 $n$과 직합 분해 $A_0^{\oplus n} = P_0 \oplus Q_0$를 택하자.
More on Algebra, Lemma \ref{more-algebra-lemma-lift-projective-module}에 의해,
동형 $C/IC \to C'/IC'$를 유도하는 에탈 환 준동형 $C \to C'$와,
$Q/IQ$가 $Q_0 \otimes_{A_0} C/IC$와 동형인 유한 사영
$C'$-가군 $Q$를 찾을 수 있다.
그러면 $D = \text{Sym}_{C'}^*(Q)$는 매끄러운 $C'$-대수이다
(More on Algebra, Lemma
\ref{more-algebra-lemma-symmetric-algebra-smooth} 참조).
그림으로 나타내면 다음과 같다.
$$
\xymatrix{
R \ar[d] \ar[rr] & &
C \ar[r] \ar[d] &
C' \ar[r] \ar[d] &
D \ar[d] \\
R/I \ar[r] &
A_0 \ar[r] &
C/IC \ar[r]^{\cong} &
C'/IC' \ar[r] &
D/ID
}
$$
$Q$를 위와 같이 택했으므로
\begin{align*}
D/ID & =
\text{Sym}_{C/IC}^*(Q_0 \otimes_{A_0} C/IC) \\
& =
\text{Sym}_{A_0}^*(Q_0) \otimes_{A_0} C/IC \\
& =
\text{Sym}_{A_0}^*(Q_0) \otimes_{A_0}
\text{Sym}_{A_0}^*(P_0) \\
& =
\text{Sym}_{A_0}^*(Q_0 \oplus P_0) \\
& =
\text{Sym}_{A_0}^*(A_0^{\oplus n}) \\
& =
A_0[x_1, \ldots, x_n]
\end{align*}
이다. $D/ID = A_0[x_1, \ldots, x_n]$에서
$x_1, \ldots, x_n$으로 보내지는 $f_1, \ldots, f_n \in D$를 택하자.
$A = D/(f_1, \ldots, f_n)$으로 놓으면 $A_0 = A/IA$이다.
$R \to A$가 $V(IA)$의 한 근방에서 신토믹임을 주장한다.
이 주장이 참이면, $1 \in A_0$로 보내지며 $A_f$가 $R$ 위에서
신토믹이 되는 $f \in A$를 찾을 수 있으므로 (1)의 증명이 끝난다.

\medskip\noindent
주장의 증명. $R \to D$는 신토믹 환 준동형 $R \to C$,
에탈 환 준동형 $C \to C'$, 매끄러운 환 준동형 $C' \to D$의
합성이므로 신토믹이다(Algebra, Lemmas
\ref{algebra-lemma-composition-syntomic} 및
\ref{algebra-lemma-smooth-syntomic}).
이 문제는 $\Spec(D)$ 위에서 국소적이므로, $D$가 상대적 전역
완전교차라고 가정해도 된다
(Algebra, Lemma \ref{algebra-lemma-syntomic}).
$D = R[y_1, \ldots, y_m]/(g_1, \ldots, g_s)$라고 하자.
$f_1, \ldots, f_n$의 올림
$f'_1, \ldots, f'_n \in R[y_1, \ldots, y_m]$를 택하자.
그러면 Algebra, Lemma
\ref{algebra-lemma-localize-relative-complete-intersection}을
적용하여 주장을 얻는다.

\medskip\noindent
(2)의 증명. 매끄러운 환 준동형은 신토믹이므로,
$A_0 = A/IA$인 신토믹 환 준동형 $R \to A$를 찾을 수 있다.
가정에 의해 $V(I)$의 소아이디얼들 위에서 $R \to A$의 섬유들은
매끄러우므로, $R \to A$는 $V(IA)$의 한 열린 근방에서 매끄럽다
(Algebra, Lemma \ref{algebra-lemma-flat-fibre-smooth}).
따라서 $A$를 국소화로 바꾸면 원하는 결과를 얻는다.
\end{proof}

\noindent
임의의 신토믹 환 준동형 $R \to A$가 국소적으로 상대적 전역
완전교차임을 알고 있다.
Algebra, Lemma \ref{algebra-lemma-syntomic}를 보라.
다음 보조정리는 $\Spec(A)$ 위의 한 벡터 다발이 상대적 전역
완전교차임을 말한다.

\begin{lemma}
\label{smoothing-lemma-syntomic-complete-intersection}
$R \to A$를 신토믹 환 준동형이라 하자.
$R$ 위의 전역 상대적 완전교차 $C$에 대해, 수축 사상을 갖는 매끄러운
$R$-대수 사상 $A \to C$가 존재한다. 즉,
$$
C \cong R[x_1, \ldots, x_n]/(f_1, \ldots, f_c)
$$
는 $R$ 위에서 평탄하고 모든 섬유의 차원이 $n - c$이다.
\end{lemma}

\begin{proof}
보조정리 \ref{smoothing-lemma-improve-presentation}을 적용하여 $A \to C$를 얻는다.
Algebra, Lemma \ref{algebra-lemma-huber}에 의해,
$f_i$가 $J/J^2$의 기저로 보내지도록
$C = R[x_1, \ldots, x_n]/(f_1, \ldots, f_c)$로 쓸 수 있다.
환 준동형 $R \to C$는 신토믹 환 준동형과 매끄러운 환 준동형의
합성이므로 신토믹이다(따라서 평탄하다).
Algebra, Lemma \ref{algebra-lemma-lci}에 의해 섬유의 차원은 $n - c$이다
(섬유들이 국소 완전교차이므로 이 보조정리를 적용할 수 있다).
\end{proof}

\begin{lemma}
\label{smoothing-lemma-smooth-standard-smooth}
$R \to A$를 매끄러운 환 준동형이라 하자.
$B$가 $R$ 위에서 표준 매끄러운 대수가 되도록, 수축 사상을 갖는
매끄러운 $R$-대수 사상 $A \to B$가 존재한다. 즉,
$$
B \cong R[x_1, \ldots, x_n]/(f_1, \ldots, f_c)
$$
이며 $\det(\partial f_j/\partial x_i)_{i, j = 1, \ldots, c}$가
$B$에서 가역이다.
\end{lemma}

\begin{proof}
보조정리 \ref{smoothing-lemma-syntomic-complete-intersection}을 적용하여,
$C = R[x_1, \ldots, x_n]/(f_1, \ldots, f_c)$가 $R$ 위의
상대적 전역 완전교차가 되도록 수축 사상을 갖는 매끄러운
$R$-대수 사상 $A \to C$를 얻는다.
$C$가 $R$ 위에서 매끄러우므로 짧은 완전열
$$
0 \to
\bigoplus\nolimits_{j = 1, \ldots, c} C f_j \to
\bigoplus\nolimits_{i = 1, \ldots, n} C\text{d}x_i \to
\Omega_{C/R} \to 0
$$
이 있다. $\Omega_{C/R}$가 사영 $C$-가군이므로 이 완전열은 분할된다.
첫 번째 사상의 왼쪽 역 $t$를 택하자.
$t(\text{d}x_i) = \sum c_{ij} f_j$라고 쓰면
$\sum_i \frac{\partial f_j}{\partial x_i} c_{i\ell} = \delta_{j\ell}$
(크로네커 델타)이다.
$$
B' = C[y_1, \ldots, y_c] =
R[x_1, \ldots, x_n, y_1, \ldots, y_c]/(f_1, \ldots, f_c)
$$
로 놓자. $R$-대수 사상 $C \to B'$는 $y_j$를 영으로 보내는
수축 사상을 갖는다. 사상
$$
R[z_1, \ldots, z_n] \longrightarrow B',\quad
z_i \longmapsto x_i - \sum\nolimits_j c_{ij} y_j
$$
이 $\Spec(C) \to \Spec(B')$의 상에 속하는 모든 점에서 에탈임을
주장한다.\hfil\break
$\Omega_{B'/R[z_1, \ldots, z_n]}$에서
$$
0 =
\text{d}f_j - \sum\nolimits_i \frac{\partial f_j}{\partial x_i} \text{d}z_i
\equiv
\sum\nolimits_{i, \ell}
\frac{\partial f_j}{\partial x_i} c_{i\ell} \text{d}y_\ell
\equiv
\text{d}y_j \bmod (y_1, \ldots, y_c)\Omega_{B'/R[z_1, \ldots, z_n]}
$$
이다.
$\sum B'\text{d}y_j + (y_1, \ldots, y_c)\Omega_{B'/R[z_1, \ldots, z_n]}$
를 법으로 하여 $0 = \text{d}z_i = \text{d}x_i$이므로,
$$
\Omega_{B'/R[z_1, \ldots, z_n]}/
(y_1, \ldots, y_c)\Omega_{B'/R[z_1, \ldots, z_n]} = 0.
$$
을 얻는다. $\Omega_{B'/R[z_1, \ldots, z_n]}$는 유한 $B'$-가군이므로,
나카야마 보조정리에 의해
$(\Omega_{B'/R[z_1, \ldots, z_n]})_g = 0$이 되는
$g \in 1 + (y_1, \ldots, y_c)$가 존재한다.
이는 $R[z_1, \ldots, z_n] \to B'_g$가 비분기임을 증명한다.
Algebra, Definition \ref{algebra-definition-unramified}를 보라.
$k$가 체인 임의의 환 준동형 $R \to k$에 대하여, 차원이 $n$인
매끄러운 $k$-대수 사이의 비분기 환 준동형
$k[z_1, \ldots, z_n] \to (B'_g) \otimes_R k$를 얻는다.
Algebra, Lemmas \ref{algebra-lemma-CM-over-regular-flat} 및
\ref{algebra-lemma-characterize-smooth-kbar}에 의해,
$k[z_1, \ldots, z_n] \to (B'_g) \otimes_R k$는 평탄하다.
섬유별 평탄성 판정법
(Algebra, Lemma \ref{algebra-lemma-criterion-flatness-fibre})으로부터
$R[z_1, \ldots, z_n] \to B'_g$가 평탄함을 얻는다.
끝으로 Algebra, Lemma \ref{algebra-lemma-characterize-etale}에 의해
$R[z_1, \ldots, z_n] \to B'_g$는 에탈이다.
$B = B'_g$로 놓자. $C \to B$는 매끄럽고 수축 사상을 가지므로
$A \to B$도 매끄럽고 수축 사상을 갖는다.
또한 $R[z_1, \ldots, z_n] \to B$는 에탈이다.
Algebra, Lemma \ref{algebra-lemma-etale-standard-smooth}에 의해
$$
B = R[z_1, \ldots, z_n, w_1, \ldots, w_m]/(g_1, \ldots, g_m)
$$
으로 쓸 수 있으며, 여기서 $\det(\partial g_j/\partial w_i)$는
$B$에서 가역이다. 이로써 보조정리가 증명된다.
\end{proof}

\begin{lemma}
\label{smoothing-lemma-colimit-standard-smooth}
$R \to \Lambda$를 환 준동형이라 하자.
$\Lambda$가 매끄러운 $R$-대수들의 여과 여극한이면,
$\Lambda$는 표준 매끄러운 $R$-대수들의 여과 여극한이다.
\end{lemma}

\begin{proof}
$A$가 $R$ 위에서 유한 표시를 갖는 $R$-대수 사상
$A \to \Lambda$를 생각하자.
Algebra, Lemma \ref{algebra-lemma-when-colimit}에 따르면,
이 사상을 표준 매끄러운 대수를 거쳐 분해하면 충분하다.
또한 $B$가 $R$ 위에서 매끄럽도록
$A \to B \to \Lambda$로 분해할 수 있음을 알고 있다.
$C$가 $R$ 위에서 표준 매끄럽고 수축 사상 $C \to B$를 갖는
$R$-대수 사상 $B \to C$를 택하자.
보조정리 \ref{smoothing-lemma-smooth-standard-smooth}를 보라.
그러면 원하는 분해는 $A \to B \to C \to B \to \Lambda$이다.
\end{proof}

\begin{lemma}
\label{smoothing-lemma-standard-smooth-include-generators}
$R \to A$를 표준 매끄러운 환 준동형이라 하자.
$E \subset A$를 원소 수가 $|E| = n$인 유한 부분집합이라 하자.
그러면 $c \geq n$이고,
$\det(\partial f_j/\partial x_i)_{i, j = 1, \ldots, c}$가 $A$에서 가역이며,
$E$가 $x_1, \ldots, x_n$의 합동류들의 집합이 되는 표시
$A = R[x_1, \ldots, x_{n + m}]/(f_1, \ldots, f_c)$가 존재한다.
\end{lemma}

\begin{proof}
$\det(\partial g_j/\partial y_i)_{i, j = 1, \ldots, d}$의 상이
$A$에서 가역이 되도록\hfil\break
$A = R[y_1, \ldots, y_m]/(g_1, \ldots, g_d)$라는 표시를 택하자.
$E = \{a_1, \ldots, a_n\}$으로 원소들을 나열하고,
$A$에서의 상이 $a_i$인 $h_i \in R[y_1, \ldots, y_m]$를 택하자.
표시
$$
A = R[x_1, \ldots, x_n, y_1, \ldots, y_m]/
(x_1 - h_1, \ldots, x_n - h_n, g_1, \ldots, g_d)
$$
를 생각하고 $c = n + d$로 놓으면 된다.
\end{proof}

\begin{lemma}
\label{smoothing-lemma-compare-standard}
$R \to A$를 유한 표시 환 준동형이라 하자.
$a \in A$라 하자. $a$에 대해 다음 조건들을 생각하자.
\begin{enumerate}
\item $A_a$는 $R$ 위에서 매끄럽다.
\item $A_a$는 $R$ 위에서 매끄럽고 $\Omega_{A_a/R}$는 안정적으로 자유이다.
\item $A_a$는 $R$ 위에서 매끄럽고 $\Omega_{A_a/R}$는 자유이다.
\item $A_a$는 $R$ 위에서 표준 매끄럽다.
\item $a$는 $R$ 위의 $A$에서 엄밀 표준이다.
\item $a$는 $R$ 위의 $A$에서 초등 표준이다.
\end{enumerate}
그러면 다음이 성립한다.
\begin{enumerate}
\item[(a)] (4) $\Rightarrow$ (3) $\Rightarrow$ (2) $\Rightarrow$ (1).
\item[(b)] (6) $\Rightarrow$ (5).
\item[(c)] (6) $\Rightarrow$ (4).
\item[(d)] (5) $\Rightarrow$ (2).
\item[(e)] (2) $\Rightarrow$ $e \geq e_0$인 원소 $a^e$들은
$R$ 위의 $A$에서 엄밀 표준이다.
\item[(f)] (4) $\Rightarrow$ $e \geq e_0$인 원소 $a^e$들은
$R$ 위의 $A$에서 초등 표준이다.
\end{enumerate}
\end{lemma}

\begin{proof}
(a)는 정의들과 Algebra, Lemma \ref{algebra-lemma-standard-smooth}에서
명백하다. (b)는 정의 \ref{smoothing-definition-strictly-standard}에서 명백하다.

\medskip\noindent
(c)의 증명.
(\ref{smoothing-equation-elementary-standard-one})과
(\ref{smoothing-equation-elementary-standard-two})가 성립하도록 표시
$A = R[x_1, \ldots, x_n]/(f_1, \ldots, f_m)$를 택하자.
$a$로 보내지는 $h \in R[x_1, \ldots, x_n]$를 택하면
$$
A_a = R[x_0, x_1, \ldots, x_n]/(x_0h - 1, f_1, \ldots, f_m).
$$
이다. $J = (x_0h - 1, f_1, \ldots, f_m)$으로 쓰자.
(\ref{smoothing-equation-elementary-standard-two})에 의해 $A_a$-가군
$J/J^2$는 $A_a$ 위에서 $x_0h - 1, f_1, \ldots, f_c$로 생성된다.
따라서 Algebra, Lemma \ref{algebra-lemma-huber}의 증명에서와 같이,
$$
A_a = R[x_0, \ldots, x_n, x_{n + 1}]/
(x_0h - 1, f_1, \ldots, f_c, gx_{n + 1} - 1).
$$
이 되도록 $g \in 1 + J$를 택할 수 있다.
이제 (\ref{smoothing-equation-elementary-standard-one})에 의해
$R \to A_a$가 표준 매끄럽다
(좌표 $x_0, x_1, \ldots, x_c, x_{n + 1}$에 대해 미분한다).

\medskip\noindent
(d)의 증명.
(\ref{smoothing-equation-strictly-standard-one})과
(\ref{smoothing-equation-strictly-standard-two})가 성립하도록 표시
$A = R[x_1, \ldots, x_n]/(f_1, \ldots, f_m)$를 택하자.
$I = (f_1, \ldots, f_m)$으로 쓰자.
$A_a$가 $R$ 위에서 매끄럽다는 것은 이미 알고 있다.
보조정리 \ref{smoothing-lemma-elkik}를 보라.
보조정리 \ref{smoothing-lemma-parse-equation-strictly-standard-one}에 의해
$(I/I^2)_a$는 $f_1, \ldots, f_c$를 기저로 하는 자유 가군이고,
$\bigoplus A_a \text{d}x_i$의 한 직합인자로 동형으로 보내진다.
$\Omega_{A_a/R} = (\Omega_{A/R})_a$는 사상
$(I/I^2)_a \to \bigoplus A_a \text{d}x_i$의 여핵이므로
안정적으로 자유이다.

\medskip\noindent
(e)의 증명. $I$가 유한 생성인 표시 $A = R[x_1, \ldots, x_n]/I$를 택하자.
가정에 의해 분할 짧은 완전열
$$
0 \to (I/I^2)_a \to \bigoplus\nolimits_{i = 1, \ldots, n} A_a\text{d}x_i \to
\Omega_{A_a/R} \to 0
$$
이 있다. 따라서 $(I/I^2)_a \oplus \Omega_{A_a/R}$는 자유 $A_a$-가군이다.
$\Omega_{A_a/R}$가 안정적으로 자유이므로 $(I/I^2)_a$도 안정적으로
자유이다. 그러므로 어떤 $r$에 대해 위의 표시를
$A = R[x_1, \ldots, x_n, x_{n + 1}, \ldots, x_{n + r}]/J$,
$J = (I, x_{n + 1}, \ldots, x_{n + r})$로 바꾸면
$(J/J^2)_a$가 (유한) 자유 가군이 되게 할 수 있다.
$(J/J^2)_a$의 기저로 보내지는 $f_1, \ldots, f_c \in J$를 택하자.
이를 생성원 목록 $f_1, \ldots, f_m \in J$로 확장하자.
표시 $A = R[x_1, \ldots, x_{n + r}]/(f_1, \ldots, f_m)$를 생각하자.
그러면 구성에 의해 충분히 큰 모든 $e$에 대해
$a^e$는 (\ref{smoothing-equation-strictly-standard-two})를 만족한다.
또한
$(J/J^2)_a \to \bigoplus\nolimits_{i = 1, \ldots, n + r} A_a\text{d}x_i$
는 분할 단사 사상이므로 $A_a$-선형 왼쪽 역을 찾을 수 있다.
이 왼쪽 역을 기저 $f_1, \ldots, f_c$에 관해 쓰고 분모를 없애면,
어떤 $e_0 \geq 1$에 대해
$$
A^{\oplus c} \xrightarrow{(f_1, \ldots, f_c)}
J/J^2 \xrightarrow{f \mapsto \text{d}f}
\bigoplus\nolimits_{i = 1, \ldots, n + r} A \text{d}x_i
\xrightarrow{\psi_0}
A^{\oplus c}
$$
가 $a^{e_0}$을 곱하는 사상이 되게 하는 선형 사상
$\psi_0 : A^{\oplus n + r} \to A^{\oplus c}$를 얻는다.
보조정리 \ref{smoothing-lemma-parse-equation-strictly-standard-one}에 의해
모든 $a^{ce_0}$에 대해 (\ref{smoothing-equation-strictly-standard-one})이 성립하며,
따라서 $e \geq ce_0$인 모든 $e$에 대해 $a^e$에도 성립한다.

\medskip\noindent
(f)의 증명.
$\det(\partial f_j/\partial x_i)_{i, j = 1, \ldots, c}$가 $A_a$에서
가역이 되도록 표시 $A_a = R[x_1, \ldots, x_n]/(f_1, \ldots, f_c)$를
택하자. 어떤 $m < n$에 대해 $x_1, \ldots, x_m$의 동치류들이,
$R$ 위에서 $A$를 생성하는 $a_1, \ldots, a_m \in A$에 대한
$a_i/1$에 대응한다고 가정해도 된다.
보조정리 \ref{smoothing-lemma-standard-smooth-include-generators}를 보라.
$m < i \leq n$에 대해 $x_i$를 $a^Nx_i$로 바꾸면,
어떤 $a_i \in A$에 대해 $x_i$의 동치류가 $a_i/1 \in A_a$라고
가정할 수 있다. 환 준동형
$$
\Psi : R[x_1, \ldots, x_n] \longrightarrow A,\quad
x_i \longmapsto a_i.
$$
를 생각하자. 이는 전사 환 준동형이다.
$f_j$를 $a^Nf_j$로 바꾸면
$f_j \in R[x_1, \ldots, x_n]$이고 $\Psi(f_j) = 0$이라고
가정할 수 있다(실제로 $A_a$에서 $f_j(a_1/1, \ldots, a_n/1) = 0$이기
때문이다). $J = \Ker(\Psi)$라 하자.
그러면 $A = R[x_1, \ldots, x_n]/J$는 하나의 표시이며,
원소 $f_1, \ldots, f_c \in J$는 $(J/J^2)_a$가
$f_1, \ldots, f_c$를 기저로 하는 자유 가군이 되게 하고,
$\det(\partial f_j/\partial x_i)_{i, j = 1, \ldots, c}$는
$A_a$의 가역원으로 보내진다. 따라서 원하는 대로,
충분히 큰 모든 $e$에 대해 $a^e$는
(\ref{smoothing-equation-elementary-standard-one})과
(\ref{smoothing-equation-elementary-standard-two})를 만족한다.
\end{proof}






\section{막간: N\'eron 특이점 해소}
\label{smoothing-section-neron}

\noindent
Popescu 정리의 일반적인 경우에 대한 논의를 잠시 멈추고,
더 쉽지만 이미 매우 흥미로운 경우, 즉 $R \to \Lambda$가
이산 값매김환의 준동형인 경우를 살펴보자.
이 경우는 \cite[Section 4]{Artin-Algebraic-Approximation}에서 다룬다.

\begin{situation}
\label{smoothing-situation-neron}
여기서 $R \subset \Lambda$는 분기 지수가 $1$인 이산 값매김환의
확장이다(More on Algebra, Definition
\ref{more-algebra-definition-extension-discrete-valuation-rings}).
$R \to A$가 평탄하고 유한형인 분해
$$
R \to A \xrightarrow{\varphi} \Lambda
$$
가 주어져 있다고 가정한다.
$\mathfrak q = \Ker(\varphi)$ 및
$\mathfrak p = \varphi^{-1}(\mathfrak m_\Lambda)$로 놓자.
\end{situation}

\noindent
상황 \ref{smoothing-situation-neron}에서 $\pi \in R$를 균일화원이라 하자.
$A$가 $R$ 위에서 평탄하다는 것은 $\pi$가 $A$에서 비영인자라는
뜻임을 상기하자(More on Algebra, Lemma
\ref{more-algebra-lemma-valuation-ring-torsion-free-flat}).
$R \subset \Lambda$에 대한 가정에 의해 $\pi$는 $\Lambda$의
균일화원으로 보내진다. $\pi \in \mathfrak p$이므로
N\'eron의 아핀 블로업 대수
(Algebra, Section \ref{algebra-section-blow-up} 참조)
$$
\varphi' :
A' = A[\textstyle{\frac{\mathfrak p}{\pi}}]
\longrightarrow
\Lambda
$$
를 생각할 수 있다. 여기에는 $a \in \mathfrak p^n$인
$a/\pi^n$을 $\Lambda$의 $\pi^{-n}\varphi(a)$로 보내는
$\Lambda$로 가는 유도 사상이 주어진다.
$A'$의 대응하는 소아이디얼들을
$\mathfrak q' \subset \mathfrak p' \subset A'$로 표기하자.
$A'$의 동형류는 균일화원의 선택에 의존하지 않음에 유의하자.
이 구성을 반복하면 열
$$
A \to A' \to A'' \to \ldots \to \Lambda
$$
을 얻는다.

\begin{lemma}
\label{smoothing-lemma-neron-functorial}
상황 \ref{smoothing-situation-neron}에서 N\'eron의 블로업은 다음 의미에서
함자적이다.
\begin{enumerate}
\item $a \in A$, $a \not \in \mathfrak p$이면,
$A_a$의 N\'eron 블로업은 $A'_a$이다.
\item $B \to A$가 핵이 $I$인 평탄 유한형 $R$-대수의 전사 사상이면,
$A'$는 $B'/IB'$를 그 $\pi$-거듭제곱 꼬임으로 나눈 몫이다.
\end{enumerate}
\end{lemma}

\begin{proof}
(1)과 (2)는 모두 Algebra, Lemma
\ref{algebra-lemma-blowup-base-change}의 특수한 경우이다.
실제로 $A_1 \to A_2 \to \Lambda$가
$\mathfrak p_1 A_2 = \mathfrak p_2$를 만족하면,
$A_2'$는 $A_1' \otimes_{A_1} A_2$를 그 $\pi$-거듭제곱 꼬임으로
나눈 몫이다.
\end{proof}

\begin{lemma}
\label{smoothing-lemma-neron-blowup-smooth}
상황 \ref{smoothing-situation-neron}에서 $R \to A$가 $\mathfrak p$에서
매끄럽고 $R/\pi R \subset \Lambda/\pi \Lambda$가 분리 체확장이라고
가정하자. 그러면 $R \to A'$는 $\mathfrak p'$에서 매끄러우며,
짧은 완전열
$$
0 \to
\Omega_{A/R} \otimes_A A'_{\mathfrak p'} \to
\Omega_{A'/R, \mathfrak p'} \to
(A'/\pi A')_{\mathfrak p'}^{\oplus c} \to 0
$$
이 존재한다. 여기서 $c = \dim((A/\pi A)_\mathfrak p)$이다.
\end{lemma}

\begin{proof}
보조정리 \ref{smoothing-lemma-neron-functorial}에 의해 $A$를 $\mathfrak p$에
속하지 않는 원소에 대한 국소화로 바꿔도 된다.
이하에서는 이를 별도 언급 없이 사용한다.
$\kappa = R/\pi R$로 쓰자.
매끄러움은 밑변환에 대해 안정적이므로
(Algebra, Lemma \ref{algebra-lemma-base-change-smooth}),
$A/\pi A$는 $\mathfrak p$에서 $\kappa$ 위로 매끄럽다.
따라서 $(A/\pi A)_\mathfrak p$는 정칙 국소환이다
(Algebra, Lemma \ref{algebra-lemma-characterize-smooth-over-field}).
$(A/\pi A)_\mathfrak p$의 정칙 매개변수계로 보내지는
$g_1, \ldots, g_c \in \mathfrak p$를 택하자.
필요하면 $A$를 국소화로 바꾸어
$\mathfrak p = (\pi, g_1, \ldots, g_c)$라고 할 수 있다.
$\pi, g_1, \ldots, g_c$는 $A_\mathfrak p$의 정칙열임에 유의하자
(먼저 $\pi$는 비영인자이고, 나머지 열에는
Algebra, Lemma \ref{algebra-lemma-regular-ring-CM}을 적용한다).
$A$를 국소화로 바꾸면 $\pi, g_1, \ldots, g_c$가 $A$의 정칙열이라고
가정할 수 있다
(Algebra, Lemma \ref{algebra-lemma-regular-sequence-in-neighbourhood}).
More on Algebra, Lemma \ref{more-algebra-lemma-blowup-regular-sequence}에
의해
$$
A' = A[y_1, \ldots, y_c]/(\pi y_1 - g_1, \ldots, \pi y_c - g_c) =
A[y_1, \ldots, y_c]/I
$$
이다. 이하에서는 소박한 여접복합체를 이용한 매끄러움의 정의
(Algebra, Definition \ref{algebra-definition-smooth})와
Algebra, Lemma \ref{algebra-lemma-smooth-at-point}의 판정법을
별도 언급 없이 사용한다.
$R \to A[y_1, \ldots, y_c] \to A'$에 대한
Algebra, Lemma \ref{algebra-lemma-exact-sequence-NL}의 완전열은
다음과 같다.
$$
0 \to H_1(\NL_{A'/R}) \to I/I^2 \to
\Omega_{A/R} \otimes_A A' \oplus
\bigoplus\nolimits_{i = 1, \ldots, c} A' \text{d}y_i \to
\Omega_{A'/R} \to 0
$$
여기서 $I/I^2$에 속하는 $\pi y_i - g_i$의 동치류는 다음 항의
$- \text{d}g_i + \pi \text{d}y_i$로 보내진다.
이때 $\NL_{A'/A[y_1, \ldots, y_c]}$를 계산하기 위해
Algebra, Lemma \ref{algebra-lemma-NL-surjection}을 사용하였다.
또한 $R \to A[y_1, \ldots, y_c]$가 매끄럽다는 사실도 사용하였다.
따라서 $H_1(\NL_{A[y_1, \ldots, y_c]/R}) = 0$이며,
$\Omega_{A[y_1, \ldots, y_c]/R}$는 유한 사영(따라서 평탄)
$A[y_1, \ldots, y_c]$-가군이다. 실제로 이 가군은
$\Omega_{A/R} \otimes_A A[y_1, \ldots, y_c]$와
$\text{d}y_i$를 기저로 하는 자유 가군의 직합이다.
증명을 마치려면 $\text{d}g_1, \ldots, \text{d}g_c$가
유한 자유 가군 $\Omega_{A/R, \mathfrak p}$의 기저의 일부를 이룸을
보이면 충분하다. 그러면 $(I/I^2)_\mathfrak p$는
$\pi y_i - g_i$를 기저로 하는 자유 가군이고, 완전열 가운데 사상의
$\mathfrak p$에서의 국소화는 단사이므로
$H_1(\NL_{A'/R})_\mathfrak p = 0$이며,
여핵 $\Omega_{A'/R, \mathfrak p}$도 유한 자유 가군임을 알 수 있다.
이를 위해서는 $\text{d}g_i$의 상들이
$\Omega_{A/R, \mathfrak p}/\pi = \Omega_{(A/\pi A)/\kappa, \mathfrak p}$
에서 $\kappa(\mathfrak p)$-선형 독립임을 보이면 충분하다
(이 등식은 Algebra, Lemma
\ref{algebra-lemma-differentials-base-change}에 의한다).
$\kappa \subset \kappa(\mathfrak p) \subset \Lambda/\pi \Lambda$이므로,
$\kappa(\mathfrak p)$는 $\kappa$ 위에서 분리적이다
(Algebra, Definition \ref{algebra-definition-separable-field-extension}).
이제 원하는 선형 독립성은
Algebra, Lemma \ref{algebra-lemma-computation-differential}에서 따른다.
\end{proof}

\begin{lemma}
\label{smoothing-lemma-neron-when-smooth}
상황 \ref{smoothing-situation-neron}에서 $R \to A$가 $\mathfrak q$에서
매끄럽고, 핵이 $I$인 $R$-대수의 전사 사상 $B \to A$가 있다고
가정하자. $R \to B$가
$\mathfrak p_B = (B \to A)^{-1}\mathfrak p$에서 매끄럽다고 가정하자.
$$
I/I^2 \otimes_A \Lambda \to \Omega_{B/R} \otimes_B \Lambda
$$
의 여핵이 자유 $\Lambda$-가군이면,
$R \to A$는 $\mathfrak p$에서 매끄럽다.
\end{lemma}

\begin{proof}
사상 $I/I^2 \to \Omega_{B/R} \otimes_B A$의 여핵은 $\Omega_{A/R}$이다.
Algebra, Lemma \ref{algebra-lemma-differential-seq}를 보라.
$d = \dim_\mathfrak q(A/R)$를 $\mathfrak q$에서의 $R \to A$의
상대차원, 즉 $\mathfrak q$에서의 $\Spec(A[1/\pi])$의 차원이라 하자.
Algebra, Definition \ref{algebra-definition-relative-dimension}을 보라.
그러면 $\Omega_{A/R, \mathfrak q}$는 $A_\mathfrak q$ 위에서
계수가 $d$인 자유 가군이다
(Algebra, Lemma \ref{algebra-lemma-characterize-smooth-over-field}).
따라서 보조정리의 가정이 성립하면
$\Omega_{A/R} \otimes_A \Lambda$는 계수가 $d$인 자유 가군이다.
그러므로 $\Omega_{A/R} \otimes_A \kappa(\mathfrak p)$의 차원은 $d$이다
($\Lambda/\pi \Lambda$와 텐서곱한 뒤에 이 사실이 성립하기 때문이다).
$R \to A$가 평탄하고 $\mathfrak p$가 $\mathfrak q$의 특수화이므로,
Algebra, Lemma
\ref{algebra-lemma-dimension-fibres-bounded-open-upstairs}에 의해
$\dim_\mathfrak p(A/R) \geq d$이다.
그러면 Algebra, Lemmas \ref{algebra-lemma-flat-fibre-smooth} 및
\ref{algebra-lemma-characterize-smooth-over-field}에 의해
$R \to A$가 $\mathfrak p$에서 매끄럽다.
\end{proof}

\begin{lemma}
\label{smoothing-lemma-neron-desingularization}
상황 \ref{smoothing-situation-neron}에서 $R \to A$가 $\mathfrak q$에서
매끄럽고 $R/\pi R \subset \Lambda/\pi \Lambda$가 분리 체확장이라고
가정하자. 그러면 유한 번의 아핀 N\'eron 블로업을 시행한 뒤에는
대수 $A$가 $\mathfrak p$에서 $R$ 위로 매끄러워진다.
\end{lemma}

\begin{proof}
$R$-대수 $B$와 전사 사상 $B \to A$를 택하자.
$\mathfrak p_B = (B \to A)^{-1}(\mathfrak p)$로 놓고,
$\mathfrak p_B$에서의 $R \to B$의 상대차원을 $r$로 표기하자.
$R \to B$가 $\mathfrak p_B$에서 매끄럽도록 $B$를 택한다.
예를 들어 $B$를 $R$ 위의 $r$변수 다항식 대수로 택할 수 있다.
보조정리 \ref{smoothing-lemma-neron-when-smooth}의 복합체
$$
I/I^2 \otimes_A \Lambda \longrightarrow \Omega_{B/R} \otimes_B \Lambda
$$
를 생각하자. $\Lambda$ 위의 유한 가군의 구조
(More on Algebra, Lemma \ref{more-algebra-lemma-modules-PID})에 의해,
여핵은 어떤 $d \geq 0$, $n \geq 0$, $e_i \geq 1$에 대해
$$
\Lambda^{\oplus d} \oplus
\bigoplus\nolimits_{i = 1, \ldots, n} \Lambda/\pi^{e_i} \Lambda
$$
의 형태임을 알 수 있다.
$d$는 $\mathfrak q$에서의 $A/R$의 상대차원임에 유의하자
(Algebra, Lemma \ref{algebra-lemma-characterize-smooth-over-field}).
결함량 $e = \sum_{i = 1, \ldots, n} e_i$가 영이면,
보조정리 \ref{smoothing-lemma-neron-when-smooth}에 의해 증명이 끝난다.

\medskip\noindent
이제 N\'eron 블로업을 시행하면 어떻게 되는지 살펴보자.
$A'$는 $B'/IB'$를 그 $\pi$-거듭제곱 꼬임으로 나눈 몫이며
(보조정리 \ref{smoothing-lemma-neron-functorial}),
$R \to B'$는 $\mathfrak p_{B'}$에서 매끄럽다
(보조정리 \ref{smoothing-lemma-neron-blowup-smooth}).
따라서 블로업 뒤에도 정확히 같은 상황이 성립한다.
그림으로 나타내면 다음과 같다.
$$
\xymatrix{
0 \ar[r] & I' \ar[r] & B' \ar[r] & A' \ar[r] & 0 \\
0 \ar[r] & I \ar[u] \ar[r] & B \ar[u] \ar[r] & A \ar[r] \ar[u] & 0
}
$$
$I \subset \mathfrak p_B$이므로 $I \to I'$는 $\pi I'$를 거쳐 분해된다.
유도되는 복합체 사상을 살펴보면
$$
\xymatrix{
I'/(I')^2 \otimes_{A'} \Lambda \ar[r] &
\Omega_{B'/R} \otimes_{B'} \Lambda \ar@{=}[r] & M' \\
I/I^2 \otimes_A \Lambda \ar[r] \ar[u] &
\Omega_{B/R} \otimes_B \Lambda \ar[u] \ar@{=}[r] & M
}
$$
를 얻는다. 그러면 $M \subset M'$는 유한 자유 $\Lambda$-가군이며,
몫 $M'/M$은 $\pi$에 의해 소멸한다.
보조정리 \ref{smoothing-lemma-neron-blowup-smooth}를 보라.
수평 사상들의 상을 $N \subset M$과 $N' \subset M'$로 놓고,
$Q = M/N$ 및 $Q' = M'/N'$로 표기하자.
가환 그림
$$
\xymatrix{
0 \ar[r] &
N' \ar[r] &
M' \ar[r] &
Q' \ar[r] &
0 \\
0 \ar[r] &
N \ar[r] \ar[u] &
M \ar[r] \ar[u] &
Q \ar[r] \ar[u] &
0
}
$$
을 얻는다. 그러면 $N \subset N'$는 계수가 $r - d$인 자유
$\Lambda$-가군이다.
$I$가 $\pi I'$ 안으로 보내지므로 $N \subset \pi N'$이다.

\medskip\noindent
$K = \Lambda_\pi$를 $\Lambda$의 분수체라 하자.
행들이 짧은 완전열인 가환 그림
$$
\xymatrix{
0 \ar[r] &
N' \ar[r] &
N'_K \cap M' \ar[r] &
Q'_{tor} \ar[r] &
0 \\
0 \ar[r] &
N \ar[r] \ar[u] &
N_K \cap M \ar[r] \ar[u] &
Q_{tor} \ar[r] \ar[u] &
0
}
$$
이 있다. 따라서 결함량의 변화는
$$
e - e' =
\text{length}(Q_{tor}) - \text{length}(Q'_{tor})
=
\text{length}(N'/N) - \text{length}(N'_K \cap M' / N_K \cap M)
$$
로 주어진다.
$M'/M$이 $\pi$에 의해 소멸하므로
$N'_K \cap M' / N_K \cap M$도 그러하며,
그 길이는 $\dim_K(N_K)$ 이하이다.
$N \subset \pi N'$이므로
$\text{length}(N'/N) \ge \dim_K(N_K)$이며,
등호가 성립할 필요충분조건은 $N = \pi N'$이다.

\medskip\noindent
증명을 마치려면 $A$가 $\mathfrak p$에서 매끄럽지 않을 때
$N$이 $\pi N'$의 진부분가군임을 보여야 한다.
이것이 N\'eron의 논증에서 수행해야 할 핵심 계산이다.
이를 위해 완전열
$$
I/I^2 \otimes_B \kappa(\mathfrak p_B)
\to \Omega_{B/R} \otimes_B \kappa(\mathfrak p_B)
\to \Omega_{A/R} \otimes_A \kappa(\mathfrak p) \to 0
$$
을 생각하자(Algebra, Lemma
\ref{algebra-lemma-differential-seq}에서 따른다).
$R \to A$가 $\mathfrak p$에서 매끄럽지 않으므로,
$\Omega_{A/R} \otimes_A \kappa(\mathfrak p)$의 차원 $s$는 $d$보다 크다.
한편 첫 번째 화살표는
Algebra, Lemma \ref{algebra-lemma-computation-differential}의 단사 사상
$$
\mathfrak p B_\mathfrak p/\mathfrak p^2 B_\mathfrak p
\to \Omega_{B/R} \otimes_B \kappa(\mathfrak p_B)
$$
을 거쳐 분해된다.
$R/\pi R \subset \Lambda/\pi \Lambda$에 대한 가정에 의해
$\kappa(\mathfrak p)$가 $k$ 위에서 분리적임에 유의하자.
따라서 $j > r - s$일 때 $g_j \in \mathfrak p^2$가 되는 생성원들
$g_1, \ldots, g_t \in I$를 찾을 수 있다.
그러면 $j > r - s$일 때 $g_j$의 $A'$에서의 상은 $\pi^2 I'$에 속한다.
$r - s < r - d$이므로, $N$의 최소 생성원들 중 적어도 하나는
$N'$에서 $\pi^2$로 나누어떨어진다.
따라서 $e$는 적어도 $1$만큼 감소하며, 증명이 끝난다.
\end{proof}

\noindent
$R \to \Lambda$가 이산 값매김환의 확장이면,
$R \to \Lambda$가 정칙일 필요충분조건은
(a) 분기 지수가 $1$이고,
(b) 분수체의 확장이 분리적이며,
(c) $R/\mathfrak m_R \subset \Lambda/\mathfrak m_\Lambda$가
분리적인 것이다.
따라서 다음 결과는 정리 \ref{smoothing-theorem-popescu}의
일반 N\'eron 특이점 해소의 특수한 경우이다.

\begin{lemma}
\label{smoothing-lemma-neron-colimit}
\begin{slogan}
이산 값매김환의 비분기 확장은 ind-매끄럽다: N\'eron 특이점 해소
\end{slogan}
$R \subset \Lambda$를 분기 지수가 $1$이고 잉여체와 분수체의
분리 확장을 유도하는 이산 값매김환의 확장이라 하자.
그러면 $\Lambda$는 매끄러운 $R$-대수들의 여과 여극한이다.
\end{lemma}

\begin{proof}
Algebra, Lemma \ref{algebra-lemma-when-colimit}에 의해,
상황 \ref{smoothing-situation-neron}의 임의의 $R \to A \to \Lambda$가
매끄러운 $R$-대수 $B$를 거쳐 $A \to B \to \Lambda$로
분해됨을 보이면 충분하다.
$A$를 $\Lambda$에서의 상으로 바꾸면,
$A$는 정역이며 그 분수체 $K$가 $\Lambda$의 분수체의 부분체라고
가정할 수 있다.
특히 가정에 의해 $A$는 $R$의 분수체 위에서 분리적이다.
그러면 Algebra, Lemma \ref{algebra-lemma-smooth-at-generic-point}에 의해
$R \to A$는 $\mathfrak q = (0)$에서 매끄럽다.
유한 번의 N\'eron 블로업 뒤에는 $R \to A$가 $\mathfrak p$에서
매끄럽다고 가정할 수 있다.
보조정리 \ref{smoothing-lemma-neron-desingularization}을 보라.
이제 $A$를 $a \in A$, $a \not \in \mathfrak p$인 원소에서
국소화하면 $R$ 위에서 매끄러워지므로 보조정리가 증명된다.
\end{proof}












\section{올림 문제}
\label{smoothing-section-lifting}

\noindent
이 절의 목표는 매끄러운 대수들의 여과 여극한인 대수들의 모임이
무한소 평탄 변형에 대해 닫혀 있음을 증명하는 것이다
(명제 \ref{smoothing-proposition-lift}).
증명은 초등적이며, \ref{smoothing-section-presentations}절의
매끄러운 대수의 표시에 관한 결과들만을 사용한다.

\begin{lemma}
\label{smoothing-lemma-lift-once}
$R \to \Lambda$를 환 준동형, $I \subset R$를 아이디얼이라 하자.
다음을 가정하자.
\begin{enumerate}
\item $I^2 = 0$이다.
\item $\Lambda/I\Lambda$는 매끄러운 $R/I$-대수들의 여과 여극한이다.
\end{enumerate}
$A$가 $R$ 위에서 유한 표시를 갖는 $R$-대수 사상
$\varphi : A \to \Lambda$를 생각하자.
그러면 분해
$$
A \to B/J \to \Lambda
$$
가 존재한다. 여기서 $B$는 매끄러운 $R$-대수이고,
$J \subset IB$는 유한 생성 아이디얼이다.
\end{lemma}

\begin{proof}
$\bar B$가 $R/I$ 위에서 표준 매끄럽도록 분해
$$
A/IA \to \bar B \to \Lambda/I\Lambda
$$
를 택하자. 가정과 보조정리 \ref{smoothing-lemma-colimit-standard-smooth}에
의해 이것이 가능하다.
$$
\bar B = A/IA[t_1, \ldots, t_r]/(\bar g_1, \ldots, \bar g_s)
$$
로 쓰고, $\bar B \to \Lambda/I\Lambda$가 $t_i$를
$I\Lambda$를 법으로 한 $\lambda_i$의 동치류로 보낸다고 하자.
$\bar g_1, \ldots, \bar g_s$를 올리는
$g_1, \ldots, g_s \in A[t_1, \ldots, t_r]$를 택하자.
어떤 $\epsilon_{ij} \in I$와 $\mu_{ij} \in \Lambda$에 대해
$\varphi(g_i)(\lambda_1, \ldots, \lambda_r) =
\sum \epsilon_{ij} \mu_{ij}$로 쓰자.
$$
A' = A[t_1, \ldots, t_r, \delta_{i, j}]/
(g_i - \sum \epsilon_{ij} \delta_{ij})
$$
로 정의하고 사상
$$
A' \longrightarrow \Lambda,\quad
a \longmapsto \varphi(a),\quad
t_i \longmapsto \lambda_i,\quad
\delta_{ij} \longmapsto \mu_{ij}
$$
을 생각하자.
$$
A'/IA' = A/IA[t_1, \ldots, t_r]/(\bar g_1, \ldots, \bar g_s)[\delta_{ij}]
\cong \bar B[\delta_{ij}]
$$
이다. $\bar B$가 표준 매끄러우므로 이 대수는 $R/I$ 위에서
표준 매끄럽다.\hfil\break
$\det(\partial \bar f_j/\partial x_i)_{i, j = 1, \ldots, c}$가
$A'/IA'$에서 가역이 되는 표시\hfil\break
$A'/IA' =\allowbreak R/I[x_1, \ldots, x_n]/(\bar f_1, \ldots, \bar f_c)$를
택하자. $\bar f_1, \ldots, \bar f_c$의 올림
$f_1, \ldots, f_c \in R[x_1, \ldots, x_n]$를 택하자. 그러면
$$
B = R[x_1, \ldots, x_n, x_{n + 1}]/
(f_1, \ldots, f_c,
x_{n + 1}\det(\partial f_j/\partial x_i)_{i, j = 1, \ldots, c} - 1)
$$
는 $R$ 위에서 매끄럽다.
매끄러운 환 준동형은 형식적으로 매끄러우므로
(Algebra, Proposition \ref{algebra-proposition-smooth-formally-smooth}),
$I$를 법으로 하여 동형인 $R$-대수 사상 $B \to A'$가 존재한다.
그러면 나카야마 보조정리
(Algebra, Lemma \ref{algebra-lemma-NAK})에 의해 $B \to A'$는 전사이다.
따라서 $A' = B/J$이며, 여기서 $J \subset IB$는 유한 생성이다
(Algebra, Lemma \ref{algebra-lemma-finite-presentation-independent} 참조).
\end{proof}

\begin{lemma}
\label{smoothing-lemma-lift-twice}
$R \to \Lambda$를 환 준동형, $I \subset R$를 아이디얼이라 하자.
다음을 가정하자.
\begin{enumerate}
\item $I^2 = 0$이다.
\item $\Lambda/I\Lambda$는 매끄러운 $R/I$-대수들의 여과 여극한이다.
\item $R \to \Lambda$는 평탄하다.
\end{enumerate}
$B$가 $R$ 위에서 매끄러운 $R$-대수 사상
$\varphi : B \to \Lambda$를 생각하자.
$J \subset IB$를 $\varphi(J) = 0$인 유한 생성 아이디얼이라 하자.
그러면 $R$-대수 사상들
$$
B \xrightarrow{\alpha} B' \xrightarrow{\beta} \Lambda
$$
이 존재하여, $B'$는 $R$ 위에서 매끄럽고 $\alpha(J) = 0$이며
$\beta \circ \alpha = \varphi$이다.
\end{lemma}

\begin{proof}
$J = (h)$인 경우에 보조정리를 증명할 수 있으면,
$J$의 생성원 수에 대한 귀납법으로 일반적인 경우도 증명할 수 있다.
즉, $J$가 $n$개의 원소 $h_1, \ldots, h_n$으로 생성될 수 있고,
$J$가 $n - 1$개의 원소로 생성되는 모든 경우에 보조정리가 성립한다고
가정하자. $n = 1$인 경우를 적용하여 첫 번째 사상이 $h_n$을 영으로
보내는 $B \to B' \to \Lambda$를 얻는다.
이제 $J'$를 $h_1, \ldots, h_{n - 1}$의 상들로 생성되는
$B'$의 아이디얼로 놓고, $n - 1$인 경우를 적용하여
$B' \to B'' \to \Lambda$를 얻는다.
$B \to B'' \to \Lambda$가 원하는 조건들을 만족함은 쉽게 확인할 수 있다.

\medskip\noindent
$J = (h)$라고 가정하고, 어떤 $\epsilon_i \in I$와 $b_i \in B$에 대해
$h = \sum \epsilon_i b_i$로 쓰자.
$0 = \varphi(h) = \sum \epsilon_i \varphi(b_i)$임에 유의하자.
$\Lambda$가 $R$ 위에서 평탄하므로, 평탄성의 방정식 판정법
(Algebra, Lemma \ref{algebra-lemma-flat-eq})에 의해
$\varphi(b_i) = \sum_j a_{ij} \lambda_j$와
$\sum_i \epsilon_i a_{ij} = 0$을 만족하는
$\lambda_j \in \Lambda$, $j = 1, \ldots, m$ 및 $a_{ij} \in R$를
찾을 수 있다.
$$
C = B[x_1, \ldots, x_m]/(b_i - \sum a_{ij} x_j)
$$
로 놓고, $C \to \Lambda$는 $\varphi$와 $x_j \mapsto \lambda_j$로
정의하자. 보조정리 \ref{smoothing-lemma-lift-once}에서와 같은 분해
$$
C \to B'/J' \to \Lambda
$$
를 택하자. $B$가 $R$ 위에서 매끄러우므로 사상
$B \to C \to B'/J'$를 사상 $\alpha : B \to B'$로 올릴 수 있다.
그러면 $\varphi = \beta \circ \alpha$이다.
증명을 마치기 위해 $\alpha(h) = 0$을 확인하자.
$\alpha$가 $B \to C \to B'/J'$의 올림이라는 사실은,
어떤 $\xi_j \in B'$와 $\theta_i \in J' \subset IB'$에 대해
$$
\alpha(b_i) = \sum a_{ij} \xi_j + \theta_i
$$
임을 뜻한다. 따라서 위의 관계식들과 $I^2 = 0$이라는 사실에 의해
$$
\alpha(h) = \alpha(\sum \epsilon_i b_i) =
\sum \epsilon_i a_{ij} \xi_j + \sum \epsilon_i \theta_i = 0
$$
이다.
\end{proof}

\begin{proposition}
\label{smoothing-proposition-lift}
\begin{slogan}
대수의 ind-매끄러움은 무한소 변형에 대해 안정적이다
\end{slogan}
$R \to \Lambda$를 환 준동형, $I \subset R$를 아이디얼이라 하자.
다음을 가정하자.
\begin{enumerate}
\item $I$는 멱영이다.
\item $\Lambda/I\Lambda$는 매끄러운 $R/I$-대수들의 여과 여극한이다.
\item $R \to \Lambda$는 평탄하다.
\end{enumerate}
그러면 $\Lambda$는 매끄러운 $R$-대수들의 여과 여극한이다.
\end{proposition}

\begin{proof}
어떤 $n$에 대해 $I^n = 0$이므로, $n$에 대한 귀납법에 의해
$I^2 = 0$인 경우만 생각하면 충분하다.
$A$가 $R$ 위에서 유한 표시를 갖는 $R$-대수 사상
$\varphi : A \to \Lambda$를 생각하자.
$B$가 $R$ 위에서 매끄러운 분해 $A \to B \to \Lambda$를 찾아야 한다.
Algebra, Lemma \ref{algebra-lemma-when-colimit}를 보라.
보조정리 \ref{smoothing-lemma-lift-once}에 의해
$A = B/J$이며 $B$가 $R$ 위에서 매끄럽고
$J \subset IB$가 유한 생성 아이디얼이라고 가정할 수 있다.
보조정리 \ref{smoothing-lemma-lift-twice}에 의해,
$B'$가 $R$ 위에서 매끄럽고 $\alpha(J) = 0$인
$R$-대수의 가환 그림
$$
\xymatrix{
B \ar[rr]_\alpha \ar[rd]_\varphi & & B' \ar[ld]^\beta \\
& \Lambda
}
$$
을 찾을 수 있다.
따라서 $\alpha$는 $B \to A \to B'$로 분해되며, 증명이 끝난다.
\end{proof}






\section{올림 보조정리}
\label{smoothing-section-lifting-lemma}

\noindent
다음은 지독할 만큼 영리한 보조정리이다.

\begin{lemma}
\label{smoothing-lemma-lifting}
$R$를 뇌터 환, $\Lambda$를 $R$-대수라 하자.
$\pi \in R$라 하고
$\text{Ann}_R(\pi) = \text{Ann}_R(\pi^2)$ 및
$\text{Ann}_\Lambda(\pi) = \text{Ann}_\Lambda(\pi^2)$라고 가정하자.
$\bar C$가 유한 표시를 갖는 $R$-대수 사상들
$R/\pi^2R \to \bar C \to \Lambda/\pi^2\Lambda$가 주어졌다고 하자.
그러면 $R$-대수 준동형 $D \to \Lambda$와 가환 그림
$$
\xymatrix{
R/\pi^2R \ar[r] \ar[d] &
\bar C \ar[r] \ar[d] &
\Lambda/\pi^2\Lambda \ar[d] \\
R/\pi R \ar[r] &
D/\pi D \ar[r] &
\Lambda/\pi \Lambda
}
$$
이 존재하며, 다음 성질들을 만족한다.
\begin{enumerate}
\item[(a)] $D$는 유한 표시를 갖는다.
\item[(b)] $\pi \not \in \mathfrak q$인 모든 소아이디얼
$\mathfrak q$에서 $R \to D$는 매끄럽다.
\item[(c)] $\pi \in \mathfrak q$이고,
$R/\pi^2 R \to \bar C$가 매끄러운 $\bar C$의 소아이디얼 위에 놓이는
모든 소아이디얼 $\mathfrak q$에서 $R \to D$는 매끄럽다.
\item[(d)] $R/\pi^2R \to \bar C$가 매끄러운 $\bar C$의
소아이디얼 위에 놓이는 모든 소아이디얼에서
$\bar C/\pi \bar C \to D/\pi D$는 매끄럽다.
\end{enumerate}
\end{lemma}

\begin{proof}
표시
$$
\bar C = R[x_1, \ldots, x_n]/(f_1, \ldots, f_m)
$$
를 택하자.
또한 $I = (f_1, \ldots, f_m)$으로 표기하고,
$R/\pi^2R[x_1, \ldots, x_n]$에서의 $I$의 상을 $\bar I$로 표기한다.
$R$가 뇌터 환이므로 $\bar C$도 뇌터 환이다.
따라서 $R/\pi^2 R \to \bar C$의 매끄러운 자취는 준콤팩트이다.
Topology, Lemma \ref{topology-lemma-Noetherian}을 보라.
보조정리 \ref{smoothing-lemma-find-strictly-standard}를 적용하여
유한한 원소 목록 $a_1, \ldots, a_r \in R[x_1, \ldots, x_n]$을
다음과 같이 택할 수 있다.
\begin{enumerate}
\item 열린 부분공간들
$\Spec(\bar C_{a_k}) \subset \Spec(\bar C)$의 합집합은
$R/\pi^2 R \to \bar C$의 매끄러운 자취를 덮는다.
\item 각 $k = 1, \ldots, r$에 대해 유한 부분집합
$E_k \subset \{1, \ldots, m\}$가 존재하여,
$(\bar I/\bar I^2)_{a_k}$는 $j \in E_k$인 $f_j$들의 동치류들을
기저로 하는 자유 가군이다.
\end{enumerate}
$I_k = (f_j, j \in E_k) \subset I$로 놓고,
$R/\pi^2R[x_1, \ldots, x_n]$에서의 $I_k$의 상을 $\bar I_k$로 표기하자.
(2)와 나카야마 보조정리에 의해,
어떤 $b'_k \in \bar I_{a_k}$에 대해
$(\bar I/\bar I_k)_{a_k}$는 $1 + b'_k$에 의해 소멸한다.
어떤 $b_k \in I$와 정수 $N$에 대해 $b'_k$가 $b_k/(a_k)^N$의
상이라고 하자. $a_k$를 $a_kb_k$로 바꾸면
\begin{enumerate}
\item[(3)] $(\bar I_k)_{a_k} = (\bar I)_{a_k}$이다.
\end{enumerate}
따라서 필요하면 $a_k$를 높은 거듭제곱으로 바꾸어,
임의의 $\ell \in \{1, \ldots, m\}$에 대해
\begin{enumerate}
\item[(4)]
$a_k f_\ell = \sum\nolimits_{j \in E_k} h_{k, \ell}^jf_j + \pi^2 g_{k, \ell}$
\end{enumerate}
로 쓸 수 있다. 여기서 어떤
$h_{i, \ell}^j, g_{i, \ell} \in R[x_1, \ldots, x_n]$를 택한다.
$\ell \in E_k$이면
$h_{k, \ell}^j = a_k\delta_{\ell, j}$ (크로네커 델타) 및
$g_{k, \ell} = 0$으로 택한다.
$$
D = R[x_1, \ldots, x_n, z_1, \ldots, z_m]/
(f_j - \pi z_j, p_{k, \ell}).
$$
로 놓자.
여기서 $j \in \{1, \ldots, m\}$,
$k \in \{1, \ldots, r\}$,
$\ell \in \{1, \ldots, m\}$이고,
$$
p_{k, \ell} = a_k z_\ell - \sum\nolimits_{j \in E_k} h_{k, \ell}^j z_j
- \pi g_{k, \ell}.
$$
이다. 위의 선택에 의해 $\ell \in E_k$이면
$p_{k, \ell} = 0$임에 유의하자.

\medskip\noindent
사상 $R \to D$는 주어진 사상이다.
$\bar C \to \Lambda/\pi^2\Lambda$가 $x_i$를
$\pi^2$를 법으로 한 $\lambda_i$의 동치류로 보낸다고 하자.
$f \in R[x_1, \ldots, x_n]$에 대해 $x_i$ 자리에 $\lambda_i$를
대입한 결과를 $f(\lambda) \in \Lambda$로 표기한다.
그러면 어떤 $\mu_j \in \Lambda$에 대해
$f_j(\lambda) = \pi^2 \mu_j$임을 안다.
$x_i \mapsto \lambda_i$ 및 $z_j \mapsto \pi\mu_j$로
$D \to \Lambda$를 정의하자. 이는 다음 식 때문에 잘 정의된다.
\begin{align*}
p_{k, \ell} & \mapsto
a_k(\lambda) \pi \mu_\ell -
\sum\nolimits_{j \in E_k} h_{k, \ell}^j(\lambda) \pi \mu_j
- \pi g_{k, \ell}(\lambda) \\
& =
\pi\left(a_k(\lambda) \mu_\ell -
\sum\nolimits_{j \in E_k} h_{k, \ell}^j(\lambda) \mu_j
- g_{k, \ell}(\lambda)\right)
\end{align*}
위의 (4)에서 $x_i = \lambda_i$를 대입하면 괄호 안의 식은
$\pi^2$에 의해 소멸함을 알 수 있다.
$\text{Ann}_\Lambda(\pi) = \text{Ann}_\Lambda(\pi^2)$라고
가정했으므로 이 식은 $\pi$에 의해서도 소멸한다.
사상 $\bar C \to D/\pi D$는 $x_i \mapsto x_i$로 결정된다
(잘 정의됨은 명백하다).
따라서 (b), (c), (d)를 증명하면 증명이 끝난다.

\medskip\noindent
(4)를 사용하면 다음 핵심 등식을 얻는다.
\begin{align*}
\pi p_{k, \ell} & =
\pi a_k z_\ell - \sum\nolimits_{j \in E_k} \pi h_{k, \ell}^jz_j
- \pi^2 g_{k, \ell} \\
& =
- a_k (f_\ell - \pi z_\ell) + a_k f_\ell +
\sum\nolimits_{j \in E_k} h_{k, \ell}^j (f_j - \pi z_j) -
\sum\nolimits_{j \in E_k} h_{k, \ell}^j f_j - \pi^2 g_{k, \ell} \\
& =
-a_k(f_\ell - \pi z_\ell) +
\sum\nolimits_{j \in E_k} h_{k, \ell}^j(f_j - \pi z_j)
\end{align*}
마지막 식은 $f_j - \pi z_j$들로 생성되는 아이디얼에 속한다.
특히 $D[1/\pi]$는
$R[1/\pi][x_1, \ldots, x_n, z_1, \ldots, z_m]/(f_j - \pi z_j)$와
동형이며, 이는 다시 $R[1/\pi][x_1, \ldots, x_n]$와 동형이므로
$R$ 위에서 매끄럽다. 이로써 (b)가 증명된다.

\medskip\noindent
고정한 $k \in \{1, \ldots, r\}$에 대해 환
$$
D_k = R[x_1, \ldots, x_n, z_1, \ldots, z_m]/
(f_j - \pi z_j, j \in E_k, p_{k, \ell})
$$
을 생각하자. $\ell \in E_k$이면 $p_{k, \ell}$가 영이므로,
방정식의 수는 $m = |E_k| + (m - |E_k|)$이다.
또한
\begin{align*}
(D_k/\pi D_k)_{a_k}
& =
R/\pi R[x_1, \ldots, x_n, 1/a_k, z_1, \ldots, z_m]/
(f_j, j \in E_k, p_{k, \ell}) \\
& =
(\bar C/\pi \bar C)_{a_k}[z_1, \ldots, z_m]/
(a_kz_\ell - \sum\nolimits_{j \in E_k} h_{k, \ell}^j z_j) \\
& \cong
(\bar C/\pi \bar C)_{a_k}[z_j, j \in E_k]
\end{align*}
임에 유의하자.
특히 $(D_k/\pi D_k)_{a_k}$는
$(\bar C/\pi \bar C)_{a_k}$ 위에서 매끄럽다.
$a_k$의 선택에 의해 $(\bar C/\pi \bar C)_{a_k}$는
$R/\pi R$ 위에서 상대차원 $n - |E_k|$로 매끄럽다. (2)를 보라.
따라서 $\pi$를 포함하고 $\Spec(\bar C_{a_k})$ 위에 놓이는
소아이디얼 $\mathfrak q_k \subset D_k$에 대해,
$R \to D_k$의 섬유환은 $\mathfrak q_k$에서 차원이 $n$이고 매끄럽다.
위에서 센 방정식의 수에 의해 $R \to D_k$는
$\mathfrak q_k$에서 신토믹이다.
Algebra, Lemma
\ref{algebra-lemma-localize-relative-complete-intersection}을 보라.
따라서 Algebra, Lemma \ref{algebra-lemma-flat-fibre-smooth}에 의해
$R \to D_k$는 $\mathfrak q_k$에서 매끄럽다.

\medskip\noindent
증명을 마치기 위해, $\pi$를 포함하고
$R/\pi^2 R \to \bar C$가 매끄러운 소아이디얼 위에 놓이는
소아이디얼 $\mathfrak q \subset D$를 택하자.
(1)에 의해 어떤 $k$에 대해 $a_k \not \in \mathfrak q$이다.
전사 사상 $D_k \to D$가 $\mathfrak q$에서의 국소환들 사이의
동형을 유도함을 보이겠다.
앞 문단에 의해 환 준동형 $\bar C/\pi \bar C \to D_k/\pi D_k$와
$R \to D_k$가 대응하는 소아이디얼 $\mathfrak q_k$에서
매끄러움을 알고 있으므로, 이것으로 (c)와 (d)가 증명되어
전체 증명이 끝난다.

\medskip\noindent
먼저 임의의 $\ell$에 대해 위에서 증명한 등식
$\pi p_{k, \ell} = -a_k(f_\ell - \pi z_\ell) +
\sum_{j \in E_k} h_{k, \ell}^j (f_j - \pi z_j)$은
$f_\ell - \pi z_\ell$가 $(D_k)_{a_k}$에서, 따라서 특히
$(D_k)_{\mathfrak q_k}$에서 영으로 보내짐을 보여 준다.
관계식 (4)에 의해 $I/I^2$에서
$a_k f_\ell = \sum_{j \in E_k} h_{k, \ell}^j f_j$이다.
$(\bar I_k/\bar I_k^2)_{a_k}$가 $j \in E_k$인 $f_j$들을
기저로 하는 자유 가군이므로, 모든 $k, k', \ell$ 및 $j \in E_k$에 대해
$$
a_{k'} h_{k, \ell}^j -
\sum\nolimits_{j' \in E_{k'}} h_{k', \ell}^{j'} h_{k, j'}^j
$$
는 $\bar C_{a_k}$에서 영이다.
따라서
$$
a_k^N\left(
a_{k'} h_{k, \ell}^j -
\sum\nolimits_{j' \in E_{k'}} h_{k', \ell}^{j'} h_{k, j'}^j
\right)
$$
가 $I_k + \pi^2R[x_1, \ldots, x_n]$에 속하게 하는
충분히 큰 정수 $N$을 찾을 수 있다.
$\pi$를 법으로 계산하면
\begin{align*}
&
a_kp_{k', \ell} - a_{k'}p_{k, \ell} + \sum h_{k', \ell}^{j'} p_{k, j'}
\\
&
=
- a_k \sum h_{k', \ell}^{j'} z_{j'}
+ a_{k'} \sum h_{k, \ell}^j z_j
+ \sum h_{k', \ell}^{j'} a_k z_{j'}
- \sum \sum h_{k', \ell}^{j'} h_{k, j'}^j z_j \\
&
=
\sum \left(
a_{k'} h_{k, \ell}^j
- \sum h_{k', \ell}^{j'} h_{k, j'}^j
\right) z_j
\end{align*}
이다. 여기서는 아인슈타인 합 규약을 사용하였다.
위의 결과들을 결합하면 $a_k^{N + 1} p_{k', \ell}$은
$R[x_1, \ldots, x_n, z_1, \ldots, z_m]$에서
$I_k$와 $\pi$로 생성되는 아이디얼에 포함된다.
따라서 $p_{k', \ell}$는 $\pi (D_k)_{a_k}$ 안으로 보내진다.
한편 등식
$$
\pi p_{k', \ell} =
-a_{k'} (f_\ell - \pi z_\ell) +
\sum\nolimits_{j' \in E_{k'}} h_{k', \ell}^{j'}(f_{j'} - \pi z_{j'})
$$
은 $\pi p_{k', \ell}$가 $(D_k)_{a_k}$에서 영임을 보여 준다.
$\text{Ann}_R(\pi) = \text{Ann}_R(\pi^2)$라고 가정하였고
$(D_k)_{\mathfrak q_k}$가 매끄러우므로 $R$ 위에서 평탄하다.
따라서
$\text{Ann}_{(D_k)_{\mathfrak q_k}}(\pi) =
\text{Ann}_{(D_k)_{\mathfrak q_k}}(\pi^2)$이다.
결국 $p_{k', \ell}$도 영으로 보내지므로
$D_{\mathfrak q} = (D_k)_{\mathfrak q_k}$이고, 증명이 끝난다.
\end{proof}






\section{특이점 해소 보조정리}
\label{smoothing-section-desingularization-lemma}

\noindent
다음은 또 하나의 지독할 만큼 영리한 보조정리이다.

\begin{lemma}
\label{smoothing-lemma-desingularize}
$R$를 뇌터 환, $\Lambda$를 $R$-대수라 하자.
$\pi \in R$라 하고
$\text{Ann}_\Lambda(\pi) = \text{Ann}_\Lambda(\pi^2)$라고 가정하자.
$A$가 유한 표시를 갖는 $R$-대수 사상 $A \to \Lambda$를 생각하자.
다음을 가정한다.
\begin{enumerate}
\item $\pi$의 상은 $R$ 위의 $A$에서 엄밀 표준이다.
\item $\Lambda/\pi^4 \Lambda$로 가는 사상과 양립하는 절단 사상
$\rho : A/\pi^4 A \to R/\pi^4 R$가 존재한다.
\end{enumerate}
그러면 $B$가 유한 표시를 갖고
$\mathfrak a B \subset H_{B/R}$가 되도록
$R$-대수 사상들 $A \to B \to \Lambda$를 찾을 수 있다.
여기서
$\mathfrak a = \text{Ann}_R(\text{Ann}_R(\pi^2)/\text{Ann}_R(\pi))$이다.
\end{lemma}

\begin{proof}
표시
$$
A = R[x_1, \ldots, x_n]/(f_1, \ldots, f_m)
$$
와 $0 \leq c \leq \min(n, m)$를 택하여,
$\pi$에 대해 (\ref{smoothing-equation-strictly-standard-one})이 성립하고
$j = 1, \ldots, m - c$에 대해
\begin{equation}
\label{smoothing-equation-star}
\pi f_{c + j} \in (f_1, \ldots, f_c) + (f_1, \ldots, f_m)^2
\end{equation}
이 성립하게 하자.
$\rho$가 $x_i$를 $r_i \in R$의 동치류로 보낸다고 하자.
그러면 $x_i$를 $x_i - r_i$로 바꿀 수 있다.
따라서 $R/\pi^4 R$에서 $\rho(x_i) = 0$이라고 가정해도 된다.
이는 $f_j(0) \in \pi^4R$이고,
$A \to \Lambda$가 어떤 $\lambda_i \in \Lambda$에 대해
$x_i$를 $\pi^4\lambda_i$로 보냄을 뜻한다.
$$
f_j = f_j(0) + \sum\nolimits_{i = 1, \ldots, n} r_{ji} x_i + \text{h.o.t.}
$$
로 쓰자. 그러면 $\partial f_j/\partial x_i$의 상수항은 $r_{ji}$이다.
$\pi$에 대한 (\ref{smoothing-equation-strictly-standard-one})에 $\rho$를
적용하면, 어떤 $r_I \in R$에 대해
$$
\pi = \sum\nolimits_{I \subset \{1, \ldots, n\},\ |I| = c}
r_I \det(r_{ji})_{j = 1, \ldots, c,\ i \in I} \bmod \pi^4R
$$
임을 알 수 있다. 따라서 어떤 $u \in 1 + \pi^3R$에 대해
$$
u\pi = \sum\nolimits_{I \subset \{1, \ldots, n\},\ |I| = c}
r_I \det(r_{ji})_{j = 1, \ldots, c,\ i \in I}
$$
이다. Algebra, Lemma \ref{algebra-lemma-matrix-left-inverse}에 의해,
다음 등식을 만족하는 $n \times c$ 행렬 $(s_{ik})$가 존재한다.
$$
u\pi \delta_{jk} = \sum\nolimits_{i = 1, \ldots, n} r_{ji}s_{ik}\quad
\text{for all } j, k = 1, \ldots, c
$$
(크로네커 델타).
보조변수 $v_1, \ldots, v_c, w_1, \ldots, w_n$을 도입하고
$$
h_i = x_i - \pi^2 \sum\nolimits_{j = 1, \ldots c} s_{ij} v_j - \pi^3 w_i
$$
로 놓자. 이하에서는
$$
R[x_1, \ldots, x_n, v_1, \ldots, v_c, w_1, \ldots, w_n]/
(h_1, \ldots, h_n) = R[v_1, \ldots, v_c, w_1, \ldots, w_n]
$$
이라는 사실을 별도 언급 없이 사용한다.\hfil\break
$R[x_1, \ldots, x_n, v_1, \ldots, v_c, w_1, \ldots, w_n]/
(h_1, \ldots, h_n)$에서
\begin{align*}
f_j & = f_j(x_1 - h_1, \ldots, x_n - h_n) \\
& =
\pi^2 \sum\nolimits_{k = 1}^c
\left(\sum\nolimits_{i = 1}^n r_{ji} s_{ik}\right) v_k
+
\pi^3 \sum\nolimits_{i = 1}^n r_{ji}w_i \bmod \pi^4 \\
& =
\pi^3 v_j + \pi^3 \sum\nolimits_{i = 1}^n r_{ji}w_i \bmod \pi^4
\end{align*}
이다($1 \leq j \leq c$).
따라서 $g_j = v_j + \sum r_{ji}w_i \bmod \pi$이고
$R$-대수\hfil\break
$R[x_1, \ldots, x_n, v_1, \ldots, v_c, w_1, \ldots, w_n]/
(h_1, \ldots, h_n)$에서 $f_j = \pi^3 g_j$가 성립하게 하는 원소들
$g_j \in R[v_1, \ldots, v_c, w_1, \ldots, w_n]$을 택할 수 있다.
$$
B = R[x_1, \ldots, x_n, v_1, \ldots, v_c, w_1, \ldots, w_n]/
(f_1, \ldots, f_m, h_1, \ldots, h_n, g_1, \ldots, g_c).
$$
로 놓자. 사상 $A \to B$는 명백하다.
$x_i \to \pi^4\lambda_i$, $v_i \mapsto 0$,
$w_i \mapsto \pi \lambda_i$로 보내어 $B \to \Lambda$를 정의하자.
그러면 원소 $f_j$와 $h_i$가 $\Lambda$에서 영으로 보내짐은 명백하다.
또한 $g_i$는 $\pi\Lambda$의 원소 $t$로 보내지며
$\pi^3t = 0$임도 명백하다
($h$들로 생성되는 아이디얼을 법으로 $f_i = \pi^3 g_i$이기 때문이다).
따라서 $\text{Ann}_\Lambda(\pi) = \text{Ann}_\Lambda(\pi^2)$라는
가정에 의해 $t = 0$이다.
이제 매끄러움에 관한 명제를 증명하면 증명이 끝난다.

\medskip\noindent
방정식 $g_i = 0$은 $f_i = 0$에서 따라오므로
$B_\pi \cong A_\pi[v_1, \ldots, v_c]$임에 유의하자.
$\pi$가 $R$ 위의 $A$에서 엄밀 표준이라는 가정에 의해
$A_\pi$가 $R$ 위에서 매끄러우므로, $B_\pi$도 $R$ 위에서 매끄럽다.
보조정리 \ref{smoothing-lemma-elkik}를 보라.

\medskip\noindent
$B' = R[v_1, \ldots, v_c, w_1, \ldots, w_n]/(g_1, \ldots, g_c)$로 놓자.
$g_i = v_i + \sum r_{ji}w_i \bmod \pi$이므로
$B'/\pi B' = R/\pi R[w_1, \ldots, w_n]$이다.
따라서 Algebra, Lemmas
\ref{algebra-lemma-localize-relative-complete-intersection} 및
\ref{algebra-lemma-flat-fibre-smooth}에 의해,
$R \to B'$는 $V(\pi)$의 모든 점에서 상대차원이 $n$인 매끄러운
사상이다(첫 번째 보조정리는 그 소아이디얼들에서 이 사상이
신토믹, 특히 평탄함을 보여 주고, 그러면 두 번째 보조정리가
매끄러움을 보여 준다).

\medskip\noindent
$\pi \in \mathfrak q$이고 어떤 $r \in \mathfrak a$에 대해
$r \not \in \mathfrak q$인 소아이디얼 $\mathfrak q \subset B$를
생각하자. $\mathfrak q' = B' \cap \mathfrak q$로 표기한다.
전사 사상 $B' \to B$가 국소환의 동형
$(B')_{\mathfrak q'} \to B_\mathfrak q$를 유도한다고 주장한다.
이것으로 보조정리의 증명이 끝난다.
$B_\mathfrak q$는 $(B')_{\mathfrak q'}$를
$j = 1, \ldots, m - c$인 $f_{c + j}$들로 생성되는 아이디얼로
나눈 몫임에 유의하자.
두 가지를 관찰하자. 첫째, $(B')_{\mathfrak q'}$에서
$f_{c + j}$의 상은 $\pi^2$로 나누어떨어진다.
둘째, (\ref{smoothing-equation-star})에 의해 $(B')_{\mathfrak q'}$에서
$\pi f_{c + j}$의 상은
$\sum b_{j_1 j_2} f_{c + j_1}f_{c + j_2}$로 쓸 수 있다.
따라서 각 $\pi f_{c + j}$의 상은
$\pi^2 f_{c + j'}$들로 생성되는 아이디얼에 포함된다.
$(B')_{\mathfrak q'}$는 뇌터 국소환이므로
$\pi f_{c + j} = 0$이다.
Algebra, Lemma
\ref{algebra-lemma-intersect-powers-ideal-module-zero}를 보라.
$R \to (B')_{\mathfrak q'}$가 평탄하므로
$$
\left(\text{Ann}_R(\pi^2)/\text{Ann}_R(\pi)\right)
\otimes_R (B')_{\mathfrak q'}
=
\text{Ann}_{(B')_{\mathfrak q'}}(\pi^2)/\text{Ann}_{(B')_{\mathfrak q'}}(\pi)
$$
이다. $r \in \mathfrak a$가 $(B')_{\mathfrak q'}$에서 가역이므로
이 가군은 영이다.
따라서 원하는 대로 $f_{c + j}$의 상은 $(B')_{\mathfrak q'}$에서
영이다.
\end{proof}

\begin{lemma}
\label{smoothing-lemma-desingularize-strictly-standard}
$R$를 뇌터 환, $\Lambda$를 $R$-대수라 하자.
$\pi \in R$라 하고
$\text{Ann}_R(\pi) = \text{Ann}_R(\pi^2)$ 및
$\text{Ann}_\Lambda(\pi) = \text{Ann}_\Lambda(\pi^2)$라고 가정하자.
$A$와 $D$가 유한 표시를 갖는 $R$-대수 사상들
$A \to \Lambda$ 및 $D \to \Lambda$를 생각하자.
다음을 가정한다.
\begin{enumerate}
\item $\pi$는 $R$ 위의 $A$에서 엄밀 표준이다.
\item $\Lambda/\pi^4 \Lambda$로 가는 사상들과 양립하는
$R$-대수 사상 $A/\pi^4 A \to D/\pi^4 D$가 존재한다.
\end{enumerate}
그러면 $B$가 유한 표시를 갖는 $R$-대수 사상 $B \to \Lambda$와,
$\Lambda$로 가는 사상들과 양립하는
$R$-대수 사상들 $A \to B$와 $D \to B$가 존재하며,
$H_{D/R}B \subset H_{B/D}$ 및 $H_{D/R}B \subset H_{B/R}$가 된다.
\end{lemma}

\begin{proof}
$D$에서의 $\pi$의 상과 사상들
$$
D \longrightarrow A \otimes_R D \longrightarrow \Lambda
$$
에 보조정리 \ref{smoothing-lemma-desingularize}를 적용한다.
보조정리 \ref{smoothing-lemma-strictly-standard-base-change}에 의해
$\pi$는 $D$ 위의 $A \otimes_R D$에서 엄밀 표준이다.
(2)의 사상에서 유도되는 사상을 절단 사상
$\rho : (A \otimes_R D)/\pi^4 (A \otimes_R D) \to D/\pi^4 D$로 택한다.
따라서 보조정리 \ref{smoothing-lemma-desingularize}를 적용할 수 있으며,
$B$가 유한 표시를 갖고
$\mathfrak a B \subset H_{B/D}$가 되는 분해
$A \otimes_R D \to B \to \Lambda$를 얻는다. 여기서
$$
\mathfrak a = \text{Ann}_D(\text{Ann}_D(\pi^2)/\text{Ann}_D(\pi)).
$$
이다.
$D_\mathfrak q$가 $R$ 위에서 평탄한 $D$의 임의의 소아이디얼
$\mathfrak q$에 대해\hfil\break
$\text{Ann}_{D_\mathfrak q}(\pi^2)/\text{Ann}_{D_\mathfrak q}(\pi) = 0$
이다. 원소의 소멸자는 평탄한 밑변환과 가환하고
$\text{Ann}_R(\pi) = \text{Ann}_R(\pi^2)$라고 가정했기 때문이다.
$D$가 뇌터 환이므로
$\text{Ann}_D(\pi^2)/\text{Ann}_D(\pi)$는 유한 $D$-가군이며,
따라서 그 소멸자를 취하는 것은 국소화와 가환한다.
그러므로 $\mathfrak a \not \subset \mathfrak q$임을 알 수 있다.
따라서 $D \to B$는 $\mathfrak q$ 위에 놓이는 $B$의 모든
소아이디얼에서 매끄럽다.
$R \to D$가 매끄러운 $D$의 모든 소아이디얼에서
$D_\mathfrak q$는 $R$ 위에서 평탄하므로
$H_{D/R}B \subset H_{B/D}$를 얻는다.
마지막 포함관계 $H_{D/R}B \subset H_{B/R}$는
매끄러운 환 준동형의 합성이 매끄럽다는 사실에서 따른다
(Algebra, Lemma \ref{algebra-lemma-compose-smooth}).
\end{proof}

\begin{lemma}
\label{smoothing-lemma-desingularize-lifting-apply}
$R$를 뇌터 환, $\Lambda$를 $R$-대수라 하자.
$\pi \in R$라 하고
$\text{Ann}_R(\pi) = \text{Ann}_R(\pi^2)$ 및
$\text{Ann}_\Lambda(\pi) = \text{Ann}_\Lambda(\pi^2)$라고 가정하자.
$A$가 유한 표시를 갖는 $R$-대수 사상 $A \to \Lambda$를 생각하고,
$\pi$가 $R$ 위의 $A$에서 엄밀 표준이라고 가정하자.
$\bar C$가 유한 표시를 갖는 분해
$$
A/\pi^8A \to \bar C \to \Lambda/\pi^8\Lambda
$$
를 생각하자. 그러면 $B$가 유한 표시를 갖고
$R_\pi \to B_\pi$가 매끄러우며
$$
H_{\bar C/(R/\pi^8 R)} \cdot \Lambda/\pi^8\Lambda
\subset
\sqrt{H_{B/R} \Lambda} \bmod \pi^8\Lambda.
$$
가 되도록 분해 $A \to B \to \Lambda$를 찾을 수 있다.
\end{lemma}

\begin{proof}
보조정리 \ref{smoothing-lemma-lifting}을 적용하여
$\bar C/\pi^4\bar C \to D/\pi^4 D \to \Lambda/\pi^4\Lambda$라는
분해를 갖는 $R \to D \to \Lambda$를 얻자.
이때 $R \to D$는 $\pi$를 포함하지 않는 모든 소아이디얼과,
$R/\pi^8 R \to \bar C$가 매끄러운
$\bar C/\pi^4\bar C$의 소아이디얼 위에 놓이는 모든 소아이디얼에서
매끄럽다.
보조정리 \ref{smoothing-lemma-desingularize-strictly-standard}에 의해,
유한 표시 $R$-대수 $B$와
$H_{D/R}B \subset H_{B/R}$가 되는 분해들
$A \to B \to \Lambda$ 및 $D \to B \to \Lambda$를 찾을 수 있다.
이것이 보조정리가 제시한 문제의 해임을 확인하는 과정은 생략한다.
\end{proof}












\section{준비 단계: 기저체로의 환원}
\label{smoothing-section-reduction}

\noindent
이 절에서는 앞 절들의 보조정리를 적용하여,
기저환이 체인 경우에 주결과를 증명하면 충분함을 보인다.
보조정리 \ref{smoothing-lemma-reduce-to-field}를 보라.

\begin{situation}
\label{smoothing-situation-global}
여기서 $R \to \Lambda$는 뇌터 환들 사이의 정칙 환 준동형이다.
\end{situation}

\noindent
$R \to \Lambda$가 상황 \ref{smoothing-situation-global}과 같다고 하자.
$\Lambda$가 매끄러운 $R$-대수들의 여과 여극한이면
{\it $R \to \Lambda$에 대해 PT가 성립한다}고 말한다.

\begin{lemma}
\label{smoothing-lemma-product}
$R_i \to \Lambda_i$, $i = 1, 2$가 상황 \ref{smoothing-situation-global}과 같다고
하자. $i = 1, 2$에 대해 $R_i \to \Lambda_i$에서 PT가 성립하면
$R_1 \times R_2 \to \Lambda_1 \times \Lambda_2$에서도 PT가 성립한다.
\end{lemma}

\begin{proof}
생략한다. 힌트: 여과 여극한들의 곱은 여과 여극한이다.
\end{proof}

\begin{lemma}
\label{smoothing-lemma-delocalize-base}
$A$가 $R$ 위에서 유한 표시를 갖는 환 준동형들
$R \to A \to \Lambda$를 생각하자.
$S \subset R$를 곱셈적 집합이라 하자.
$B'$가 $S^{-1}R$ 위에서 매끄러운 분해
$S^{-1}A \to B' \to S^{-1}\Lambda$를 생각하자.
그러면 어떤 $s \in S$가 $R$ 위의 $B$에서 초등 표준 원소가 되도록
(정의 \ref{smoothing-definition-strictly-standard})
분해 $A \to B \to \Lambda$를 찾을 수 있다.
\end{lemma}

\begin{proof}
먼저 $S^{-1}R \to B'$에 보조정리
\ref{smoothing-lemma-smooth-standard-smooth}를 적용한다.
따라서 $B'$가 $S^{-1}R$ 위에서 표준 매끄럽다고 가정해도 된다.
$A = R[x_1, \ldots, x_n]/(g_1, \ldots, g_t)$로 쓰고,
$\Lambda$에서 $x_i \mapsto \lambda_i$라고 하자.
어떤 $c \geq n$에 대해
$B' = S^{-1}R[x_1, \ldots, x_{n + m}]/(f_1, \ldots, f_c)$로
쓸 수 있다. 여기서
$\det(\partial f_j/\partial x_i)_{i, j = 1, \ldots, c}$는 $B'$에서
가역이고, $A \to B'$는 $x_i \mapsto x_i$로 주어진다.
보조정리 \ref{smoothing-lemma-standard-smooth-include-generators}를 보라.
$i > n$인 $x_i$에 $S$의 한 원소를 곱하고 방정식 $f_j$를 그에 맞게
수정하면, $i > n$에 대해 어떤 $\lambda_i \in \Lambda$가 존재하여
$B' \to S^{-1}\Lambda$가 $x_i$를 $\lambda_i/1$로 보낸다고
가정할 수 있다.
어떤 $a_j \in S^{-1}R[x_1, \ldots, x_{n + m}]$에 대해 관계식
$$
1 =
a_0 \det(\partial f_j/\partial x_i)_{i, j = 1, \ldots, c}
+
\sum\nolimits_{j = 1, \ldots, c} a_jf_j
$$
을 택하자. $S$의 각 원소는 $B'$에서 가역이므로
(분모를 없애어) $f_j, a_j \in R[x_1, \ldots, x_{n + m}]$이고
어떤 $s_0 \in S$에 대해
$$
s_0 = a_0 \det(\partial f_j/\partial x_i)_{i, j = 1, \ldots, c}
+
\sum\nolimits_{j = 1, \ldots, c} a_jf_j
$$
라고 가정할 수 있다.
$g_j$가
$S^{-1}R[x_1, \ldots, x_{n + m}]/(f_1, \ldots, x_c)$에서
영으로 보내지므로,
$R[x_1, \ldots, x_{n + m}]/(f_1, \ldots, f_c)$에서
$s_j g_j = 0$이 되게 하는 원소 $s_j \in S$를 찾을 수 있다.
$f_j$가 $S^{-1}\Lambda$에서 영으로 보내지므로,
$\Lambda$에서
$s'_j f_j(\lambda_1, \ldots, \lambda_{n + m}) = 0$이 되게 하는
$s'_j \in S$를 찾을 수 있다. 환
$$
B = R[x_1, \ldots, x_{n + m}]/
(s'_1f_1, \ldots, s'_cf_c, g_1, \ldots, g_t)
$$
과, $B \to \Lambda$가 $x_i \mapsto \lambda_i$로 주어지는 분해
$A \to B \to \Lambda$를 생각하자.
$s = s_0s_1 \ldots s_ts'_1 \ldots s'_c$가
$R$ 위의 $B$에서 초등 표준이라고 주장한다.
이것으로 증명이 끝난다.
실제로 $s_j g_j \in (f_1, \ldots, f_c)$이므로
$sg_j \in (s'_1f_1, \ldots, s'_cf_c)$이다. 끝으로
$$
a_0\det(\partial s'_jf_j/\partial x_i)_{i, j = 1, \ldots, c}
+
\sum\nolimits_{j = 1, \ldots, c}
(s'_1 \ldots \hat{s'_j} \ldots s'_c) a_j s'_jf_j
=
s_0s'_1\ldots s'_c
$$
이며, 원하는 대로 오른쪽 항은 $s$를 나눈다.
\end{proof}

\begin{lemma}
\label{smoothing-lemma-reduce-to-field}
\begin{slogan}
Popescu 근사를 증명하는 문제는 체 위의 대수인 경우로 환원된다
\end{slogan}
$R$가 체인 모든 상황 \ref{smoothing-situation-global}에서 PT가 성립하면,
일반적인 경우에도 PT가 성립한다.
\end{lemma}

\begin{proof}
$R$가 체인 모든 상황 \ref{smoothing-situation-global}에서 PT가 성립한다고
가정하자. 상황 \ref{smoothing-situation-global}과 같은 임의의
$R \to \Lambda$를 생각하자.
$R/I \to \Lambda/I\Lambda$도 뇌터 환들 사이의 정칙 환 준동형임에
유의하자. More on Algebra, Lemma
\ref{more-algebra-lemma-regular-base-change}를 보라.
아이디얼들의 집합
$$
\mathcal{I} = \{I \subset R \mid R/I \to \Lambda/I\Lambda
\text{ does not have PT}\}
$$
를 생각하자. $\mathcal{I}$가 공집합임을 보여야 한다.
이 집합이 공집합이 아니면 $R$가 뇌터 환이므로 극대 원소를 갖는다.
$R$를 $R/I$로, $\Lambda$를 $\Lambda/I$로 바꾸면,
$R$의 임의의 영이 아닌 아이디얼에 대해
$R/I \to \Lambda/I\Lambda$에서 PT가 성립하는 상황을 얻는다.
특히 명제 \ref{smoothing-proposition-lift}를 적용하면 $R$가 축약환임을 알 수 있다.

\medskip\noindent
$A \to \Lambda$를, $A$가 유한 표시를 갖는
$R$-대수 준동형이라고 하자.
$B$가 $R$ 위에서 매끄러운 분해 $A \to B \to \Lambda$를 찾아야 한다.
Algebra, Lemma \ref{algebra-lemma-when-colimit}를 보라.

\medskip\noindent
$S \subset R$를 비영인자들의 집합이라 하고,
$R$의 전분수환 $Q = S^{-1}R$을 생각하자.
Algebra, Lemmas
\ref{algebra-lemma-total-ring-fractions-no-embedded-points} 및
\ref{algebra-lemma-Noetherian-irreducible-components}에 의해,
$Q = K_1 \times \ldots \times K_n$은 체들의 곱이다.
보조정리 \ref{smoothing-lemma-product}와 가정에 의해
환 준동형 $S^{-1}R \to S^{-1}\Lambda$에서 PT가 성립한다.
따라서 $B'$가 $S^{-1}R$ 위에서 매끄러운 분해
$S^{-1}A \to B' \to S^{-1}\Lambda$를 찾을 수 있다.

\medskip\noindent
보조정리 \ref{smoothing-lemma-delocalize-base}를 적용하면,
어떤 $\pi \in S$가 $R$ 위의 $B$에서 초등 표준이 되는 분해
$A \to B \to \Lambda$를 얻는다.
$A$를 $B$로 바꾸어, $\pi$가 $A$에서 초등 표준이고 따라서
엄밀 표준이라고 가정할 수 있다.
$R/\pi^8R \to \Lambda/\pi^8\Lambda$에서 PT가 성립함을 알고 있다.
따라서 $R/\pi^8 R \to \bar C$가 매끄러운 분해
$R/\pi^8 R \to A/\pi^8A \to \bar C \to \Lambda/\pi^8\Lambda$를
찾을 수 있다.
보조정리 \ref{smoothing-lemma-lifting}에 의해 $D$가 $R$ 위에서 매끄럽고
분해
$R/\pi^4 R \to A/\pi^4A \to D/\pi^4D \to \Lambda/\pi^4\Lambda$를
갖는 $R$-대수 사상 $D \to \Lambda$를 찾을 수 있다.
보조정리 \ref{smoothing-lemma-desingularize-strictly-standard}에 의해
$B$가 $R$ 위에서 매끄러운 $A \to B \to \Lambda$를 찾을 수 있으며,
이로써 증명이 끝난다.
\end{proof}










\section{국소적 기법}
\label{smoothing-section-local}


\begin{situation}
\label{smoothing-situation-local}
뇌터 환 $R$, $R$-대수 사상 $A \to \Lambda$, 그리고 소 아이디얼
$\mathfrak q \subset \Lambda$가 주어져 있다고 하자. $A$가 $R$ 위에서
유한 표시를 갖는다고 가정한다. 이 상황에서는
$\mathfrak h_A = \sqrt{H_{A/R} \Lambda}$로 쓴다.
\end{situation}

\noindent
상황 \ref{smoothing-situation-local}과 같은
$R \to A \to \Lambda \supset \mathfrak q$를 생각하자. $B$가 유한 표시를
가지며 $\mathfrak h_A \subset \mathfrak h_B \not \subset \mathfrak q$를
만족하는 분해 $A \to B \to \Lambda$가 존재하면
{\it $R \to A \to \Lambda \supset \mathfrak q$는 해소 가능하다}고 한다.
이 경우 분해 $A \to B \to \Lambda$를
{\it $R \to A \to \Lambda \supset \mathfrak q$의 해소}라고 부른다.

\begin{lemma}
\label{smoothing-lemma-lift-solution}
상황 \ref{smoothing-situation-local}과 같은
$R \to A \to \Lambda \supset \mathfrak q$를 생각하자. $r \geq 1$이라 하고,
$\pi_1, \ldots, \pi_r \in R$의 상이 $\mathfrak q$의 원소라고 하자. 다음을 가정한다.
\begin{enumerate}
\item $i = 1, \ldots, r$에 대해
$$
\text{Ann}_{R/(\pi_1^8, \ldots, \pi_{i - 1}^8)R}(\pi_i)
=
\text{Ann}_{R/(\pi_1^8, \ldots, \pi_{i - 1}^8)R}(\pi_i^2)
$$
이고
$$
\text{Ann}_{\Lambda/(\pi_1^8, \ldots, \pi_{i - 1}^8)\Lambda}(\pi_i)
=
\text{Ann}_{\Lambda/(\pi_1^8, \ldots, \pi_{i - 1}^8)\Lambda}(\pi_i^2)
$$
\item $i = 1, \ldots, r$에 대해 원소 $\pi_i$의 상은 $R$ 위의 $A$에서
엄밀 표준 원소이다.
\end{enumerate}
그러면
$$
R/(\pi_1^8, \ldots, \pi_r^8)R \to A/(\pi_1^8, \ldots, \pi_r^8)A
\to \Lambda/(\pi_1^8, \ldots, \pi_r^8)\Lambda \supset
\mathfrak q/(\pi_1^8, \ldots, \pi_r^8)\Lambda
$$
가 해소 가능하면 $R \to A \to \Lambda \supset \mathfrak q$도 해소 가능하다.
\end{lemma}

\begin{proof}
$r$에 대한 귀납법으로 증명한다.

\medskip\noindent
$r = 1$인 경우를 보자. 이때 가정에 의해 $\pi_1^8$을 법으로 한 상황을
해소하는 분해 $A/\pi_1^8 \to \bar C \to \Lambda/\pi_1^8$이 존재한다.
조건 (1)과 (2)는 보조정리 \ref{smoothing-lemma-desingularize-lifting-apply}를
적용하는 데 필요한 가정이다. 따라서 해소 $\bar C$를
``끌어올려'' $R \to A \to \Lambda \supset \mathfrak q$의 해소를 얻을 수 있다.

\medskip\noindent
$r > 1$인 경우를 보자. 이때는 다음 상황에 $r - 1$에 대한 귀납 가정을 적용한다.
$R/\pi_1^8 \to A/\pi_1^8 \to \Lambda/\pi_1^8
\supset \mathfrak q/\pi_1^8\Lambda$.
조건 (2)는 보조정리 \ref{smoothing-lemma-strictly-standard-base-change}에 의해
보존됨에 유의하자.
\end{proof}

\begin{lemma}
\label{smoothing-lemma-delocalize-weak}
\begin{reference}
\cite[Lemma 12.2]{swan} 또는
\cite[Lemma 2]{popescu-GND}
\end{reference}
상황 \ref{smoothing-situation-local}과 같은
$R \to A \to \Lambda \supset \mathfrak q$를 생각하자.
$\mathfrak p = R \cap \mathfrak q$라 하자. $\mathfrak q$가
$\mathfrak h_A$ 위에서 극소이고
$R_\mathfrak p \to A_\mathfrak p \to \Lambda_\mathfrak q
\supset \mathfrak q\Lambda_\mathfrak q$가 해소 가능하다고 가정하자.
그러면 $C$가 유한 표시를 가지며
$H_{C/R} \Lambda \not \subset \mathfrak q$를 만족하는 분해
$A \to C \to \Lambda$가 존재한다.
\end{lemma}

\begin{proof}
다음 해소를 생각하자.
$A_\mathfrak p \to C \to \Lambda_\mathfrak q$,
$R_\mathfrak p \to A_\mathfrak p \to \Lambda_\mathfrak q
\supset \mathfrak q\Lambda_\mathfrak q$의 해소이다.
$\mathfrak q$가 $\mathfrak h_A$ 위에서 극소라는 가정에 의해 이는
$H_{C/R_\mathfrak p} \Lambda_\mathfrak q = \Lambda_\mathfrak q$를 뜻한다.
보조정리 \ref{smoothing-lemma-final-solve}에 의해 $C$가 $R_\mathfrak p$ 위에서
매끄럽다고 가정해도 된다. 보조정리 \ref{smoothing-lemma-smooth-standard-smooth}에 의해
$C$가 $R_\mathfrak p$ 위에서 표준 매끄럽다고 가정해도 된다.
$A = R[x_1, \ldots, x_n]/(g_1, \ldots, g_t)$로 쓰고,
$A \to \Lambda$가 $x_i \mapsto \lambda_i$로 주어진다고 하자.
어떤 $c \geq n$에 대해
$C = R_\mathfrak p[x_1, \ldots, x_{n + m}]/(f_1, \ldots, f_c)$로 쓰자.
여기서 $A \to C$는 $x_i$를 $x_i$로 보내고,
$\det(\partial f_j/\partial x_i)_{i, j = 1, \ldots, c}$
는 $C$에서 가역이다. 보조정리
\ref{smoothing-lemma-standard-smooth-include-generators}를 보라.
분모를 제거하여 $f_1, \ldots, f_c$가
$R[x_1, \ldots, x_{n + m}]$의 원소라고 가정해도 된다. 물론
$\det(\partial f_j/\partial x_i)_{i, j = 1, \ldots, c}$
는 $R[x_1, \ldots, x_{n + m}]/(f_1, \ldots, f_c)$에서 가역이 아니지만,
어떤 원소 $s_0 \in R$, $s_0 \not \in \mathfrak p$를 가역화하면 가역이 된다.
$R[x_1, \ldots, x_n] \to A \to C$ 아래에서 $g_j$의 상이 영이므로,
$R[x_1, \ldots, x_{n + m}]/(f_1, \ldots, f_c)$에서 $s_jg_j$가 영이 되는
$s_j \in R$, $s_j \not \in \mathfrak p$를 찾을 수 있다.
다항식
$F_j \in R[x_1, \ldots, x_n, X_{n + 1}, \ldots, X_{n + m + 1}]$
이 $X_{n + 1}, \ldots, X_{n + m + 1}$에 관해 동차이고
$f_j = F_j(x_1, \ldots, x_{n + m}, 1)$을 만족하도록 잡자.
$i = 1, \ldots, m + 1$에 대해 $\lambda_{n + i} \in \Lambda$를
$\lambda_{n + m + 1} \not \in \mathfrak q$이고, $\Lambda_\mathfrak q$에서
$x_{n + i}$의 상이 $\lambda_{n + i}/\lambda_{n + m + 1}$이 되도록 잡자. 그러면
\begin{align*}
F_j(\lambda_1, \ldots, \lambda_{n + m + 1})
& =
(\lambda_{n + m + 1})^{\deg(F_j)} F_j(\lambda_1, \ldots, \lambda_n,
\frac{\lambda_{n + 1}}{\lambda_{n + m + 1}}, \ldots,
\frac{\lambda_{n + m}}{\lambda_{n + m + 1}}, 1) \\
& =
(\lambda_{n + m + 1})^{\deg(F_j)} f_j(\lambda_1, \ldots, \lambda_n,
\frac{\lambda_{n + 1}}{\lambda_{n + m + 1}}, \ldots,
\frac{\lambda_{n + m}}{\lambda_{n + m + 1}}) \\
& = 0
\end{align*}
가 $\Lambda_\mathfrak q$에서 성립한다. 따라서 $\Lambda$에서
$\lambda_0 F_j(\lambda_1, \ldots, \lambda_{n + m + 1}) = 0$을 만족하는
$\lambda_0 \in \Lambda$, $\lambda_0 \not \in \mathfrak q$를 찾을 수 있다.
이제 $B$를
$$
R[x_0, \ldots, x_{n + m + 1}]/
(g_1, \ldots, g_t, x_0F_1(x_1, \ldots, x_{n + m + 1}), \ldots,
x_0F_c(x_1, \ldots, x_{n + m + 1}))
$$
로 두고 $x_i$를 $\lambda_i$로 보내어 $\Lambda$로 사상한다.
$b$를 $B$에서 $x_0 x_{n + m + 1} s_0 s_1 \ldots s_t$의 상이라 하자.
그러면 $B_b$는
$$
R_{s_0s_1 \ldots s_t}[x_0, x_1, \ldots, x_{n + m + 1}, 1/x_0x_{n + m + 1}]/
(f_1, \ldots, f_c)
$$
와 동형이고, 이는 구성상 $R$ 위에서 매끄럽다.
$b$의 상은 $\mathfrak q$의 원소가 아니므로 원하는 결론을 얻는다.
\end{proof}

\begin{lemma}
\label{smoothing-lemma-delocalize-height-zero}
상황 \ref{smoothing-situation-local}과 같은
$R \to A \to \Lambda \supset \mathfrak q$를 생각하자.
$\mathfrak p = R \cap \mathfrak q$라 하자. 다음을 가정한다.
\begin{enumerate}
\item $\mathfrak q$는 $\mathfrak h_A$ 위에서 극소이다.
\item $R_\mathfrak p \to A_\mathfrak p \to \Lambda_\mathfrak q
\supset \mathfrak q\Lambda_\mathfrak q$는 해소 가능하다.
\item $\dim(\Lambda_\mathfrak q) = 0$.
\end{enumerate}
그러면 $R \to A \to \Lambda \supset \mathfrak q$는 해소 가능하다.
\end{lemma}

\begin{proof}
(3)에 의해 환 $\Lambda_\mathfrak q$는 아르틴 국소환이므로
$\mathfrak q\Lambda_\mathfrak q$는 멱영이다. 따라서 어떤 $N > 0$에 대해
$(\mathfrak h_A)^N \Lambda_\mathfrak q = 0$이다. 그러므로 $\Lambda$에서
$\lambda (\mathfrak h_A)^N = 0$을 만족하는
$\lambda \in \Lambda$, $\lambda \not \in \mathfrak q$가 존재한다.
$H_{A/R} = (a_1, \ldots, a_r)$라 쓰면 $\Lambda$에서
$\lambda a_i^N = 0$이다. 보조정리 \ref{smoothing-lemma-delocalize-weak}에 의해
$C$가 유한 표시를 가지며 $\mathfrak h_C \not \subset \mathfrak q$를
만족하는 분해 $A \to C \to \Lambda$를 찾을 수 있다.
$C = A[x_1, \ldots, x_n]/(f_1, \ldots, f_m)$로 쓰고
$$
B = A[x_1, \ldots, x_n, y_1, \ldots, y_r, z, t_{ij}]/
(f_j - \sum y_i t_{ij}, zy_i)
$$
로 두자. 여기서 $t_{ij}$는 $rm$개의 변수로 이루어진 집합이다.
$t_{ij}$를 영으로 두어 정의되는 사상 $B \to C[y_i, z]/(y_iz)$가 있음에
유의하자. 사상 $B \to \Lambda$는 합성
$B \to C[y_i, z]/(y_iz) \to \Lambda$이다. 여기서
$C[y_i, z]/(y_iz) \to \Lambda$는 주어진 사상 $C \to \Lambda$이고,
$z$를 $\lambda$로 보내며 $y_i$를 $\Lambda$에서 $a_i^N$의 상으로 보낸다.

\medskip\noindent
$B$가 $R \to A \to \Lambda \supset \mathfrak q$의 해를 준다고 주장한다.
먼저 $B_z$는 $C[y_1, \ldots, y_r, z, z^{-1}]$와 동형이므로 매끄럽다.
한편
$B_{y_\ell} \cong A[x_i, y_i, y_\ell^{-1}, t_{ij}, i \not = \ell]$
는 $A$ 위에서 매끄럽다. 따라서 $z$와 $a_\ell y_\ell$은 모두
$H_{B/R}$의 원소임을 알 수 있다. 여기서는 매끄러운 사상들의 합성이
매끄럽다는 사실을 썼다. 이로써 보조정리가 증명되었다.
\end{proof}




\section{분리적 잉여류체}
\label{smoothing-section-separable}

\noindent
이 절에서는 잉여류체 확대가 분리적인 경우에 국소 문제를 푸는 방법을 설명한다.

\begin{lemma}[Ogoma]
\label{smoothing-lemma-ogoma}
뇌터 환 $A$와 유한 $A$-모듈 $M$을 생각하자.
$S \subset A$를 곱셈적 집합이라 하자. $\pi \in A$이고
$\Ker(\pi : S^{-1}M \to S^{-1}M) =
\Ker(\pi^2 : S^{-1}M \to S^{-1}M)$
이면, 임의의 $n > 0$에 대해
$\Ker(s^n\pi : M \to M) = \Ker((s^n\pi)^2 : M \to M)$.
을 만족하는 $s \in S$가 존재한다.
\end{lemma}

\begin{proof}
$K = \Ker(\pi : M \to M)$,
$K' = \{m \in M \mid \pi^2 m = 0\text{ in }S^{-1}M\}$,
$Q = K'/K$라 하자. 가정에 의해 $S^{-1}Q = 0$임에 유의하자.
$A$가 뇌터 환이므로 $Q$는 유한 $A$-모듈이다.
따라서 $s$가 $Q$를 소멸시키는 $s \in S$를 찾을 수 있다.
이 $s$가 원하는 원소이다.
\end{proof}

\begin{lemma}
\label{smoothing-lemma-find-sequence}
$\Lambda$를 뇌터 환이라 하자. $I \subset \Lambda$를 아이디얼이라 하고,
$I \subset \mathfrak q$인 소 아이디얼을 생각하자.
$n, e$를 양의 정수라 하자. $\mathfrak q^n\Lambda_\mathfrak q \subset I\Lambda_\mathfrak q$이고
$\Lambda_\mathfrak q$가 차원 $d$인 정칙 국소환이라고 가정하자.
그러면 어떤 $n > 0$과 $\pi_1, \ldots, \pi_d \in \Lambda$가 존재하여
\begin{enumerate}
\item $(\pi_1, \ldots, \pi_d)\Lambda_\mathfrak q =
\mathfrak q\Lambda_\mathfrak q$,
\item $\pi_1^n, \ldots, \pi_d^n \in I$이고,
\item $i = 1, \ldots, d$에 대해
$$
\text{Ann}_{\Lambda/(\pi_1^e, \ldots, \pi_{i - 1}^e)\Lambda}(\pi_i) =
\text{Ann}_{\Lambda/(\pi_1^e, \ldots, \pi_{i - 1}^e)\Lambda}(\pi_i^2).
$$
\end{enumerate}
\end{lemma}

\begin{proof}
$S = \Lambda \setminus \mathfrak q$로 두면
$\Lambda_\mathfrak q = S^{-1}\Lambda$이다.
먼저 (1)을 만족하는 $\pi_1, \ldots, \pi_d$를 잡자.
$\Lambda_\mathfrak q$가 정칙이므로 이는 가능하다. 가정에 의해
$\pi_i^n \in I\Lambda_\mathfrak q$이다. 따라서 $s_i\pi_i^n \in I$를 만족하는
$s_1, \ldots, s_d \in S$를 찾을 수 있다.
$\pi_i$를 $s_i\pi_i$로 바꾸면 (2)를 얻는다.
(1)과 (2)는 $S$의 원소를 더 곱해도 보존됨에 유의하자.
어떤 $t \in \{0, \ldots, d\}$에 대해 $i = 1, \ldots, t$일 때
(3)이 성립한다고 하자. Algebra, Lemma
\ref{algebra-lemma-regular-ring-CM}에 의해
$\pi_1, \ldots, \pi_d$는 $S^{-1}\Lambda$에서 정칙열이다.
특히 Algebra, Lemma \ref{algebra-lemma-regular-sequence-powers}에 의해
$\pi_1^e, \ldots, \pi_t^e, \pi_{t + 1}$은
$S^{-1}\Lambda = \Lambda_\mathfrak q$에서 정칙열이다. 따라서
$$
\text{Ann}_{S^{-1}\Lambda/(\pi_1^e, \ldots, \pi_{i - 1}^e)}(\pi_i) =
\text{Ann}_{S^{-1}\Lambda/(\pi_1^e, \ldots, \pi_{i - 1}^e)}(\pi_i^2).
$$
보조정리 \ref{smoothing-lemma-ogoma}에 의해 어떤 $s \in S$에 대해
$\pi_{t + 1}$을 $s\pi_{t + 1}$로 바꾸면 $i = t + 1$일 때의 (3)을 얻는다.
$t$에 대한 귀납법으로 (1), (2), (3)을 만족하는 열을 얻는다.
\end{proof}

\begin{lemma}
\label{smoothing-lemma-resolve-special}
상황 \ref{smoothing-situation-local}과 같은
$k \to A \to \Lambda \supset \mathfrak q$를 생각하자. 다음을 가정한다.
\begin{enumerate}
\item $k$는 체이다.
\item $\Lambda$는 뇌터 환이다.
\item $\mathfrak q$는 $\mathfrak h_A$ 위에서 극소이다.
\item $\Lambda_\mathfrak q$는 정칙 국소환이다.
\item 체 확대 $\kappa(\mathfrak q)/k$는 분리적이다.
\end{enumerate}
그러면 $k \to A \to \Lambda \supset \mathfrak q$는 해소 가능하다.
\end{lemma}

\begin{proof}
$d = \dim \Lambda_\mathfrak q$로 두자. $R = k[x_1, \ldots, x_d]$로 두자.
다음을 만족하는 $n > 0$을 잡자.
$\mathfrak q^n\Lambda_\mathfrak q \subset \mathfrak h_A\Lambda_\mathfrak q$
$\mathfrak q$가 $\mathfrak h_A$ 위에서 극소이므로 이는 가능하다.
$H_{A/R}$의 생성원 $a_1, \ldots, a_r$을 잡고 다음과 같이 두자.
$$
B = A[x_1, \ldots, x_d, z_{ij}]/(x_i^n - \sum z_{ij}a_j)
$$
$B_{a_j}$는 각각 $R$ 위에서 매끄럽다. 실제로 이는
$A_{a_j}[x_1, \ldots, x_d]$ 위의 다항식 대수이고, $A_{a_j}$는 $k$ 위에서 매끄럽다.
따라서 $B_{x_i}$는 $R$ 위에서 매끄럽다. 보조정리
\ref{smoothing-lemma-improve-presentation}에서 구성한 $R$-대수 사상을
$B \to C$라 하자. 여기에는 $R$-대수 수축 사상 $C \to B$가 딸려 있다.
특히 위 도식에 들어맞는 사상 $C \to \Lambda$가 있다.
구성상 $C_{x_i}$는 $\Omega_{C_{x_i}/R}$가 자유인 매끄러운 $R$-대수이다.
따라서 $x_i^c$가 $C/R$에서 엄밀 표준이 되는 $c > 0$을 찾을 수 있다.
보조정리 \ref{smoothing-lemma-compare-standard}를 보라.
이제 보조정리 \ref{smoothing-lemma-find-sequence}에서처럼
$\pi_1, \ldots, \pi_d \in \Lambda$를 잡자. 여기서
$n = n$, $e = 8c$, $\mathfrak q = \mathfrak q$, $I = \mathfrak h_A$이다.
어떤 $\pi_{ij} \in \Lambda$에 대해
$\pi_i^n = \sum \lambda_{ij} a_j$로 쓰자.
$x_i \mapsto \pi_i$ 및 $z_{ij} \mapsto \lambda_{ij}$로 주어지는
사상 $B \to \Lambda$가 있다. $R = k[x_1, \ldots, x_d]$로 두자.
다음 도식을 생각하자.
$$
\xymatrix{
R \ar[r] & B \ar[rd] \\
k \ar[u] \ar[r] & A \ar[u] \ar[r] & \Lambda
}
$$
이제 $R \to C \to \Lambda \supset \mathfrak q$와 $R$의 원소열
$x_1^c, \ldots, x_d^c$에 보조정리 \ref{smoothing-lemma-lift-solution}을 적용한다.
가정 (2)는 명백하다. 가정 (1)은 $R$에 대해서는 직접 확인할 수 있고,
$\Lambda$에 대해서는 $\pi_1, \ldots, \pi_d$의 선택에 의해 성립한다.
(다음에 유의하자. 만일
$\text{Ann}_\Lambda(\pi) = \text{Ann}_\Lambda(\pi^2)$이면
$\text{Ann}_\Lambda(\pi) = \text{Ann}_\Lambda(\pi^c)$가 모든 $c > 0$에 대해 성립한다.)
따라서 다음을 해소하면 충분하다.
$$
R/(x_1^e, \ldots, x_d^e) \to C/(x_1^e, \ldots, x_d^e) \to
\Lambda/(\pi_1^e, \ldots, \pi_d^e) \supset
\mathfrak q/(\pi_1^e, \ldots, \pi_d^e)
$$
$e = 8c$로 둔다. 보조정리 \ref{smoothing-lemma-delocalize-height-zero}에 의해
$\mathfrak q$에서 국소화한 뒤 이를 해소하면 충분하다.
그런데 $x_1, \ldots, x_d$의 상은 $\Lambda_\mathfrak q$에서 정칙열이므로
$R_\mathfrak p \to \Lambda_\mathfrak q$가 평탄임을 알 수 있다.
Algebra, Lemma \ref{algebra-lemma-flat-over-regular}를 보라. 따라서
$$
R_\mathfrak p/(x_1^e, \ldots, x_d^e) \to
\Lambda_\mathfrak q/(\pi_1^e, \ldots, \pi_d^e)
$$
는 아르틴 국소환들 사이의 평탄한 환 준동형이다.
더욱이 가정에 의해 이 사상은 잉여류체 위에서 분리적 체 확대를 유도한다.
따라서 Algebra, Lemma \ref{algebra-lemma-colimit-syntomic}와
명제 \ref{smoothing-proposition-lift}에 의해 이 사상은 매끄러운 대수들의 여과 여극한이다.
원하는 해의 존재는 Algebra, Lemma \ref{algebra-lemma-when-colimit}에서 따른다.
\end{proof}






\section{비분리적 잉여류체}
\label{smoothing-section-inseparable}

\noindent
이 절에서는 잉여류체 확대가 비분리적인 경우에 국소 문제를 푸는 방법을 설명한다.

\begin{lemma}
\label{smoothing-lemma-helper}
$k$를 표수 $p > 0$인 체라 하자.
$(\Lambda, \mathfrak m, K)$를 아르틴 국소 $k$-대수라 하자.
$\dim H_1(L_{K/k}) < \infty$라고 가정하자.
그러면 $\Lambda$는 아르틴 국소 $k$-대수 $A$들의 여과 여극한으로 쓸 수 있다.
여기서 각 사상 $A \to \Lambda$는 평탄하고,
$\mathfrak m_A \Lambda = \mathfrak m$이며,
$A$는 $k$ 위에서 본질적으로 유한형이다.
\end{lemma}

\begin{proof}
$A \to \Lambda$의 평탄성은 $A \to \Lambda$가 단사임을 함의함에 유의하자.
따라서 이 보조정리는 실제로 $\Lambda$가 이러한 부분환 $A \subset \Lambda$들의
유향 합집합이라고 말한다. $\mathfrak m^n = 0$을 만족하는 최소 정수를 $n$이라 하자.
$n$에 대한 귀납법으로 보조정리를 증명한다. 체 확대는 유한 생성 체 확대들의
합집합이므로 $n = 1$인 경우는 명백하다.

\medskip\noindent
$\mathfrak m$을 생성하는 $\lambda_1, \ldots, \lambda_d \in \mathfrak m$을 잡자.
$K$는 $\mathbf{F}_p$ 위에서 형식적으로 매끄러우므로(Algebra, Lemma
\ref{algebra-lemma-formally-smooth-extensions-easy} 참조), 몫사상
$\Lambda \to K$의 절단인 환 준동형 $\sigma : K \to \Lambda$를 찾을 수 있다.
일반적으로 $\sigma$는 $k$-대수 사상이 {\bf 아니다}. $\sigma$가 주어지면
$$
\Psi_\sigma : K[x_1, \ldots, x_d] \longrightarrow \Lambda
$$
를 $K$의 원소에는 $\sigma$를 적용하고 $x_i$는 $\lambda_i$로 보내도록 정의한다.
주장: 어떤 $\sigma : K \to \Lambda$와 $k$ 위에서 유한 생성인 부분체
$k \subset F \subset K$가 존재하여, $\Lambda$에서 $k$의 상이
$\Psi_\sigma(F[x_1, \ldots, x_d])$에 포함된다.

\medskip\noindent
$\mathfrak m^n = 0$을 만족하는 최소 정수 $n$에 대한 귀납법으로 주장을 증명한다.
$n = 1$인 경우는 명백하다. $n > 1$이면
$I = \mathfrak m^{n - 1}$ 및 $\Lambda' = \Lambda/I$로 두자.
귀납 가정에 의해, $k \to \Lambda \to \Lambda'$의 상이
$A' = \Psi_{\sigma'}(F'[x_1, \ldots, x_d])$에 포함되도록 하는
$\sigma' : K \to \Lambda'$와 유한 생성 부분체
$k \subset F' \subset K$가 주어져 있다고 해도 된다.
유도된 사상을 $\tau' : k \to A'$로 표기하자.
$\sigma'$의 올림 $\sigma : K \to \Lambda$를 잡자. 이는 앞서 언급한
$K/\mathbf{F}_p$의 형식적 매끄러움에 의해 가능하다.
나중에 쓰기 위해, 어떤 미분 $D : K \to I$에 대해 $\sigma$를
$\sigma + D$로 바꿀 수 있음에 유의하자.
$A = F[x_1, \ldots, x_d]/(x_1, \ldots, x_d)^n$로 두자.
그러면 $\Psi_\sigma$는 환 준동형 $\Psi_\sigma : A \to \Lambda$를 유도한다.
이를 몫사상 $\Lambda \to \Lambda'$와 합성하면 멱영 핵을 갖는 전사 사상
$A \to A'$이 유도된다. $\tau'$의 올림 $\tau : k \to A$를 잡자.
$k/\mathbf{F}_p$가 형식적으로 매끄러우므로 이는 가능하다.
이로써 두 사상 $k \to \Lambda$, 즉
$\Psi_\sigma \circ \tau : k \to \Lambda$와 주어진 사상
$i : k \to \Lambda$를 얻는다. 이들은 $I$를 법으로 일치하므로 그 차이는 미분
$\theta = i - \Psi_\sigma \circ \tau : k \to I$이다.
$\sigma$를 $\sigma + D$로 바꾸면 $\theta$는 $\theta - D|_k$로 바뀜에 유의하자.

\medskip\noindent
$k$의 원소들의 집합 $\{y_j\}_{j \in J}$를 그 미분 $\text{d}y_j$들이
$\Omega_{k/\mathbf{F}_p}$의 기저를 이루도록 잡자.
$\mathbf{F}_p \subset k \subset K$에 대한 Jacobi--Zariski 열은
$$
0 \to H_1(L_{K/k}) \to \Omega_{k/\mathbf{F}_p} \otimes K \to
\Omega_{K/\mathbf{F}_p} \to \Omega_{K/k} \to 0
$$
$\dim H_1(L_{K/k}) < \infty$이므로 첫 사상의 상이
$\bigoplus_{j \in J_0} K\text{d}y_j$에 포함되게 하는 유한 부분집합
$J_0 \subset J$를 찾을 수 있다. 따라서 원소 $\text{d}y_j$,
$j \in J \setminus J_0$의 상들은 $\Omega_{K/\mathbf{F}_p}$에서
$K$-선형 독립이다. 그러므로 $\theta - D|_k = \xi \circ \text{d}$가 되도록
$D : K \to I$를 고를 수 있다. 여기서 $\xi$는 합성
$$
\Omega_{k/\mathbf{F}_p} = \bigoplus\nolimits_{j \in J} k \text{d}y_j
\longrightarrow \bigoplus\nolimits_{j \in J_0} k \text{d}y_j
\longrightarrow I
$$
$j \in J_0$에 대해 $f_j = \xi(\text{d}y_j) \in I$라 하자.
위와 같이 $\sigma$를 $\sigma + D$로 바꾸자. 그러면
$\Omega_{k/\mathbf{F}_p}$에서 $\text{d}a = \sum a_j \text{d}y_j$일 때,
$a \in k$에 대해 $\theta(a) = \sum_{j \in J_0} a_j f_j$이다.
$I$는 $\lambda_1, \ldots, \lambda_d$에 관한 전체 차수
$|E| = \sum e_i = n - 1$의 단항식
$\lambda^E = \lambda_1^{e_1} \ldots \lambda_d^{e_d}$들로 생성됨에 유의하자.
$c_{j, E} \in K$인 $f_j = \sum_E c_{j, E} \lambda^E$로 쓰자.
$F'$를 $F = F'(c_{j, E})$로 바꾸면 주장이 성립한다.

\medskip\noindent
주장에서와 같이 $\sigma$와 $F$를 잡자. $\Psi_\sigma$의 핵은 유한 개의 다항식
$g_1, \ldots, g_t \in K[x_1, \ldots, x_d]$로 생성된다. 유한 개의 원소를
첨가하여 $F$를 확대하면, 이 다항식들의 계수가 $F$에 속한다고
가정해도 된다. 이 경우 사상
$A = F[x_1, \ldots, x_d]/(g_1, \ldots, g_t) \to
K[x_1, \ldots, x_d]/(g_1, \ldots, g_t) = \Lambda$는 평탄함이 명백하다.
주장에 의해 $A$는 $\Lambda$의 $k$-부분대수이다.
$K$가 부분체 $F$들의 여과 합집합이므로 $\Lambda$는 이 대수들의
여과 여극한임이 명백하다. 마지막으로 Algebra, Lemma
\ref{algebra-lemma-essentially-of-finite-type-into-artinian-local}에 의해
이 대수들은 $k$ 위에서 본질적으로 유한형이다.
\end{proof}

\begin{lemma}
\label{smoothing-lemma-solution-modulo}
$k$를 표수 $p > 0$인 체라 하자.
$\Lambda$를 기하적으로 정칙인 뇌터 $k$-대수라 하자.
$\mathfrak q \subset \Lambda$를 소 아이디얼이라 하자.
$n \geq 1$을 정수라 하고,
$E \subset \Lambda_\mathfrak q/\mathfrak q^n\Lambda_\mathfrak q$를
유한 부분집합이라 하자. 그러면 다음 성질을 갖는 $m \geq 0$과
$\varphi : k[y_1, \ldots, y_m] \to \Lambda$를 찾을 수 있다.
\begin{enumerate}
\item $\mathfrak p = \varphi^{-1}(\mathfrak q)$로 두면
$\mathfrak q\Lambda_\mathfrak q = \mathfrak p \Lambda_\mathfrak q$
이고 $k[y_1, \ldots, y_m]_\mathfrak p \to \Lambda_\mathfrak q$는 평탄하다.
\item 아르틴 국소환의 준동형들에 의한 분해
$$
k[y_1, \ldots, y_m]_\mathfrak p/\mathfrak p^n k[y_1, \ldots, y_m]_\mathfrak p
\to D \to
\Lambda_\mathfrak q/\mathfrak q^n\Lambda_\mathfrak q
$$
가 존재하며, 첫째 화살표는 본질적으로 매끄럽고 둘째 화살표는 평탄하다.
\item $\mathfrak q^n\Lambda_\mathfrak q$를 법으로 $E$는 $D$에 포함된다.
\end{enumerate}
\end{lemma}

\begin{proof}
$\bar \Lambda = \Lambda_\mathfrak q/\mathfrak q^n\Lambda_\mathfrak q$로 두자.
More on Algebra, Proposition
\ref{more-algebra-proposition-characterization-geometrically-regular}에 의해
$\dim H_1(L_{\kappa(\mathfrak q)/k}) < \infty$임에 유의하자.
$E$를 포함하는 $A \subset \bar \Lambda$를 다음과 같이 잡자. $A$는
$k$ 위에서 본질적으로 유한형인 아르틴 국소환이고, 사상
$A \to \bar \Lambda$는 평탄하며, $\mathfrak m_A$는 $\bar \Lambda$의
극대 아이디얼을 생성한다. 보조정리 \ref{smoothing-lemma-helper}를 보라.
잉여류체를 $F = A/\mathfrak m_A$로 표기하면 $k \subset F \subset K$이다.
$\bar \Lambda$에서 $A$의 원소로 보내지고, 또한
$\text{d}\lambda_1, \ldots, \text{d}\lambda_t$의 상이
$\Omega_{F/k}$의 기저를 이루도록
$\lambda_1, \ldots, \lambda_t \in \Lambda$를 잡자.
$y_j$를 $\lambda_j$로 보내는 사상
$\varphi' : k[y_1, \ldots, y_t] \to \Lambda$를 생각하자.
$\mathfrak p' = (\varphi')^{-1}(\mathfrak q)$로 두자.
More on Algebra, Lemma
\ref{more-algebra-lemma-geometrically-regular-over-field}에 의해 환 준동형
$k[y_1, \ldots, y_t]_{\mathfrak p'} \to \Lambda_\mathfrak q$는 평탄하고,
$\Lambda_\mathfrak q/\mathfrak p' \Lambda_\mathfrak q$는 정칙이다.
따라서 $A \subset \bar \Lambda$로 상이 들어가고,
$\Lambda_\mathfrak q/\mathfrak p' \Lambda_\mathfrak q$의 정칙 매개변수계로
보내지는 원소 $\lambda_{t + 1}, \ldots, \lambda_m \in \Lambda$를 더 고를 수 있다.
이로써 성질 (1)을 갖는 $\varphi : k[y_1, \ldots, y_m] \to \Lambda$를 얻고,
$k[y_1, \ldots, y_m]_\mathfrak p/\mathfrak p^n k[y_1, \ldots, y_m]_\mathfrak p
\to \bar\Lambda$
는 $A$를 통해 분해된다. 따라서
$k[y_1, \ldots, y_m]_\mathfrak p/\mathfrak p^n k[y_1, \ldots, y_m]_\mathfrak p
\to A$는 Algebra, Lemma
\ref{algebra-lemma-flatness-descends-more-general}에 의해 평탄하다.
구성상 잉여류체 확대 $F/\kappa(\mathfrak p)$는 유한 생성이고
$\Omega_{F/\kappa(\mathfrak p)} = 0$이다. 따라서 More on Algebra, Lemma
\ref{more-algebra-lemma-cartier-equality}에 의해 유한 분리적이다. 그러므로
$k[y_1, \ldots, y_m]_\mathfrak p/\mathfrak p^n k[y_1, \ldots, y_m]_\mathfrak p
\to A$
는 Algebra, Lemma
\ref{algebra-lemma-essentially-of-finite-type-into-artinian-local}에 의해 유한이다.
마지막으로 Algebra, Lemma \ref{algebra-lemma-characterize-etale}에 의해
이 사상이 에탈임을 알 수 있다. 에탈 환 준동형은 물론 본질적으로
매끄러우므로 원하는 결론을 얻는다.
\end{proof}

\begin{lemma}
\label{smoothing-lemma-enlarge-solution-modulo}
$\varphi : k[y_1, \ldots, y_m] \to \Lambda$, $n$, $\mathfrak q$,
$\mathfrak p$ 및
$$
k[y_1, \ldots, y_m]_\mathfrak p/\mathfrak p^n \to
D \to \Lambda_\mathfrak q/\mathfrak q^n \Lambda_\mathfrak q
$$
가 보조정리 \ref{smoothing-lemma-solution-modulo}에서와 같다고 하자. 그러면 임의의
$\lambda \in \Lambda \setminus \mathfrak q$에 대해 어떤 정수 $q > 0$과 분해
$$
k[y_1, \ldots, y_m]_\mathfrak p/\mathfrak p^n \to
D \to D' \to \Lambda_\mathfrak q/\mathfrak q^n \Lambda_\mathfrak q
$$
가 존재한다. 여기서 $D \to D'$는 아르틴 국소환들 사이의 본질적으로
매끄러운 사상이고, 마지막 화살표는 평탄하며, $\lambda^q$는 $D'$에 속한다.
\end{lemma}

\begin{proof}
$\bar \Lambda = \Lambda_\mathfrak q/\mathfrak q^n\Lambda_\mathfrak q$로 두자.
$\bar \lambda$를 $\bar \Lambda$에서 $\lambda$의 상이라 하고,
$\alpha \in \kappa(\mathfrak q)$를 잉여류체에서 $\lambda$의 상이라 하자.
$D$의 잉여류체를 $k \subset F \subset \kappa(\mathfrak q)$라 하자.
$\alpha$가 $F$에 속하면 $x \bar\lambda = 1 \bmod \mathfrak q$를 만족하는
$x \in D$를 찾을 수 있다. 따라서 $q$가 $p$로 나누어떨어지면
$(x \bar \lambda)^q = 1 \bmod (\mathfrak q)^q$이다. 그러므로
$\bar\lambda^q$는 $D$에 속한다. $\alpha$가 $F$ 위에서 초월적이면,
$D$와 $\bar \lambda$로 생성한 부분환을
$\mathfrak m = D[\bar \lambda] \cap \mathfrak q \bar \Lambda$에서 국소화한
$D' = (D[\bar \lambda])_\mathfrak m$으로 둘 수 있다. 이 경우
$D[\bar \lambda]$는 실제로 $D$ 위의 다항식 대수이므로 이것으로 충분하다.
마지막으로 $\lambda \bmod \mathfrak q$가 $F$ 위에서 대수적이면,
$\alpha^q$가 $F$ 위에서 분리 대수적이 되게 하는 $p$-거듭제곱 $q$를
찾을 수 있다. Fields, Section \ref{fields-section-algebraic}을 보라.
$D$와 $\bar\Lambda$는 헨젤 국소환임에 유의하자. Algebra, Lemma
\ref{algebra-lemma-local-dimension-zero-henselian}을 보라.
잉여류체 확대가 $F(\alpha^q)/F$인 유한 에탈 확대를 $D \to D'$라 하자.
Algebra, Lemma \ref{algebra-lemma-henselian-cat-finite-etale}를 보라.
$\bar\Lambda$가 헨젤이고 $F(\alpha^q)$가 그 잉여류체에 포함되므로
분해 $D' \to \bar \Lambda$를 찾을 수 있다. 논증의 첫 부분에 의해
어떤 $q' > 0$에 대해 $\bar\lambda^{qq'} \in D'$임을 알 수 있다.
\end{proof}

\begin{lemma}
\label{smoothing-lemma-resolve-general}
상황 \ref{smoothing-situation-local}과 같은
$k \to A \to \Lambda \supset \mathfrak q$를 생각하자. 다음을 가정한다.
\begin{enumerate}
\item $k$는 표수 $p > 0$인 체이다.
\item $\Lambda$는 뇌터 환이고 $k$ 위에서 기하적으로 정칙이다.
\item $\mathfrak q$는 $\mathfrak h_A$ 위에서 극소이다.
\end{enumerate}
그러면 $k \to A \to \Lambda \supset \mathfrak q$는 해소 가능하다.
\end{lemma}

\begin{proof}
다음 단계들을 주어진 순서로 수행하여 보조정리를 증명한다.
각 단계의 정당화는 아래에서 제시한다.
\begin{enumerate}
\item
\label{smoothing-item-power}
다음을 만족하는 정수 $N > 0$을 잡는다.
$\mathfrak q^N\Lambda_\mathfrak q \subset H_{A/k}\Lambda_\mathfrak q$.
\item
\label{smoothing-item-generators}
아이디얼 $H_{A/R}$의 생성원 $a_1, \ldots, a_t \in A$를 잡는다.
\item
\label{smoothing-item-dimension}
$d = \dim(\Lambda_\mathfrak q)$로 둔다.
\item
\label{smoothing-item-standardizer}
$B = A[x_1, \ldots, x_d, z_{ij}]/(x_i^{2N} - \sum z_{ij}a_j)$로 둔다.
\item
\label{smoothing-item-elkik}
$B$를 $k[x_1, \ldots, x_d]$-대수로 보고, 보조정리
\ref{smoothing-lemma-improve-presentation}에서와 같은 $B \to C$를 잡는다.
또한 절단 $C \to B$를 얻는다.
\item
\label{smoothing-item-strictly-standard}
각 $x_i^c$가 $k[x_1, \ldots, x_d]$ 위의 $C$에서 엄밀 표준이 되게 하는
$c > 0$을 잡는다.
\item
\label{smoothing-item-set-n}
$n = N + dc$ 및 $e = 8c$로 둔다.
\item
\label{smoothing-item-set-E}
$E \subset \Lambda_\mathfrak q/\mathfrak q^n\Lambda_\mathfrak q$를
$k$-대수로서 $A$의 생성원들의 상으로 이루어진 집합이라 하자.
\item
\label{smoothing-item-NP}
정수 $m$, $k$-대수 사상
$\varphi : k[y_1, \ldots, y_m] \to \Lambda$, 그리고 아르틴 국소환들에 의한 분해
$$
k[y_1, \ldots, y_m]_\mathfrak p/\mathfrak p^n k[y_1, \ldots, y_m]_\mathfrak p
\to D \to
\Lambda_\mathfrak q/\mathfrak q^n\Lambda_\mathfrak q
$$
를 다음과 같이 잡는다. 첫째 화살표는 본질적으로 매끄럽고, 둘째 화살표는
평탄하며, $E$는 $D$에 포함된다. 또한 $\mathfrak p = \varphi^{-1}(\mathfrak q)$로 두면
$k[y_1, \ldots, y_m]_\mathfrak p \to \Lambda_\mathfrak q$가 평탄하고
$\mathfrak p \Lambda_\mathfrak q = \mathfrak q \Lambda_\mathfrak q$이다.
\item
\label{smoothing-item-choose-pii}
$k[y_1, \ldots, y_m]_\mathfrak p$의 정칙 매개변수계로 보내지는
$\pi_1, \ldots, \pi_d \in \mathfrak p$를 잡는다.
\item
\label{smoothing-item-set-R}
$R = k[y_1, \ldots, y_m, t_1, \ldots, t_m]$ 및 $\gamma_i = \pi_i t_i$로 둔다.
\item
\label{smoothing-item-modify-pii}
필요하면 $i = 1, \ldots, d$에 대해 다음이 성립하도록 $\pi_i$의 선택을 바꾼다.
$$
\text{Ann}_{R/(\gamma_1^e, \ldots, \gamma_{i - 1}^e)R}(\gamma_i)
=
\text{Ann}_{R/(\gamma_1^e, \ldots, \gamma_{i - 1}^e)R}(\gamma_i^2)
$$
\item
\label{smoothing-item-choose-deltai}
$\delta_1, \ldots, \delta_d \in \Lambda$, $\delta_i \not \in \mathfrak q$와 분해
$D \to D' \to \Lambda_\mathfrak q/\mathfrak q^n\Lambda_\mathfrak q$
가 존재한다. 여기서 $D'$는 아르틴 국소환이고, $D \to D'$는 본질적으로
매끄러우며, 사상
$D' \to \Lambda_\mathfrak q/\mathfrak q^n\Lambda_\mathfrak q$
는 평탄하다. 또한 $\pi_i' = \delta_i \pi_i$로 두면
$i = 1, \ldots, d$에 대해 다음이 성립한다.
\begin{enumerate}
\item
$\Lambda$에서 $(\pi_i')^{2N} = \sum a_j\lambda_{ij}$이고,
$\lambda_{ij} \bmod \mathfrak q^n\Lambda_\mathfrak q$는 $D'$의 원소이다.
\item $\text{Ann}_{\Lambda/({\pi'}_1^e, \ldots, {\pi'}_{i - 1}^e)}({\pi'}_i) =
\text{Ann}_{\Lambda/({\pi'}_1^e, \ldots, {\pi'}_{i - 1}^e)}({\pi'}_i^2)$,
\item $\delta_i \bmod \mathfrak q^n\Lambda_\mathfrak q$는 $D'$의 원소이다.
\end{enumerate}
\item
\label{smoothing-item-map-B-Lambda}
$x_i$를 $\pi'_i$로, $z_{ij}$를 위에서 구한 $\lambda_{ij}$로 보내어
$B \to \Lambda$를 정의한다. 사상 $B \to \Lambda$를 수축 사상
$C \to B$와 합성하여 $C \to \Lambda$를 정의한다.
\item
\label{smoothing-item-set-map-R}
$k[y_1, \ldots, y_m]$에는 $\varphi$를 적용하고 $t_i$는 $\delta_i$로 보내어
$R \to \Lambda$를 정의한다. 또한 사상
$$
k[x_1, \ldots, x_d]
\longrightarrow
R = k[y_1, \ldots, y_m, t_1, \ldots, t_d]
$$
을 $x_i$를 $\gamma_i = \pi_i t_i$로 보내도록 정의한다.
\item
\label{smoothing-item-first-solve}
다음을 해소하면 충분하다.
$$
R
\to
C \otimes_{k[x_1, \ldots, x_d]} R
\to
\Lambda \supset \mathfrak q
$$
\item
\label{smoothing-item-set-I}
$I = (\gamma_1^e, \ldots, \gamma_d^e) \subset R$로 둔다.
\item
\label{smoothing-item-second-resolve}
다음을 해소하면 충분하다.
$$
R/I
\to
C \otimes_{k[x_1, \ldots, x_d]} R/I
\to
\Lambda/I\Lambda \supset \mathfrak q/I\Lambda
$$
\item
\label{smoothing-item-set-r}
$\mathfrak q$의 역상을
$\mathfrak r \subset R = k[y_1, \ldots, y_m, t_1, \ldots, t_d]$로 표기한다.
\item
\label{smoothing-item-third-resolve}
다음을 해소하면 충분하다.
$$
(R/I)_\mathfrak r \to
C \otimes_{k[x_1, \ldots, x_d]} (R/I)_\mathfrak r \to
\Lambda_\mathfrak q/I\Lambda_\mathfrak q
\supset
\mathfrak q\Lambda_\mathfrak q/I\Lambda_\mathfrak q
$$
\item
\label{smoothing-item-fourth-resolve}
$k[y_1, \ldots, y_m]$에서 $J = (\pi_1^e, \ldots, \pi_d^e)$로 둔다.
\item
\label{smoothing-item-fifth-resolve}
다음을 해소하면 충분하다.
$$
(R/JR)_\mathfrak p \to
C \otimes_{k[x_1, \ldots, x_d]} (R/JR)_\mathfrak p \to
\Lambda_\mathfrak q/J\Lambda_\mathfrak q
\supset
\mathfrak q\Lambda_\mathfrak q/J\Lambda_\mathfrak q
$$
\item
\label{smoothing-item-sixth-resolve}
다음을 해소하면 충분하다.
$$
(R/\mathfrak p^nR)_\mathfrak p \to
C \otimes_{k[x_1, \ldots, x_d]} (R/\mathfrak p^nR)_\mathfrak p \to
\Lambda_\mathfrak q/\mathfrak q^n\Lambda_\mathfrak q
\supset
\mathfrak q\Lambda_\mathfrak q/\mathfrak q^n\Lambda_\mathfrak q
$$
\item
\label{smoothing-item-seventh-resolve}
다음을 해소하면 충분하다.
$$
(R/\mathfrak p^nR)_\mathfrak p \to
B \otimes_{k[x_1, \ldots, x_d]} (R/\mathfrak p^nR)_\mathfrak p \to
\Lambda_\mathfrak q/\mathfrak q^n\Lambda_\mathfrak q
\supset
\mathfrak q\Lambda_\mathfrak q/\mathfrak q^n\Lambda_\mathfrak q
$$
\item
\label{smoothing-item-eighth-resolve}
주어진 사상
$k[y_1, \ldots, y_m]_\mathfrak p/\mathfrak p^n k[y_1, \ldots, y_m]_\mathfrak p
\to D'$와 $t_i$를 $t_i$로 보내는 사상으로 $D'[t_1, \ldots, t_d]$에
$R_\mathfrak p/\mathfrak p^nR_\mathfrak p$-대수 구조를 준다. 분해
$$
B \otimes_{k[x_1, \ldots, x_d]} (R/\mathfrak p^nR)_\mathfrak p
\to D'[t_1, \ldots, t_d] \to
\Lambda_\mathfrak q/\mathfrak q^n\Lambda_\mathfrak q
$$
를 찾으면 충분하다. 여기서 둘째 화살표는 $t_i$를 $\delta_i$로 보내고,
주어진 준동형 $D' \to \Lambda_\mathfrak q/\mathfrak q^n\Lambda_\mathfrak q$를 유도한다.
\item
\label{smoothing-item-done}
위에서 $D'$를 선택한 방식에 의해 이러한 분해가 존재한다.
\end{enumerate}
이제 각 단계를 정당화한다. 단, 단지 표기만 도입하는 단계는 생략한다.

\medskip\noindent
단계 (\ref{smoothing-item-power}). $\mathfrak q$가
$\mathfrak h_A = \sqrt{H_{A/k}\Lambda}$ 위에서 극소이므로 가능하다.

\medskip\noindent
단계 (\ref{smoothing-item-strictly-standard}). $A_{a_i}$가 $k$ 위에서 매끄러움에
유의하자. 따라서 $A_{a_j}[x_1, \ldots, x_d]$ 위의 다항식 대수와 동형인
$B_{a_j}$는 $k[x_1, \ldots, x_d]$ 위에서 매끄럽다. 그러므로
$B_{x_i}$는 $k[x_1, \ldots, x_d]$ 위에서 매끄럽다. 보조정리
\ref{smoothing-lemma-improve-presentation}에 의해 $C_{x_i}$는
$k[x_1, \ldots, x_d]$ 위에서 매끄럽고 그 미분 모듈은 유한 자유이다.
따라서 보조정리 \ref{smoothing-lemma-compare-standard}에 의해 $x_i$의 어떤 거듭제곱은
$k[x_1, \ldots, x_n]$ 위의 $C$에서 엄밀 표준이다.

\medskip\noindent
단계 (\ref{smoothing-item-NP}). 보조정리 \ref{smoothing-lemma-solution-modulo}를 적용하면 된다.

\medskip\noindent
단계 (\ref{smoothing-item-choose-pii}). 구성상
$k[y_1, \ldots, y_m]_\mathfrak p \to \Lambda_\mathfrak q$는
평탄하고 $\mathfrak p \Lambda_\mathfrak q = \mathfrak q \Lambda_\mathfrak q$이다.
따라서 Algebra, Lemma
\ref{algebra-lemma-dimension-base-fibre-equals-total}에 의해
$\dim(k[y_1, \ldots, y_m]_\mathfrak p) = d$임을 알 수 있다.
그러므로 $\Lambda_\mathfrak q$의 정칙 매개변수계로 보내지는
$\pi_1, \ldots, \pi_d \in \Lambda$를 찾을 수 있다.

\medskip\noindent
단계 (\ref{smoothing-item-modify-pii}). Algebra, Lemma
\ref{algebra-lemma-regular-ring-CM}에 의해 열 $\pi_1, \ldots, \pi_d$의
임의의 순열은 $k[y_1, \ldots, y_m]_\mathfrak p$에서 정칙열이다.
따라서 $\gamma_1 = \pi_1 t_1, \ldots, \gamma_d = \pi_d t_d$는
$R_\mathfrak p = k[y_1, \ldots, y_m]_\mathfrak p[t_1, \ldots, t_d]$에서
정칙열이다. Algebra, Lemma
\ref{algebra-lemma-regular-sequence-in-polynomial-ring}을 보라.
$S = k[y_1, \ldots, y_m] \setminus \mathfrak p$로 두면
$R_\mathfrak p = S^{-1}R$이다. $\pi_i$에 $S$의 원소들을 곱해도
$\pi_1, \ldots, \pi_d$와 $\gamma_1, \ldots, \gamma_d$가 정칙열인 성질은
보존됨에 유의하자. 이제
$$
\text{Ann}_{R/(\gamma_1^e, \ldots, \gamma_{i - 1}^e)R}(\gamma_i)
=
\text{Ann}_{R/(\gamma_1^e, \ldots, \gamma_{i - 1}^e)R}(\gamma_i^2)
$$
가 어떤 $t \in \{0, \ldots, d\}$에 대해 $i = 1, \ldots, t$일 때
성립한다고 하자. Algebra, Lemma
\ref{algebra-lemma-regular-sequence-powers}에 의해
$\gamma_1^e, \ldots, \gamma_t^e, \gamma_{t + 1}$은
$S^{-1}R$에서 정칙열이다. 따라서
$$
\text{Ann}_{S^{-1}R/(\gamma_1^e, \ldots, \gamma_{i - 1}^e)}(\gamma_i) =
\text{Ann}_{S^{-1}R/(\gamma_1^e, \ldots, \gamma_{i - 1}^e)}(\gamma_i^2).
$$
그러므로
$$
\text{Ann}_{R/(\gamma_1^e, \ldots, \gamma_t^e)R}(\gamma_{t + 1})
=
\text{Ann}_{R/(\gamma_1^e, \ldots, \gamma_t^e)R}(\gamma_{t + 1}^2)
$$
를 얻는다. 여기서는 보조정리 \ref{smoothing-lemma-ogoma}에 따라 어떤 $s \in S$에 대해
$\pi_{t + 1}$을 $s\pi_{t + 1}$로 바꾸었다. $t$에 대한 귀납법으로
원하는 열을 얻는다.

\medskip\noindent
단계 (\ref{smoothing-item-choose-deltai}). $S = \Lambda \setminus \mathfrak q$로 두면
$\Lambda_\mathfrak q = S^{-1}\Lambda$이다.
$\bar \Lambda = \Lambda_\mathfrak q/\mathfrak q^n \Lambda_\mathfrak q$로 두자.
어떤 $t \in \{0, \ldots, d\}$, $\delta_1, \ldots, \delta_t \in S$ 및
(\ref{smoothing-item-choose-deltai})에서와 같은 분해 $D \to D' \to \bar \Lambda$가
주어져 있고, $i = 1, \ldots, t$에 대해 (a), (b), (c)가 성립한다고 하자.
(\ref{smoothing-item-power})에 의해
$\mathfrak q^N \Lambda_\mathfrak q \subset H_{A/k}\Lambda_\mathfrak q$이므로
$\pi_{t + 1}^N \in H_{A/k}\Lambda_\mathfrak q$이다. 따라서
$\pi_{t + 1}^N \in H_{A/k} \bar\Lambda$이다. 또한
$D' \to \bar \Lambda$가 충실히 평탄하므로
$\pi_{t + 1}^N \in H_{A/k}D'$이다. Algebra, Lemma
\ref{algebra-lemma-faithfully-flat-universally-injective}를 보라.
$H_{A/k} = (a_1, \ldots, a_t)$임을 상기하자.
$D'$에서 $\pi_{t + 1}^N = \sum a_j d_j$라 쓰고,
$d_j \in D'$의 올림 $c_j \in \Lambda_\mathfrak q$를 잡자. 그러면
$\epsilon \in \mathfrak q^n\Lambda_\mathfrak q \subset
\mathfrak q^{n - N}H_{A/k}\Lambda_\mathfrak q$인
$\pi_{t + 1}^N = \sum c_j a_j + \epsilon$을 얻는다. 어떤
$c'_j \in \mathfrak q^{n - N}\Lambda_\mathfrak q$에 대해
$\epsilon = \sum a_j c'_j$로 쓰자. 따라서
$\pi_{t + 1}^{2N} = \sum (\pi_{t + 1}^N c_j + \pi_{t + 1}^N c'_j) a_j$이다.
$\pi_{t + 1}^Nc'_j$의 상은 $\bar \Lambda$에서 영임에 유의하자.
이 자명하지만 핵심적인 관찰은 나중에 (a)가 성립함을 보장한다.
이제 다음 두 성질을 만족하는 $s \in S$와
$\mu_{t + 1j} \in \Lambda$를 잡자. 한편으로
$S^{-1}\Lambda$에서
$\pi_{t + 1}^N c_j + \pi_{t + 1}^N c'_j = \mu_{t + 1j}/s^{2N}$이고,
다른 한편으로 $\Lambda$에서
$(s \pi_{t + 1})^{2N} = \sum \mu_{t + 1j}a_j$이다
(사소한 세부사항은 생략한다). 더 나아가 $s$를 그 거듭제곱으로 바꾸고
$D'$를 확대하여 $s$의 상이 $D'$의 원소라고 해도 된다.
이 선택들 아래에서 $\mu_{t + 1j}$의 상은 $D'$의 원소인 $s^{2N}d_j$이다.
$\varphi$의 선택에 의해 $\pi_1, \ldots, \pi_d$는
$S^{-1}\Lambda$에서 매개변수들의 정칙열임에 유의하자. 따라서 Algebra, Lemma
\ref{algebra-lemma-regular-ring-CM}에 의해 $\pi_1, \ldots, \pi_d$는
$\Lambda_\mathfrak q$에서 정칙열을 이룬다. 그러므로 Algebra, Lemma
\ref{algebra-lemma-regular-sequence-powers}에 의해
${\pi'}_1^e, \ldots, {\pi'}_t^e, s\pi_{t + 1}$은
$S^{-1}\Lambda$에서 정칙열이다. 따라서
$$
\text{Ann}_{S^{-1}\Lambda/({\pi'}_1^e, \ldots, {\pi'}_t^e)}(s\pi_{t + 1}) =
\text{Ann}_{S^{-1}\Lambda/({\pi'}_1^e, \ldots, {\pi'}_t^e)}((s\pi_{t + 1})^2).
$$
보조정리 \ref{smoothing-lemma-ogoma}를 적용하면 다음을 만족하는 $s' \in S$를 찾을 수 있다.
$$
\text{Ann}_{\Lambda/({\pi'}_1^e, \ldots, {\pi'}_t^e)}((s')^qs\pi_{t + 1})
=
\text{Ann}_{\Lambda/({\pi'}_1^e, \ldots, {\pi'}_t^e)}(((s')^qs\pi_{t + 1})^2).
$$
이 식은 임의의 $q > 0$에 대해 성립한다. 보조정리
\ref{smoothing-lemma-enlarge-solution-modulo}에 의해 $q$를 선택하고 $D'$를 확대하여
$(s')^q$의 상이 $D'$의 원소라고 할 수 있다.
$\delta_{t + 1} = (s')^qs$로 두면 $i = 1, \ldots, t + 1$에 대해
(a), (b), (c)가 성립한다. (a)에 대해서는
$\lambda_{t + 1j} = (s')^{2Nq}\mu_{t + 1j}$가 조건을 만족함에 유의하자.
$t$에 대한 귀납법으로 원하는 결론을 얻는다.

\medskip\noindent
단계 (\ref{smoothing-item-first-solve}). 구성상
$H_{(C \otimes_{k[x_1, \ldots, x_d]} R)/R} \Lambda$의 근기는
$\mathfrak h_A$를 포함한다. 실제로 원소 $a_j \in H_{A/k}$는
$H_{B/k[x_1, \ldots, x_n]}$의 원소로 보내지고, 따라서
$H_{C/k[x_1, \ldots, x_n]}$의 원소로 보내지며, 따라서
$a_j \otimes 1$은 $H_{C \otimes_{k[x_1, \ldots, x_d]} R/R}$의 원소로 보내진다.
더욱이 다음 상황의 해 $C \otimes_{k[x_1, \ldots, x_n]} R \to T \to \Lambda$가
주어져 있다고 하자.
$$
R
\to
C \otimes_{k[x_1, \ldots, x_d]} R
\to
\Lambda \supset \mathfrak q
$$
그러면 $R$가 $k$ 위에서 매끄러우므로 $H_{T/R} \subset H_{T/k}$이다.
따라서 $T$는 원래 상황 $k \to A \to \Lambda \supset \mathfrak q$의 해이기도 하다.

\medskip\noindent
단계 (\ref{smoothing-item-second-resolve}).
$R \to C \otimes_{k[x_1, \ldots, x_d]} R
\to \Lambda \supset \mathfrak q$와 원소열
$\gamma_1^c, \ldots, \gamma_d^c$에 보조정리 \ref{smoothing-lemma-lift-solution}을 적용하면 된다.
$x_i^c$가 $k[x_1, \ldots, x_d]$ 위의 $C$에서 엄밀 표준이므로,
보조정리 \ref{smoothing-lemma-strictly-standard-base-change}에 의해 원소 $\gamma_i^c$는
$R$ 위의 $C \otimes_{k[x_1, \ldots, x_d]} R$에서 엄밀 표준임에 유의하자.
보조정리 \ref{smoothing-lemma-lift-solution}의 다른 가정은 단계
(\ref{smoothing-item-modify-pii})와 (\ref{smoothing-item-choose-deltai})에 의해 성립한다.

\medskip\noindent
단계 (\ref{smoothing-item-third-resolve}). 보조정리 \ref{smoothing-lemma-delocalize-height-zero}를
단계 (\ref{smoothing-item-second-resolve})의 상황에 적용한다. 나머지 논증에서는
목표 환이 아르틴 국소환이므로, 원천 환 위에서 매끄러운 대수 $T$를
통한 분해를 찾는다.

\medskip\noindent
단계 (\ref{smoothing-item-fifth-resolve}).
$C \otimes_{k[x_1, \ldots, x_d]} (R/JR)_\mathfrak p \to
T \to \Lambda_\mathfrak q/J\Lambda_\mathfrak q$가 다음 상황의 해라고 하자.
$$
(R/JR)_\mathfrak p \to
C \otimes_{k[x_1, \ldots, x_d]} (R/JR)_\mathfrak p \to
\Lambda_\mathfrak q/J\Lambda_\mathfrak q
\supset
\mathfrak q\Lambda_\mathfrak q/J\Lambda_\mathfrak q
$$
그러면 $C \otimes_{k[x_1, \ldots, x_d]} (R/I)_\mathfrak r \to T_\mathfrak r \to
\Lambda_\mathfrak q/I\Lambda_\mathfrak q$는 단계
(\ref{smoothing-item-third-resolve})의 상황에 대한 해이다.

\medskip\noindent
단계 (\ref{smoothing-item-sixth-resolve}). 우리가 택한 $n = N + dc$는 충분히 커서
$\mathfrak p^nk[y_1, \ldots, y_m]_\mathfrak p \subset J_\mathfrak p$
이고 $\mathfrak q^n \Lambda_\mathfrak q \subset J\Lambda_\mathfrak q$이다.
따라서
$C \otimes_{k[x_1, \ldots, x_d]} (R/\mathfrak p^nR)_\mathfrak p \to
T \to \Lambda_\mathfrak q/\mathfrak q^n\Lambda_\mathfrak q$
가 (\ref{smoothing-item-fifth-resolve}의 해라면, $T/JT$를
(\ref{smoothing-item-sixth-resolve})의 해로 취할 수 있다.

\medskip\noindent
단계 (\ref{smoothing-item-seventh-resolve}). $R$-대수의 범주에서 절단
$C \to B$를 갖기 때문에 성립한다.

\medskip\noindent
단계 (\ref{smoothing-item-eighth-resolve}). $D'$는 아르틴 국소환
$k[y_1, \ldots, y_m]_\mathfrak p/\mathfrak p^n k[y_1, \ldots, y_m]_\mathfrak p$
위에서 본질적으로 매끄럽고,
$$
R_\mathfrak p/\mathfrak p^nR_\mathfrak p =
k[y_1, \ldots, y_m]_\mathfrak p/
\mathfrak p^n k[y_1, \ldots, y_m]_\mathfrak p[t_1, \ldots, t_d].
$$
이기 때문에 성립한다. 따라서 $D'[t_1, \ldots, t_d]$는 매끄러운
$R_\mathfrak p/\mathfrak p^nR_\mathfrak p$-대수들의 여과 여극한이고,
$B \otimes_{k[x_1, \ldots, x_d]} (R_\mathfrak p/\mathfrak p^nR_\mathfrak p)$
는 그중 하나를 통해 분해된다.

\medskip\noindent
단계 (\ref{smoothing-item-done}). 증명의 마지막 묘점은 단순히 사상
$B \to D'$를 사용할 수 없다는 것이다. 이 사상은 $x_i$를 $D'$에서
$\pi_i'$의 상으로 보내고 $z_{ij}$를 $D'$에서 $\lambda_{ij}$의 상으로 보낸다.
그 이유는 도식
$$
\xymatrix{
B \ar[r] & D'[t_1, \ldots, t_d] \\
k[x_1, \ldots, x_d] \ar[r] \ar[u] &
R_\mathfrak p/\mathfrak p^nR_\mathfrak p \ar[u]
}
$$
이 가환해야 하고, 합성
$B \to D'[t_1, \ldots, t_d] \to
\Lambda_\mathfrak q/\mathfrak q^n\Lambda_\mathfrak q$
이 (\ref{smoothing-item-map-B-Lambda})의 사상이어야 하기 때문이다.
이를 위해서는 $x_i$를 $D'[t_1, \ldots, t_d]$에서
$\pi_i t_i$의 상으로 보내야 한다. 따라서 $z_{ij}$를
$D'[t_1, \ldots, t_d]$에서 $\lambda_{ij} t_i^{2N} / \delta_i^{2N}$의
상으로 보내면 모든 것이 명백해진다.
\end{proof}







\section{주정리}
\label{smoothing-section-main}

\noindent
이 절에서는 지금까지의 논의를 마무리한다.

\begin{theorem}[Popescu]
\label{smoothing-theorem-popescu}
뇌터 환의 임의의 정칙 준동형은 매끄러운 환 사상들의 여과 여극한이다.
\end{theorem}

\begin{proof}
보조정리 \ref{smoothing-lemma-reduce-to-field}에 의해 $\Lambda$가 뇌터이고
$k$ 위에서 기하적으로 정칙인 경우의 $k \to \Lambda$에 대해 증명하면 충분하다.
$A$가 유한형 $k$-대수인 분해 $k \to A \to \Lambda$를 잡자.
보조정리 \ref{smoothing-lemma-final-solve}에 따라 $B$가 유한형이고
$\mathfrak h_B = \Lambda$인 분해 $A \to B \to \Lambda$를 구성하면 충분하다.
따라서 아이디얼 $\mathfrak h_A$에 대해 뇌터 귀납법을 적용할 수 있다.
소아이디얼 $\mathfrak q \supset \mathfrak h_A$를 잡자. 이때
$\mathfrak q$는 $\mathfrak h_A$ 위에서 극소이다.
이제 $k \to A \to \Lambda \supset \mathfrak q$를 해소하면 충분하다
(상황 \ref{smoothing-situation-local} 뒤의 본문에서 정의한 의미로).
$k$의 표수가 영이면 이는 보조정리 \ref{smoothing-lemma-resolve-special}에서 따른다.
$k$의 표수가 $p > 0$이면 이는 보조정리 \ref{smoothing-lemma-resolve-general}에서 따른다.
\end{proof}




\section{G-환의 근사 성질}
\label{smoothing-section-approximation-G-rings}

\noindent
$R$를 뇌터 국소환이라 하자. 이 경우 $R$가 G-환일 필요충분조건은
환 사상 $R \to R^\wedge$가 정칙인 것이다. More on Algebra, Lemma
\ref{more-algebra-lemma-check-G-ring-maximal-ideals}를 보라.
이 경우 $R$의 헨젤화 $R^h$와 엄밀 헨젤화 $R^{sh}$도 G-환이다.
More on Algebra, Lemma \ref{more-algebra-lemma-henselization-G-ring}을 보라.
더욱이 체, 완비 국소환, $\mathbf{Z}$ 또는 표수 영인 데데킨트 환 위에서
본질적으로 유한형인 임의의 대수는 G-환이다. More on Algebra, Proposition
\ref{more-algebra-proposition-ubiquity-G-ring}을 보라.
따라서 아래 결과를 적용할 수 있는 환은 충분히 많다.

\medskip\noindent
$R$를 환이라 하고 $f_1, \ldots, f_m \in R[x_1, \ldots, x_n]$이라 하자.
$S$를 $R$-대수라 하자. 이 상황에서 벡터
$(a_1, \ldots, a_n) \in S^n$가 {\it $S$에서의 해}라는 것은 다음이
성립한다는 것과 동치이다.
$$
f_j(a_1, \ldots, a_n) = 0 \text{ in } S, \text{ for }
j = 1, \ldots, m
$$
물론 대수기하학에서 중요한 문제 하나는 다항방정식계가 언제 해를 갖는지
밝히는 것이다. 다음 정리에 따르면 뇌터 국소환의 완비화에서 해가 존재하는
것만으로 그 환의 헨젤화에서 해가 존재함을 보이기에 충분한 경우가 많다.

\begin{theorem}
\label{smoothing-theorem-approximation-property}
$R$를 뇌터 국소환이라 하고
$f_1, \ldots, f_m \in R[x_1, \ldots, x_n]$이라 하자.
$(a_1, \ldots, a_n) \in (R^\wedge)^n$가 $R^\wedge$에서의 해라고 하자.
$R$가 헨젤 G-환이면, 모든 정수 $N$에 대해
$a_i - b_i \in \mathfrak m^NR^\wedge$를 만족하는 $R$에서의 해
$(b_1, \ldots, b_n) \in R^n$가 존재한다.
\end{theorem}

\begin{proof}
$a_i - c_i \in \mathfrak m^N$을 만족하는 원소 $c_i \in R$를 잡자.
생성원 $\mathfrak m^N = (d_1, \ldots, d_M)$을 잡고
$a_i = c_i + \sum a_{i, l} d_l$로 쓰자.
다항식환 $R[x_{i, l}]$과 다음 원소들을 생각하자.
$$
g_j = f_j(c_1 + \sum x_{1, l} d_l , \ldots, c_n + \sum x_{n, l} d_{n, l})
\in R[x_{i, l}]
$$
연립방정식 $g_j = 0$은 해 $(a_{i, l})$을 갖는다.
$g_j$가 $R$에서 해 $(b_{i, l})$을 가짐을 보일 수 있다고 하자.
그러면 $b_i = c_i + \sum b_{i, l}d_l$은 $\mathfrak m^N$을 법으로
$a_i$와 합동인 $f_j = 0$의 해이다. 따라서 $R^\wedge$ 위의 해의 존재가
$R$ 위의 해의 존재를 함의함을 보이면 충분하다.

\medskip\noindent
$A \subset R^\wedge$를 $a_1, \ldots, a_n$으로 생성되는 $R$-부분대수라 하자.
$R$가 G-환, 즉 $R \to R^\wedge$가 정칙이라고 가정했으므로 다음 분해가 존재한다.
$$
A \to B \to R^\wedge
$$
여기서 $B$는 $R$ 위에서 매끄럽다. 정리 \ref{smoothing-theorem-popescu}를 보라.
잉여류체를 $\kappa = R/\mathfrak m$로 표기하자. 이는 $R^\wedge$의
잉여류체이기도 하므로 다음 가환 도식을 얻는다.
$$
\xymatrix{
B \ar[rd] \ar@{..>}[r] & R' \ar@{..>}[d] \\
R \ar[r] \ar[u] & \kappa
}
$$
수직 화살표가 매끄러우므로 More on Algebra, Lemma
\ref{more-algebra-lemma-lift-section-smooth-morphism}에 의해 동형
$R/\mathfrak m \to R'/\mathfrak mR'$을 유도하는 에탈 환 사상
$R \to R'$와 위 도식을 가환이 되게 하는 $R$-대수 사상 $B \to R'$가 존재한다.
$R$가 헨젤이므로 $R \to R'$는 절단을 갖는다. Algebra, Lemma
\ref{algebra-lemma-characterize-henselian}을 보라.
$b_i \in R$를 환 사상 $A \to B \to R' \to R$에 의한 $a_i$의 상이라 하자.
이 사상들이 모두 $R$-대수 사상이므로 $(b_1, \ldots, b_n)$은 $R$에서의 해이다.
\end{proof}

\noindent
뇌터 국소환 $(R, \mathfrak m)$, 에탈 환 사상 $R \to R'$ 및
$\kappa(\mathfrak m) = \kappa(\mathfrak m')$를 만족하며 $\mathfrak m$ 위에 놓인
극대 아이디얼 $\mathfrak m' \subset R'$가 주어지면 다음 포함관계가 성립한다.
$$
R \subset R_{\mathfrak m'} \subset R^h \subset R^\wedge,
$$
이는 Algebra, Lemma \ref{algebra-lemma-henselian-functorial-prepare}와
More on Algebra, Lemma \ref{more-algebra-lemma-henselization-noetherian}에서 따른다.

\begin{theorem}
\label{smoothing-theorem-approximation-property-variant}
$R$를 뇌터 국소환이라 하고
$f_1, \ldots, f_m \in R[x_1, \ldots, x_n]$이라 하자.
$(a_1, \ldots, a_n) \in (R^\wedge)^n$가 해라고 하자.
$R$가 G-환이면 모든 정수 $N$에 대해 다음이 존재한다.
\begin{enumerate}
\item 에탈 환 사상 $R \to R'$,
\item $\mathfrak m$ 위에 놓인 극대 아이디얼 $\mathfrak m' \subset R'$,
\item $R'$에서의 해 $(b_1, \ldots, b_n) \in (R')^n$.
\end{enumerate}
이들은 $\kappa(\mathfrak m) = \kappa(\mathfrak m')$ 및
$a_i - b_i \in (\mathfrak m')^NR^\wedge$를 만족한다.
\end{theorem}

\begin{proof}
정리 \ref{smoothing-theorem-approximation-property}에서 이 정리를 도출할 수도 있다.
실제로 More on Algebra, Lemma \ref{more-algebra-lemma-henselization-G-ring}에 의해
헨젤화 $R^h$가 G-환이고, $R^h$를 에탈 확대 $R'$들의 유향 여극한으로 쓸 수 있다.
여기서는 이전 정리의 증명을 이 경우에 다시 수행하여 증명한다.

\medskip\noindent
$a_i - c_i \in \mathfrak m^N$을 만족하는 원소 $c_i \in R$를 잡자.
생성원 $\mathfrak m^N = (d_1, \ldots, d_M)$을 잡고
$a_i = c_i + \sum a_{i, l} d_l$로 쓰자.
다항식환 $R[x_{i, l}]$과 다음 원소들을 생각하자.
$$
g_j = f_j(c_1 + \sum x_{1, l} d_l , \ldots, c_n + \sum x_{n, l} d_{n, l})
\in R[x_{i, l}]
$$
연립방정식 $g_j = 0$은 해 $(a_{i, l})$을 갖는다.
$\kappa(\mathfrak m) = \kappa(\mathfrak m')$를 만족하는 극대 아이디얼
$\mathfrak m'$이 딸린 어떤 에탈 환 사상 $R \to R'$에 대해 $g_j$가
$R'$에서 해 $(b_{i, l})$을 가짐을 보일 수 있다고 하자.
그러면 $b_i = c_i + \sum b_{i, l}d_l$은 $(\mathfrak m')^N$을 법으로
$a_i$와 합동인 $f_j = 0$의 해이다. 따라서 $R^\wedge$ 위의 해의 존재가
$\mathfrak m$ 위의 어떤 소아이디얼에서 자명한 잉여류체 확대를 유도하는
어떤 에탈 환 확대 위의 해의 존재를 함의함을 보이면 충분하다.

\medskip\noindent
$A \subset R^\wedge$를 $a_1, \ldots, a_n$으로 생성되는 $R$-부분대수라 하자.
$R$가 G-환, 즉 $R \to R^\wedge$가 정칙이라고 가정했으므로 다음 분해가 존재한다.
$$
A \to B \to R^\wedge
$$
여기서 $B$는 $R$ 위에서 매끄럽다. 정리 \ref{smoothing-theorem-popescu}를 보라.
잉여류체를 $\kappa = R/\mathfrak m$로 표기하자. 이는 $R^\wedge$의
잉여류체이기도 하므로 다음 가환 도식을 얻는다.
$$
\xymatrix{
B \ar[rd] \ar@{..>}[r] & R' \ar@{..>}[d] \\
R \ar[r] \ar[u] & \kappa
}
$$
수직 화살표가 매끄러우므로 More on Algebra, Lemma
\ref{more-algebra-lemma-lift-section-smooth-morphism}에 의해 동형
$R/\mathfrak m \to R'/\mathfrak mR'$을 유도하는 에탈 환 사상
$R \to R'$와 위 도식을 가환이 되게 하는 $R$-대수 사상 $B \to R'$가 존재한다.
$b_i \in R'$를 환 사상 $A \to B \to R'$에 의한 $a_i$의 상이라 하자.
이 사상들이 모두 $R$-대수 사상이므로 $(b_1, \ldots, b_n)$은 $R'$에서의 해이다.
\end{proof}

\begin{example}
\label{smoothing-example-describe-henselian}
$(R, \mathfrak m)$을 헨젤화가 $R^h$인 뇌터 국소환이라 하자.
완비화 사이의 사상 $R^\wedge \to (R^h)^\wedge$은 동형이다. More on Algebra,
Lemma \ref{more-algebra-lemma-henselization-noetherian}을 보라.
$R^h$도 뇌터이므로(같은 곳 참조), 완비화가 충실히 평탄하다는 사실에 의해
$R^h$를 그 완비화의 부분환으로 생각할 수 있다. 따라서 $R^h$를
$R^\wedge$의 부분환과 동일시할 수 있다.

\medskip\noindent
$R^\wedge$의 어떤 원소가 $R^h$에 속하는지 알아보자. 단순성을 위해
$R$가 분수체 $K$를 갖는 정역이라고 가정한다. 분명히 $R^h$의 모든 원소
$f$는 $R$ 위에서 대수적이다. 즉, $n > 0$이고 $a_n \not = 0$인 어떤
$a_i \in R$에 대해 $a_n f^n + \ldots + a_1 f + a_0 = 0$ 꼴의 방정식이 존재한다.

\medskip\noindent
거꾸로 $f \in R^\wedge$, $n \in \mathbf{N}$이고 $a_n \not = 0$인
$a_0, \ldots, a_n \in R$가
$a_n f^n + \ldots + a_1 f + a_0 = 0$을 만족한다고 하자. $R$가 G-환이면
모든 $N > 0$에 대해 $a_n g^n + \ldots + a_1 g + a_0 = 0$ 및
$f - g \in \mathfrak m^N R^\wedge$를 만족하는 원소 $g \in R^h$가 존재한다.
정리 \ref{smoothing-theorem-approximation-property-variant}를 보라.
$N \gg 0$일 때 $f = g$라고 결론내리고자 한다.
그렇지 않다면 $R^h$에서 $P(T)$의 근 $g$를 무한히 많이 얻는다.
그러나 (1) $R^h \subset R^h \otimes_R K$이고 (2) $R^h \otimes_R K$가
$K$의 유한 개 체 확대의 곱이므로 이는 불가능하다. 실제로 $R \to K$는 단사이고
$R \to R^h$는 평탄하므로 $R^h \to R^h \otimes_R K$는 단사이며,
(2)는 More on Algebra, Lemma \ref{more-algebra-lemma-fibres-henselization}에서 따른다.

\medskip\noindent
결론: $R$가 분수체 $K$를 갖는 뇌터 국소 정역이자 G-환이면,
$R^h \subset R^\wedge$는 $K$ 위에서 대수적인 모든 원소의 집합이다.
\end{example}

\noindent
다음은 이 절의 주정리의 또 다른 변형이다.

\begin{lemma}
\label{smoothing-lemma-approximation-property-variant}
$R$를 뇌터 환이라 하고 $\mathfrak p \subset R$를 소아이디얼이라 하자.
$f_1, \ldots, f_m \in R[x_1, \ldots, x_n]$이라 하자.
$(a_1, \ldots, a_n) \in ((R_\mathfrak p)^\wedge)^n$가 해라고 하자.
$R_\mathfrak p$가 G-환이면 모든 정수 $N$에 대해 다음이 존재한다.
\begin{enumerate}
\item 에탈 환 사상 $R \to R'$,
\item $\mathfrak p$ 위에 놓인 소아이디얼 $\mathfrak p' \subset R'$,
\item $R'$에서의 해 $(b_1, \ldots, b_n) \in (R')^n$.
\end{enumerate}
이들은 $\kappa(\mathfrak p) = \kappa(\mathfrak p')$ 및
$a_i - b_i \in (\mathfrak p')^N(R'_{\mathfrak p'})^\wedge$를 만족한다.
\end{lemma}

\begin{proof}
정리 \ref{smoothing-theorem-approximation-property-variant}에 의해 다음 조건을 만족하는
해 $(b'_1, \ldots, b'_n)$을 어떤 환 $R''$에서 찾을 수 있다. 이 환은
$R_\mathfrak p$ 위에서 에탈이고, $\mathfrak p$ 위에 놓이며
$\kappa(\mathfrak p) = \kappa(\mathfrak p'')$ 및
$a_i - b'_i \in (\mathfrak p'')^N(R''_{\mathfrak p''})^\wedge$를 만족하는
소아이디얼 $\mathfrak p''$이 딸려 있다. 어떤 에탈 $R$-대수 $R'$에 대해
$R'' = R' \otimes_R R_\mathfrak p$로 쓸 수 있다(Algebra, Lemma
\ref{algebra-lemma-etale}를 보라). 필요하면 $R'$을 주국소화로 바꾸어
$(b'_1, \ldots, b'_n)$이 $R'$의 해 $(b_1, \ldots, b_n)$에서 온다고
가정해도 된다. $\mathfrak p' = R' \cap \mathfrak p''$로 두면
$R''_{\mathfrak p''} = R'_{\mathfrak p'}$이므로 증명이 끝난다.
\end{proof}



\section{헨젤 쌍의 근사}
\label{smoothing-section-approximation-pairs}

\noindent
절 \ref{smoothing-section-approximation-G-rings}의 논의를 헨젤 쌍의 경우로 일반화할 수 있다.
헨젤 쌍은 More on Algebra, Section
\ref{more-algebra-section-henselian-pairs}에서 정의하였다.

\begin{lemma}
\label{smoothing-lemma-henselian-pair}
\begin{slogan}
헨젤 쌍의 근사.
\end{slogan}
$(A, I)$를 $A$가 뇌터인 헨젤 쌍이라 하자.
$A^\wedge$를 $A$의 $I$-진 완비화라 하자.
다음 조건 가운데 적어도 하나가 성립한다고 가정하자.
\begin{enumerate}
\item $A \to A^\wedge$는 정칙 환 사상이다.
\item $A$는 뇌터 G-환이다.
\item $(A, I)$는 $B$가 뇌터 G-환인 쌍 $(B, J)$의 헨젤화이다
(More on Algebra, Lemma \ref{more-algebra-lemma-henselization}).
\end{enumerate}
$f_1, \ldots, f_m \in A[x_1, \ldots, x_n]$ 및
$\hat{a}_1, \ldots, \hat{a}_n \in A^\wedge$가 주어져 있고
$j = 1, \ldots, m$에 대해 $f_j(\hat{a}_1, \ldots, \hat{a}_n) = 0$이라 하자.
그러면 모든 $N \geq 1$에 대해 $\hat{a}_i - a_i \in I^N$이고
$j = 1, \ldots, m$에 대해 $f_j(a_1, \ldots, a_n) = 0$인
$a_1, \ldots, a_n \in A$가 존재한다.
\end{lemma}

\begin{proof}
More on Algebra, Lemma \ref{more-algebra-lemma-henselization-pair-G-ring}에 의해
(3)은 (2)를 함의한다. More on Algebra, Lemma
\ref{more-algebra-lemma-map-G-ring-to-completion-regular}에 의해 (2)는 (1)을 함의한다.
따라서 $A \to A^\wedge$가 정칙 환 사상인 경우에 보조정리를 증명하면 충분하다.

\medskip\noindent
$\hat{a}_1, \ldots, \hat{a}_n$을 보조정리의 명제에서와 같이 잡자.
정리 \ref{smoothing-theorem-popescu}에 의해 $A \to P$가 매끄럽고
$B$에서 $f_j(b_1, \ldots, b_n) = 0$인 $b_1, \ldots, b_n \in B$가 딸린
분해 $A \to B \to A^\wedge$를 찾을 수 있다.
합성을 $\sigma : B \to A^\wedge \to A/I^N$로 표기하자.
More on Algebra, Lemma \ref{more-algebra-lemma-lift-section-smooth-morphism}에 의해
동형 $A/I^N \to A'/I^NA'$을 유도하는 에탈 환 사상 $A \to A'$와
$\sigma$를 올리는 $A$-대수 사상 $\tilde \sigma : B \to A'$를 찾을 수 있다.
$(A, I)$가 헨젤이므로 $A$-대수 사상 $\chi : A' \to A$가 존재한다.
More on Algebra, Lemma \ref{more-algebra-lemma-characterize-henselian-pair}를 보라.
이제 $a_i = \chi(\tilde \sigma(b_i))$로 두면 해를 얻는다.
\end{proof}
